% Section file for Chapter: Twisted Holography and Quantum Gravity
% Result (G1): Perturbative Finiteness
%
% Develops the HS-sewing framework, shadow free energies,
% convergence of the gravitational partition function, the
% Heisenberg prototype, finiteness for the standard landscape,
% and implications for 3d quantum gravity.

\section{Perturbative finiteness}
\label{sec:thqg-perturbative-finiteness}
\label{ch:thqg-perturbative-finiteness}% alias for dangling \ref{ch:...}
\index{perturbative finiteness|textbf}
\index{UV finiteness!twisted holography}

% =====================================
% Local macros: provided only if absent from the ambient manuscript.
% =====================================
\providecommand{\HS}{\operatorname{HS}}
\providecommand{\Tr}{\operatorname{Tr}}
\providecommand{\Det}{\operatorname{Det}}
\providecommand{\Fred}{\operatorname{Fred}}
\providecommand{\rank}{\operatorname{rank}}
\providecommand{\Sym}{\operatorname{Sym}}
\providecommand{\id}{\mathrm{id}}
\providecommand{\ad}{\operatorname{ad}}
\providecommand{\wt}{\operatorname{wt}}
\providecommand{\Aut}{\operatorname{Aut}}
\providecommand{\Res}{\operatorname{Res}}
\providecommand{\Defcyc}{\operatorname{Def}_{\mathrm{cyc}}}
\providecommand{\gAmod}{\mathfrak{g}_{\cA}^{\mathrm{mod}}}
\providecommand{\Gmod}{\mathcal{G}^{\mathrm{mod}}}

Perturbative finiteness of twisted quantum gravity is the analytic form of the gravitational shadow theorem. Three implications form a chain: $D_\cA^2 = 0$ (Theorem~\ref*{V1-thm:mc2-bar-intrinsic}) furnishes $\Theta_\cA \in \mathrm{MC}(\gAmod)$ at all genera without truncation; the shadow representation $\operatorname{Sh}_{g,0}(\Theta_\cA)$ produces a cohomology class on the compact orbifold $\overline{\mathcal{M}}_g$ whose fixed-degree integral is finite; on the uniform-weight scalar lane, Bernoulli asymptotics of the Faber--Pandharipande coefficients give convergence with radius $4\pi^2$. Standard perturbative quantum gravity, by contrast, diverges at genus $g \geq 2$: the integration domain is noncompact, the integrand distributional, the expansion asymptotic. In the twisted shadow setting, the Maurer--Cartan equation is algebraic and Fulton--MacPherson compactification provides a geometric regulator. The scalar genus series converges on the proved scalar lane; the full non-scalar genus series carries the additional shadow-depth growth problem.

% ======================================================================
% SUBSECTION 1: THE HS-SEWING FRAMEWORK FOR TWISTED GRAVITY
% ======================================================================

\subsection{The HS-sewing framework for twisted gravity}
\label{subsec:thqg-I-hs-sewing-framework}
\index{HS-sewing!twisted gravity framework}

The passage from finite-dimensional algebraic data (the OPE coefficients of~$\cA$) to infinite-genus partition functions requires a bridge between algebra and analysis. The HS-sewing condition provides this bridge: it guarantees that sewing amplitudes on pairs of pants compose to trace-class operators on the sewing Hilbert space.

\subsubsection{The sewing Hilbert space}

\begin{definition}[Sewing Hilbert space; cf.\ {\cite{costello-renormalization}}]
% label removed: def:thqg-I-sewing-hilbert-space
\index{sewing Hilbert space|textbf}
Let $\cA$ be a positive-energy chiral algebra with conformal grading $\cA = \bigoplus_{n \geq 0} H_n$, where $\dim H_n < \infty$ for all~$n$. For a sewing parameter $q \in (0,1)$, the \emph{sewing Hilbert space} is the $q$-weighted $\ell^2$-completion
\begin{equation}% label removed: eq:thqg-I-sewing-hilbert
\mathcal{H}_q(\cA) \;:=\; \widehat{\bigoplus}_{n \geq 0} q^{n/2}\,H_n \;=\; \Bigl\{ \sum_{n=0}^\infty v_n \;\Big|\; v_n \in H_n,\; \sum_{n=0}^\infty q^n \|v_n\|^2 < \infty \Bigr\},
\end{equation}
where $\|\cdot\|$ is any fixed Hermitian norm on each finite-dimensional $H_n$. The inner product is $\langle v, w \rangle_q = \sum_{n \geq 0} q^n \langle v_n, w_n \rangle_{H_n}$.
\end{definition}

\begin{remark}[Independence of norms]
% label removed: rem:thqg-I-norm-independence
Since each $H_n$ is finite-dimensional, different choices of Hermitian norm on $H_n$ give equivalent topologies on $\mathcal{H}_q(\cA)$. The $q$-weighting ensures that the completion is sensitive to conformal weight, with high-weight states exponentially suppressed. For $q \to 1^-$, the Hilbert space degenerates; for $q \to 0^+$, only the vacuum sector $H_0$ survives.
\end{remark}

\begin{lemma}[Completeness and separability; \ClaimStatusProvedHere]
% label removed: lem:thqg-I-completeness
\index{sewing Hilbert space!completeness}
$\mathcal{H}_q(\cA)$ is a separable Hilbert space for every $q \in (0,1)$. A Cauchy sequence $\{v^{(k)}\}_{k \geq 1}$ converges if and only if $\sum_{n \geq 0} q^n \|v^{(k)}_n - v^{(\ell)}_n\|^2 \to 0$ as $k, \ell \to \infty$.
\end{lemma}

\begin{proof}
Each $H_n$ is finite-dimensional, hence complete. The $q$-weighted $\ell^2$-completion is a weighted direct sum of Hilbert spaces with weights $q^n$; completeness of the direct sum follows from that of each summand together with the summability condition. Separability follows from taking a countable dense subset of each $H_n$ (finite-dimensional spaces are separable) and forming finite linear combinations.
\end{proof}

\begin{proposition}[Inclusion maps; \ClaimStatusProvedHere]
% label removed: prop:thqg-I-inclusion-maps
\index{sewing Hilbert space!inclusion}
For $0 < q_1 < q_2 < 1$, the identity map on $\cA$ extends to a bounded inclusion $\iota_{q_1,q_2} \colon \mathcal{H}_{q_2}(\cA) \hookrightarrow \mathcal{H}_{q_1}(\cA)$ with operator norm $\|\iota_{q_1,q_2}\| = 1$. In particular, for every $v \in \mathcal{H}_{q_2}(\cA)$:
\begin{equation}% label removed: eq:thqg-I-inclusion
\|v\|_{q_1}^2 \;=\; \sum_{n \geq 0} q_1^n \|v_n\|^2 \;\leq\; \sum_{n \geq 0} q_2^n \|v_n\|^2 \;=\; \|v\|_{q_2}^2.
\end{equation}
\end{proposition}

\begin{proof}
Since $q_1 < q_2$, we have $q_1^n \leq q_2^n$ for all $n \geq 0$, giving the inequality termwise.
\end{proof}

\subsubsection{The HS-sewing condition}

\begin{definition}[HS-sewing condition]
\label{def:thqg-I-hs-sewing}
\index{HS-sewing!definition|textbf}
Let $\cA$ be a positive-energy chiral algebra with OPE structure constants $C^{c,k}_{a,i;\,b,j}$, where $a,b,c$ are conformal weights and $i,j,k$ index bases of $H_a$, $H_b$, $H_c$ respectively. The \emph{pair-of-pants multiplication operator} at weight $(a,b) \to c$ is the linear map $m^c_{a,b} \colon H_a \otimes H_b \to H_c$ with matrix entries $(m^c_{a,b})_{ij}^k = C^{c,k}_{a,i;\,b,j}$. The algebra~$\cA$ satisfies the \emph{HS-sewing condition} at parameter $q \in (0,1)$ if
\begin{equation}% label removed: eq:thqg-I-hs-sewing-condition
\boxed{ \sum_{a,b,c \geq 0} q^{a+b+c}\, \|m^c_{a,b}\|_{\HS}^2 \;<\; \infty,}
\end{equation}
where $\|m^c_{a,b}\|_{\HS}^2 = \sum_{i,j,k} |C^{c,k}_{a,i;\,b,j}|^2$ is the Hilbert--Schmidt norm of the multiplication operator.
\end{definition}

\begin{remark}[Physical content]
% label removed: rem:thqg-I-hs-physical
The HS-sewing condition controls the convergence of the sewing construction: a pair-of-pants decomposition of a genus-$g$ surface with $n$ punctures involves $2g - 2 + n$ pairs of pants, and each sewing circle contributes one composition of pair-of-pants operators. The HS condition guarantees that each such composition is Hilbert--Schmidt, and that the composition of two Hilbert--Schmidt operators is trace class, the key requirement for the trace that computes the partition function. The $q$-parameter is the sewing modulus $q = e^{-t}$ where $t > 0$ is the collar length.
\end{remark}

\begin{definition}[Hilbert--Schmidt and trace-class norms]
% label removed: def:thqg-I-hs-trace-norms
\index{Hilbert--Schmidt norm}
\index{trace-class!operator}
For a linear map $T \colon V \to W$ between finite-dimensional Hilbert spaces, the Hilbert--Schmidt norm is $\|T\|_{\HS}^2 = \Tr(T^*T) = \sum_{i,j} |T_{ij}|^2$. An operator $K$ on a Hilbert space $H$ is \emph{trace class} if $\|K\|_1 := \Tr\sqrt{K^*K} < \infty$. The key composition inequality is:
\begin{equation}% label removed: eq:thqg-I-hs-composition
\|AB\|_1 \;\leq\; \|A\|_{\HS}\,\|B\|_{\HS}
\end{equation}
for Hilbert--Schmidt operators $A, B$. Hence the composition of two HS operators is trace class, with trace bounded by the product of HS norms.
\end{definition}

\begin{lemma}[Schatten class hierarchy; \ClaimStatusProvedHere]
\label{lem:thqg-I-schatten}
\index{Schatten class}
Let $H$ be a separable Hilbert space and $S_p(H)$ the $p$-th Schatten class ($p \geq 1$). Then:
\begin{enumerate}[label=\textup{(\roman*)}]
\item $S_1(H) \subset S_2(H) \subset S_p(H) \subset \mathcal{B}(H)$ for $1 \leq p \leq \infty$, where $\mathcal{B}(H)$ is bounded operators.
\item $S_1(H)$ is a two-sided ideal in $\mathcal{B}(H)$: if $K \in S_1(H)$ and $T \in \mathcal{B}(H)$, then $KT, TK \in S_1(H)$.
\item If $A \in S_p(H)$ and $B \in S_q(H)$ with $1/p + 1/q = 1/r \leq 1$, then $AB \in S_r(H)$.
\end{enumerate}
The sewing application uses $p = q = 2$, $r = 1$: two Hilbert--Schmidt operators compose to trace class.
\end{lemma}

\begin{proof}
Standard functional analysis; see Reed--Simon~\cite{RS78} or Simon~\cite{Simon05}. Part~(iii) is the generalized H\"older inequality for Schatten classes.
\end{proof}

\begin{proposition}[Pair-of-pants as HS operator; \ClaimStatusProvedElsewhere]
\label{prop:thqg-I-pants-hs}
\index{pair of pants!HS operator}
Under HS-sewing~\eqref{V1-eq:thqg-I-hs-sewing-condition}, the $q$-weighted multiplication operator $m_q \colon \mathcal{H}_q(\cA) \,\widehat{\otimes}\, \mathcal{H}_q(\cA) \to \mathcal{H}_q(\cA)$ is Hilbert--Schmidt, with
\begin{equation}% label removed: eq:thqg-I-mq-hs-bound
\|m_q\|_{\HS}^2 \;=\; \sum_{a,b,c \geq 0} q^{a+b+c}\,\|m^c_{a,b}\|_{\HS}^2 \;<\; \infty.
\end{equation}
Consequently, every closed amplitude on a genus-$g$ surface obtained by sewing is trace class.
\end{proposition}

\begin{proof}
Direct from the definition: $\|m_q\|_{\HS}^2$ is precisely the left-hand side of~\eqref{V1-eq:thqg-I-hs-sewing-condition}. A genus-$g$ surface with $n$ punctures is obtained by sewing $2g - 2 + n$ pairs of pants along $3g - 3 + n$ circles. Each sewing circle contributes one pair-of-pants composition. By the Hilbert--Schmidt composition inequality~\eqref{V1-eq:thqg-I-hs-composition}, each composition of two HS operators gives a trace-class operator. Since $g \geq 1$ requires at least two compositions, and trace-class operators form an ideal, the result follows.
\end{proof}

\begin{remark}[Sewing topology]
% label removed: rem:thqg-I-sewing-topology
\index{sewing!topology}
The pair-of-pants decomposition of a genus-$g$ surface $\Sigma_{g,n}$ into $2g - 2 + n$ trinions and $3g - 3 + n$ sewing circles is determined by a pants graph $\Gamma_P$. Different choices of $\Gamma_P$ give different factorizations of the sewing operator, but the trace (the partition function) is independent of $\Gamma_P$ by the factorization property. The sewing topology on the space of pants decompositions is captured by the flip graph, whose vertices are pants decompositions and whose edges are elementary moves (Dehn twists and handle slides). The partition function is invariant under flips by the pentagon and hexagon identities of the chiral algebra.
\end{remark}

\subsubsection{The polynomial-OPE / subexponential-sector criterion}

The analytic input is the criterion of Theorem~\ref{thm:general-hs-sewing}: polynomial OPE growth and subexponential sector growth together imply HS-sewing for all $q \in (0,1)$.

\begin{theorem}[HS-sewing from polynomial OPE growth; \ClaimStatusProvedHere]
\label{thm:general-hs-sewing}
\index{HS-sewing!polynomial OPE criterion}
Let\/ $\cA$ be a positive-energy chiral algebra satisfying:
\begin{enumerate}[label=\textup{(\alph*)}]
\item \emph{Subexponential sector growth:} $\log\dim H_n = o(n)$ as $n \to \infty$.
\item \emph{Polynomial OPE growth:} There exist constants $K > 0$ and $N \in \mathbb{Z}_{\geq 0}$ such that $|C^{c,k}_{a,i;\,b,j}| \leq K(a + b + c + 1)^N$ for all conformal weights $a,b,c$ and all basis indices $i,j,k$.
\end{enumerate}
Then $\cA$ satisfies the HS-sewing condition for every $q \in (0,1)$:
\begin{equation}% label removed: eq:thqg-I-hs-criterion-bound
\sum_{a,b,c \geq 0} q^{a+b+c}\,\|m^c_{a,b}\|_{\HS}^2 \;\leq\; K^2 \sum_{n \geq 0} P_3(n)\,n^{2N}\, \bigl(\!\dim H_n\bigr)^3\,q^n \;<\; \infty,
\end{equation}
where $P_3(n) = \binom{n+2}{2}$ counts the number of triples $(a,b,c)$ with $a + b + c = n$.
\end{theorem}

\begin{proof}
The HS norm at weight level $(a,b,c)$ satisfies
\[
\|m^c_{a,b}\|_{\HS}^2 \;=\; \sum_{i=1}^{\dim H_a} \sum_{j=1}^{\dim H_b} \sum_{k=1}^{\dim H_c} |C^{c,k}_{a,i;\,b,j}|^2 \;\leq\; \dim H_a \cdot \dim H_b \cdot \dim H_c \cdot K^2(a{+}b{+}c{+}1)^{2N}.
\]
Summing over triples with $a + b + c = n$:
\[
\sum_{\substack{a,b,c \geq 0 \\ a+b+c = n}} q^n\,\|m^c_{a,b}\|_{\HS}^2 \;\leq\; K^2 q^n n^{2N} \sum_{\substack{a,b,c \geq 0 \\ a+b+c = n}} \dim H_a \cdot \dim H_b \cdot \dim H_c.
\]
The inner sum over triples is bounded by $\bigl(\sum_{j=0}^n \dim H_j\bigr)^3$. By the subexponential growth hypothesis, for every $\varepsilon > 0$ there exists $C_\varepsilon$ with $\dim H_j \leq C_\varepsilon \,e^{\varepsilon j}$. Hence $\sum_{j=0}^n \dim H_j \leq C_\varepsilon\,n\,e^{\varepsilon n}$. The total sum is therefore bounded by $K^2 C_\varepsilon^3 \sum_{n \geq 0} n^{2N+3}\, e^{3\varepsilon n}\,q^n$. Choosing $\varepsilon < |\log q|/4$, the exponential factor $e^{3\varepsilon n} q^n = e^{(3\varepsilon + \log q)n}$ decays geometrically, ensuring convergence of the power series.
\end{proof}

\begin{remark}[Optimality of the bound]
% label removed: rem:thqg-I-hs-optimality
The polynomial-OPE criterion is not optimal: examples with super-polynomial but sub-exponential OPE growth (such as certain completed W-algebras) may still satisfy HS-sewing. Nevertheless, the criterion covers the entire standard landscape of finitely strongly generated universal chiral algebras. The bound~\eqref{V1-eq:thqg-I-hs-criterion-bound} is loose by factors of order $n^3$ due to the crude estimate $\sum_{a+b+c=n} \dim H_a \dim H_b \dim H_c \leq (\sum_{j \leq n} \dim H_j)^3$; tighter bounds are available for specific families.
\end{remark}

\begin{remark}[The gap between HS-sewing and convergence]
% label removed: rem:thqg-I-gap-hs
\index{HS-sewing!necessary vs sufficient}
The HS-sewing condition is \emph{sufficient} for absolute convergence of the genus-$g$ partition function but not \emph{necessary}. Weaker conditions (e.g., trace-class sewing without the HS factorization) may suffice at specific genera. The HS condition provides uniform control across all genera: a single verification at the level of OPE data yields finiteness at every genus.
\end{remark}

\subsubsection{Explicit HS norms for the standard landscape}

\begin{computation}[HS norms for Heisenberg]
\label{comp:thqg-I-hs-heisenberg}
\index{Heisenberg!HS norms}
For the Heisenberg algebra $\mathcal{H}_k$ of rank~$1$ at level $k \in \mathbb{C}^{\times}$:
\begin{itemize}
\item $\dim H_n = p(n)$ (partitions of~$n$);
\item OPE: $a(z)a(w) \sim k/(z-w)^2$, giving $C^c_{a,b} = k\,\delta_{a+b,c+1}$ for the fundamental OPE, and all composite OPE grow polynomially with degree $N = 1$ (the Wick contraction is bilinear).
\item Growth: $\log p(n) \sim \pi\sqrt{2n/3}$ (Hardy--Ramanujan), so $\log \dim H_n = O(\sqrt{n})$, which is subexponential.
\end{itemize}
In the sewing normalization used by \ref{item:thqg-I-H2}, Wick's
theorem gives a constant $C_k$ such that the weighted multiplication
satisfies
\begin{equation}% label removed: eq:thqg-I-hs-heisenberg-norm
\|m_q\|_{\HS}^2
\;\leq\;
C_k^2\sum_{n \geq 0} P_3(n)\,n^2
\!\sum_{\substack{a+b+c=n}} p(a)p(b)p(c)\,q^n
\;<\;\infty ,
\end{equation}
for every $q \in (0,1)$. The scalar $k$ enters $C_k$; it is not an
oscillator colour.

The one-particle operator has eigenvalues $q^n$, giving
$\|T_q\|_1 = \sum_{n=1}^\infty q^n = q/(1-q)$. The two-leg
state-counting identity is
\begin{equation}% label removed: eq:thqg-I-hs-heisenberg-factorization
\sum_{a,b \geq 0} p(a)p(b)q^{a+b}
\;=\;
\prod_{m=1}^\infty \frac{1}{(1-q^m)^2}.
\end{equation}
It is a dimension series for two Fock legs, not the square series
$\sum_n p(n)^2q^n$ and not a source of powers of $k$.
\end{computation}

\begin{computation}[HS norms for affine Kac--Moody]
% label removed: comp:thqg-I-hs-affine
\index{affine Kac--Moody!HS norms}
For $\widehat{\fg}_k$ at non-critical level $k \neq -h^\vee$, with $\fg$ simple of dimension~$d$ and dual Coxeter number $h^\vee$:
\begin{itemize}
\item $\dim H_n \sim C\,n^{(d-\rank)/2}\,e^{\pi\sqrt{2dn/3}}$ (Kac--Peterson), so $\log \dim H_n = O(\sqrt{n})$, subexponential.
\item OPE: $J^a(z)J^b(w) \sim k\delta^{ab}/(z-w)^2 + f^{ab}{}_c J^c(w)/(z-w)$, polynomial of degree $N = 1$.
\item The HS norm satisfies
 \begin{equation}% label removed: eq:thqg-I-hs-affine-norm
 \|m_q\|_{\HS}^2 \;\leq\; C_k^2 \prod_{n=1}^\infty \frac{1}{(1-q^n)^{2d}}
 \end{equation}
 where $C_k$ depends on $k$ and the structure constants of $\fg$.
\end{itemize}
The key estimate for $C_k$ is obtained from the OPE normalization. The structure constants of $\widehat{\fg}_k$ at weight level $n$ grow as $|C^c_{a,b}| \leq \|f\|_\infty \cdot (k+h^\vee)$, where $\|f\|_\infty = \max_{a,b,c} |f^{ab}{}_c|$ is the maximum structure constant. For $\fg = \mathfrak{sl}_N$, one has $\|f\|_\infty = 1$ in the Chevalley basis, giving $C_k = |k + h^\vee|$.
\end{computation}

\begin{computation}[HS norms for Virasoro]
% label removed: comp:thqg-I-hs-virasoro
\index{Virasoro!HS norms}
For the Virasoro algebra $\mathrm{Vir}_c$ at central charge $c \in \mathbb{C}$:
\begin{itemize}
\item $\dim H_n = p(n)$ (Verma module dimensions);
\item OPE: $T(z)T(w) \sim c/(2(z-w)^4) + 2T(w)/(z-w)^2 + \partial T(w)/(z-w)$, polynomial of degree $N = 2$ (the $(z-w)^{-4}$ pole gives OPE coefficients growing at most quadratically in conformal weight).
\item Growth: $\log p(n) \sim \pi\sqrt{2n/3}$, subexponential.
\item The HS sum converges:
 \begin{equation}% label removed: eq:thqg-I-hs-virasoro-norm
 \|m_q\|_{\HS}^2 \;\leq\; C_c^2 \sum_{n \geq 0} p(n)^3\,n^{4}\,q^n \;<\; \infty \quad\text{for all } q \in (0,1).
 \end{equation}
\end{itemize}
The quartic growth of $n^4$ in the bound reflects the quartic OPE pole of $T(z)T(w)$. For the Virasoro algebra, the structure constants $C^c_{a,b}$ at weight level $n = a + b - c$ satisfy $|C^c_{a,b}| \leq P_2(n) := (n+1)(n+2)/2$ when $c$ is bounded; the precise bound involves the central charge through $|c|/2$ from the fourth-order pole.
\end{computation}

\begin{computation}[HS norms for $\beta\gamma$ system]
% label removed: comp:thqg-I-hs-betagamma
\index{beta-gamma system@$\beta\gamma$ system!HS norms}
For the $\beta\gamma$ system of conformal weights $(1,0)$:
\begin{itemize}
\item $\dim H_n = p(n)^2$ (two independent partition functions, one for $\beta$ modes and one for $\gamma$ modes).
\item OPE: $\beta(z)\gamma(w) \sim 1/(z-w)$, a simple pole. All composite OPE grow polynomially with degree $N = 1$.
\item Growth: $\log p(n)^2 = 2\log p(n) \sim 2\pi\sqrt{2n/3}$, subexponential.
\item The HS norm:
 \begin{equation}% label removed: eq:thqg-I-hs-betagamma-norm
 \|m_q\|_{\HS}^2 \;\leq\; C^2\prod_{n=1}^\infty \frac{1}{(1-q^n)^{4}} \;<\; \infty \quad\text{for all } q \in (0,1).
 \end{equation}
\end{itemize}
The exponent $-4$ in the product reflects the two bosonic degrees of freedom ($\beta$ and $\gamma$), each contributing $\prod_{n=1}^\infty (1-q^n)^{-2}$.
\end{computation}

\begin{computation}[HS norms for $\mathcal{W}_N$]
% label removed: comp:thqg-I-hs-wn
\index{W-algebra@$\mathcal{W}$-algebra!HS norms}
For the principal $\mathcal{W}$-algebra $\mathcal{W}_k(\mathfrak{sl}_N)$ at generic level:
\begin{itemize}
\item $\dim H_n$ is the number of $N$-colored partitions of~$n$ with colors of weight $2, 3, \ldots, N$; the growth is $\log \dim H_n = O(n^{(N-1)/N})$, subexponential for all~$N$.
\item OPE growth: the maximal singular order of the $W^{(s)}$ self-OPE is $2s$, so $N_{\max} = 2N$ for the $W^{(N)}$ field. All OPE coefficients grow at most polynomially with degree $\leq 2N$.
\item Convergence:
 \begin{equation}% label removed: eq:thqg-I-hs-wn-norm
 \|m_q\|_{\HS}^2 \;\leq\; K_N^2(c) \sum_{n \geq 0} (\dim H_n)^3\,n^{4N}\,q^n \;<\; \infty \quad\text{for all } q \in (0,1).
 \end{equation}
\end{itemize}
The $N$-colored partition function has generating series $\prod_{s=2}^N \prod_{m=s}^\infty (1-q^m)^{-1}$, and its growth exponent is controlled by the Meinardus theorem: $\log(\dim H_n) \sim C_N \cdot n^{(N-1)/N}$ with $C_N = N \cdot (\Gamma(1/N) \zeta(1+1/N))^{N/(N+1)} / (N+1)$. For $N = 2$ (Virasoro), $C_2 = \pi\sqrt{2/3}$; for $N = 3$ ($\mathcal{W}_3$), the growth is $O(n^{2/3})$.
\end{computation}

\begin{computation}[HS norms for lattice VOAs]
% label removed: comp:thqg-I-hs-lattice
\index{lattice VOA!HS norms}
For the lattice vertex algebra $V_\Lambda$ of an even integral lattice $\Lambda$ of rank~$r$:
\begin{itemize}
\item $\dim H_n = |\Lambda^*/\Lambda| \cdot p_r(n)$, where $p_r(n)$ is the $r$-colored partition function. Growth: $\log \dim H_n = O(\sqrt{n})$, subexponential.
\item OPE: inherited from rank-$r$ Heisenberg plus lattice vertex operators $V_\alpha(z)V_\beta(w) \sim (z-w)^{\langle\alpha,\beta\rangle}:\!V_{\alpha+\beta}(w)\!:$; the OPE growth is polynomial of degree $N = \max(r, \max_{\alpha \in \Lambda}|\langle\alpha,\alpha\rangle|/2)$.
\item The HS norm factors as
 \begin{equation}% label removed: eq:thqg-I-hs-lattice-norm
 \|m_q\|_{\HS}^2 \;\leq\; C_\Lambda^2\prod_{n=1}^\infty \frac{1}{(1-q^n)^{2r}} \,\Theta_\Lambda(q)^2 \;<\; \infty \quad\text{for all } q \in (0,1),
 \end{equation}
 where $\Theta_\Lambda(q) = \sum_{\alpha \in \Lambda} q^{\langle\alpha,\alpha\rangle/2}$ is the lattice theta function.
\end{itemize}
The factorization reflects the tensor-product structure $V_\Lambda = \mathcal{H}^{\otimes r} \otimes \mathbb{C}[\Lambda]$: the rank-$r$ Heisenberg part contributes $\prod_{n=1}^\infty (1-q^n)^{-2r}$ and the lattice group ring contributes $\Theta_\Lambda(q)^2$ (one factor from the bra side, one from the ket side of the HS norm).
\end{computation}

\begin{computation}[HS norms for free fermions]
\label{comp:thqg-I-hs-fermion}
\index{free fermions!HS norms}
For the free-fermion VOA of rank~$n$ (i.e., $n$ chiral fermions $\psi_1, \ldots, \psi_n$ of conformal weight $1/2$):
\begin{itemize}
\item $\dim H_m = \binom{n}{m}$ for $m \leq n$, $\dim H_m = 0$ for $m > n$ when counting weight-$m/2$ states; more precisely, the full Fock space has $\dim H_m$ growing as partitions into distinct parts, with generating function $\prod_{j=1}^n (1 + q^{j-1/2})$.
\item OPE: $\psi_i(z)\psi_j(w) \sim \delta_{ij}/(z-w)$, simple pole. OPE degree $N = 0$ (the structure constants are $0$ or $\pm 1$).
\item The HS norm:
 \begin{equation}% label removed: eq:thqg-I-hs-fermion-norm
 \|m_q\|_{\HS}^2 \;\leq\; \prod_{j=1}^{n} (1 + q^{j-1/2})^2 \;<\; \infty \quad\text{for all } q \in (0,1).
 \end{equation}
 This bound is finite because it is a finite product of bounded terms.
\end{itemize}
\end{computation}

\begin{table}[ht]
\centering
\caption{HS-sewing growth rates for the standard landscape}
% label removed: tab:thqg-I-hs-growth
\index{HS-sewing!growth rates}
\renewcommand{\arraystretch}{1.4}
{\footnotesize
\begin{tabular}{|l|c|c|c|c|}
\hline
\textbf{Family $\cA$} & $\boldsymbol{\dim H_n}$ & \textbf{OPE growth $N$} & \textbf{Sector growth} & \textbf{HS bound} \\
\hline
$\mathcal{H}_k$ (rank 1) & $p(n)$ & $1$ & $O(\sqrt{n})$ & $C_k^2\sum_n p(n)^3n^4q^n$ \\
\hline
$\widehat{\fg}_k$ & $\sim e^{\pi\sqrt{2dn/3}}$ & $1$ & $O(\sqrt{n})$ & $C_k^2\prod_{n=1}^\infty (1-q^n)^{-2d}$ \\
\hline
$\mathrm{Vir}_c$ & $p(n)$ & $2$ & $O(\sqrt{n})$ & $C_c^2\sum p(n)^3 n^4 q^n$ \\
\hline
$\beta\gamma$ system & $p(n)^2$ & $1$ & $O(\sqrt{n})$ & $\prod_{n=1}^\infty (1-q^n)^{-4}$ \\
\hline
$\mathcal{W}_N$ & $N$-colored parts & $2N$ & $O(n^{(N-1)/N})$ & $K_N^2(c)\sum (\dim H_n)^3 n^{4N} q^n$ \\
\hline
$V_\Lambda$ (rank $r$) & $|\Lambda^*/\Lambda|\,p_r(n)$ & $r$ & $O(\sqrt{n})$ & $C_\Lambda^2\prod_{n=1}^\infty (1-q^n)^{-2r}\Theta_\Lambda^2$ \\
\hline
Free fermions & $\leq 2^n$ & $0$ & polynomial & $\prod (1+q^{j-1/2})^2$ \\
\hline
\end{tabular}}
\end{table}

\subsubsection{Connection to trace-class operators on sewing Hilbert spaces}

\begin{proposition}[Trace-class amplitudes; \ClaimStatusProvedElsewhere]
% label removed: prop:thqg-I-trace-class-amplitudes
\index{trace class!sewing amplitudes}
Let\/ $\cA$ satisfy the HS-sewing condition at parameter $q$. Let $\Sigma_{g,n}$ be a genus-$g$ surface with $n$ punctures and a pants decomposition $\mathcal{P}$ with collar parameters $q_1, \ldots, q_{3g-3+n} \in (0,q]$. Then the sewing amplitude
\begin{equation}% label removed: eq:thqg-I-sewing-amplitude
Z_{\mathcal{P}}(\cA;\, q_1, \ldots, q_{3g-3+n}) \;=\; \Tr_{\mathcal{H}_q(\cA)^{\otimes n}} \Bigl( \prod_{\text{sewing circles}} m_{q_i} \Bigr)
\end{equation}
is well defined and absolutely convergent. The amplitude depends on the sewing moduli $(q_1, \ldots, q_{3g-3+n})$ but not on the choice of pants decomposition (by the factorization property of the chiral algebra).
\end{proposition}

\begin{proof}
Each pants composition $m_{q_i}$ is Hilbert--Schmidt by Proposition~\ref{prop:thqg-I-pants-hs}. A genus-$g$ surface requires $2g - 2 + n$ pairs of pants and $3g - 3 + n$ sewing circles. The sewing amplitude is a trace of a composition of $3g - 3 + n$ Hilbert--Schmidt operators. Since $3g - 3 + n \geq 2$ for $g \geq 1$ (and $\geq 3$ for $n = 0$), at least one pair of consecutive HS operators composes to give a trace-class operator. The remaining HS operators are bounded, and trace class times bounded is still trace class (Lemma~\ref{lem:thqg-I-schatten}(ii)). Hence the trace converges absolutely. Independence from the pants decomposition follows from the associativity and factorization axioms of the chiral algebra~\cite{BD04, FBZ04}.
\end{proof}

\begin{proposition}[Uniform bound on the trace; \ClaimStatusProvedHere]
\label{prop:thqg-I-uniform-trace-bound}
\index{trace class!uniform bound}
Under HS-sewing at parameter $q$, the sewing amplitude on a genus-$g$ surface with $n$ punctures satisfies the uniform bound
\begin{equation}% label removed: eq:thqg-I-uniform-trace
|Z_{\mathcal{P}}(\cA;\,q_1,\ldots,q_{3g-3+n})| \;\leq\; \|m_q\|_{\HS}^{3g-3+n}
\end{equation}
for all $q_1, \ldots, q_{3g-3+n} \in (0,q]$. In particular, the bound is uniform in the sewing moduli.
\end{proposition}

\begin{proof}
Each factor $m_{q_i}$ satisfies $\|m_{q_i}\|_{\HS} \leq \|m_q\|_{\HS}$ (since $q_i \leq q$ and the HS norm is monotone increasing in $q$). The trace of a composition of operators satisfies $|\Tr(A_1 \cdots A_k)| \leq \|A_1\|_{\HS}\cdots\|A_k\|_{\HS}$ when $k \geq 2$ (by iterated application of the HS composition inequality and the trace bound $|\Tr(K)| \leq \|K\|_1$).
\end{proof}

\begin{corollary}[Genus-$g$ partition function]
\label{cor:thqg-I-genus-g-partition}
\index{partition function!genus $g$}
For $\cA$ satisfying HS-sewing, the genus-$g$ partition function
\begin{equation}% label removed: eq:thqg-I-Z-g
Z_g(\cA) \;:=\; \int_{\overline{\mathcal{M}}_g} Z_{\mathcal{P}}(\cA;\,\mathbf{q}) \;\mu_{\mathrm{WP}}
\end{equation}
is well defined as an integral over the Deligne--Mumford compactification, where $\mu_{\mathrm{WP}}$ denotes the Weil--Petersson volume form.
\end{corollary}

\begin{proof}
The trace-class norm provides a pointwise bound on the integrand: $|Z_{\mathcal{P}}(\cA;\,\mathbf{q})| \leq \|T_{\mathbf{q}}\|_1$ where $T_{\mathbf{q}}$ is the sewing operator. By the HS-bound (Proposition~\ref{prop:thqg-I-uniform-trace-bound}), this norm is bounded uniformly in the interior of the moduli space. Near the boundary, the Fulton--MacPherson compactification provides the regulator: the sewing moduli $q_i$ vanish on the boundary divisors, and the $q$-weighted HS norms degenerate to zero. The integral over the compact orbifold $\overline{\mathcal{M}}_g$ is therefore finite.
\end{proof}

\begin{remark}[The Deligne--Mumford boundary]
% label removed: rem:thqg-I-dm-boundary
\index{Deligne--Mumford!boundary}
The boundary $\partial\overline{\mathcal{M}}_g = \overline{\mathcal{M}}_g \setminus \mathcal{M}_g$ consists of stable curves with nodes. Near a boundary stratum $\delta_\Gamma$ indexed by a dual graph $\Gamma$, the sewing moduli satisfy $q_e \to 0$ for each edge $e$ of $\Gamma$. The integrand $Z_{\mathcal{P}}(\cA;\,\mathbf{q})$ decomposes as a product of amplitudes on the irreducible components of the stable curve, multiplied by propagator insertions at the nodes. The vanishing of $q_e$ at the boundary ensures that the propagator insertions are trace class with trace-class norm tending to zero, giving the regularity needed for integrability over $\overline{\mathcal{M}}_g$.
\end{remark}


% ======================================================================
% SUBSECTION 2: SHADOW FREE ENERGIES AND THE BERNOULLI GENERATING FUNCTION
% ======================================================================

\subsection{Shadow free energies and the Bernoulli generating function}
\label{subsec:thqg-I-shadow-free-energies}
\index{shadow free energy|textbf}
\index{Bernoulli numbers!generating function}

The scalar projection of the Maurer--Cartan element $\Theta_\cA$ onto genus~$g$ is the shadow free energy $F_g(\cA)$. This is the degree-$0$, genus-$g$ component of the shadow algebra $\cA^{\mathrm{sh}} = H_\bullet(\Defcyc^{\mathrm{mod}}(\cA))$ (Definition~\ref{V1-def:shadow-algebra}).

\subsubsection{Definition of the shadow free energy}

\begin{definition}[Shadow free energy]
% label removed: def:thqg-I-shadow-free-energy
\index{shadow free energy!definition|textbf}
For a modular Koszul chiral algebra $\cA$ with universal Maurer--Cartan class $\Theta_\cA \in \mathrm{MC}(\gAmod)$ (Theorem~\ref*{V1-thm:mc2-bar-intrinsic}), the \emph{genus-$g$ shadow free energy} is the degree-zero shadow representation:
\begin{equation}% label removed: eq:thqg-I-Fg-def
F_g(\cA) \;:=\; \int_{\overline{\mathcal{M}}_g} \operatorname{Sh}_{g,0}(\Theta_\cA).
\end{equation}
On the uniform-weight lane (Vol~I, Definition~\ref*{V1-def:scalar-lane}),
genus universality identifies this
with $\kappaChHodge(\cA)\lambda_g^{\mathrm{FP}}$, where
$\lambda_g^{\mathrm{FP}}
=\int_{\overline{\mathcal{M}}_{g,1}}\psi_1^{2g-2}\lambda_g$; for
arbitrary modular Koszul algebras, the unconditional scalar statement
is the genus-$1$ identity $F_1(\cA)=\kappaChHodge(\cA)/24$.
\end{definition}

\begin{remark}[Relation to the shadow algebra]
% label removed: rem:thqg-I-shadow-algebra-relation
The shadow free energy $F_g(\cA)$ is the $(r,g) = (0,g)$ component of the shadow algebra $\cA^{\mathrm{sh}} = H_\bullet(\Defcyc^{\mathrm{mod}}(\cA))$ in the bigrading by degree~$r$ and genus~$g$. The higher-degree shadows at genus $g$ (the cubic shadow $\mathfrak{C}_g$ at $r = 3$, the quartic resonance class $\mathfrak{Q}_g$ at $r = 4$) encode the gravitational interactions beyond the free-energy level. The shadow free energy is the simplest invariant, controlled entirely by the modular characteristic $\kappaChHodge(\cA)$.
\end{remark}

\begin{remark}[Physical interpretation: vacuum energy]
% label removed: rem:thqg-I-vacuum-energy
\index{shadow free energy!physical interpretation}
The shadow free energy $F_g(\cA)$ is the genus-$g$ scalar
vacuum-energy shadow of the twisted holographic theory: it is the
connected, no-insertion coefficient of the algebraic shadow
determinant expansion. Its physical reading as a genus-$g$
cosmological constant requires a chosen bulk comparison datum; it is
not a contribution to an independently constructed path integral. On the
proved scalar lane, $F_g$ is determined by the modular
characteristic $\kappaChHodge(\cA)$ through the Faber--Pandharipande integral;
for arbitrary modular Koszul algebras, the unconditional scalar
statement is the genus-$1$ value $F_1(\cA)=\kappaChHodge(\cA)/24$. Beyond that
lane, the scalar formula at genus~$g \geq 2$ receives a
cross-channel correction for multi-weight algebras
\textup{(}Volume~I, Theorem~\textup{\ref*{V1-thm:multi-weight-genus-expansion}}\textup{)}.
\end{remark}

\subsubsection{The Faber--Pandharipande integral}

The key input from algebraic geometry is the evaluation of the Hodge integral:

\begin{proposition}[Faber--Pandharipande coefficient; \ClaimStatusProvedElsewhere]
\label{prop:thqg-I-faber-pandharipande}
\index{Faber--Pandharipande!coefficient}
\index{Hodge integral}
The Faber--Pandharipande marked Hodge integral is
\begin{equation}% label removed: eq:thqg-I-fp-integral
\lambda_g^{\mathrm{FP}} \;:=\; \int_{\overline{\mathcal{M}}_{g,1}}\psi_1^{2g-2}\lambda_g \;=\; \frac{2^{2g-1} - 1}{2^{2g-1}} \cdot \frac{|B_{2g}|}{(2g)!}\,,
\end{equation}
where $B_{2g}$ is the $2g$-th Bernoulli number. This formula is due to Faber--Pandharipande~\cite{FP00}.
\end{proposition}

\begin{proof}
The marked Hodge integral
$\lambda_g^{\mathrm{FP}}
=\int_{\overline{\mathcal{M}}_{g,1}}\psi_1^{2g-2}\lambda_g$
is computed by localization together with the string equation and
the Witten--Kontsevich recursion for $\psi$-intersections. The
Bernoulli numbers enter through the Mumford--Faber--Pandharipande
evaluation of this integral~\cite{FP00}.
\end{proof}

\begin{remark}[Bernoulli numbers]
% label removed: rem:thqg-I-bernoulli
\index{Bernoulli numbers!values}
The first ten Bernoulli numbers: $B_0 = 1$, $B_1 = -1/2$, $B_2 = 1/6$, $B_4 = -1/30$, $B_6 = 1/42$, $B_8 = -1/30$, $B_{10} = 5/66$, $B_{12} = -691/2730$, $B_{14} = 7/6$, $B_{16} = -3617/510$, $B_{18} = 43867/798$, $B_{20} = -174611/330$. The absolute values grow super-exponentially: $|B_{2g}| \sim 2 \cdot (2g)! / (2\pi)^{2g}$ (Euler's formula). The alternating signs $(-1)^{g+1} B_{2g} > 0$ reflect the alternation of Hodge classes.
\end{remark}

\begin{proposition}[Shadow scalar coefficient in closed form; \ClaimStatusProvedElsewhere{} {\textup{Theorem~\ref{V1-thm:modular-characteristic}}}]
\label{prop:thqg-I-Fg-closed-form}
\index{shadow free energy!closed form}
For any modular Koszul chiral algebra $\cA$ with
generators of the same conformal weight (uniform-weight lane):
\begin{equation}% label removed: eq:thqg-I-Fg-formula
\boxed{ F_g^{\mathrm{sc}}(\cA) \;=\; \kappaChHodge(\cA) \cdot \frac{2^{2g-1}-1}{2^{2g-1}} \cdot \frac{|B_{2g}|}{(2g)!}\,. }
\end{equation}
In particular, on this scalar lane $F_g^{\mathrm{sc}}$ depends only on $g$
and the single scalar $\kappaChHodge(\cA)$. At genus~$1$, the same formula
holds for every modular Koszul algebra.
\end{proposition}

\begin{proof}
This is the combination of genus universality on the
scalar lane
($F_g^{\mathrm{sc}}(\cA)=\kappaChHodge(\cA)\lambda_g^{\mathrm{FP}}$, cf.\
Theorem~\ref{V1-thm:genus-universality}) with the
Faber--Pandharipande evaluation of the Hodge integral
\textup{(}Proposition~\ref{prop:thqg-I-faber-pandharipande}\textup{)}.
\end{proof}

\begin{remark}[Universality]
% label removed: rem:thqg-I-universality
\index{shadow free energy!universality}
On the scalar lane, the genus-$g$ free energy depends only on
the single scalar $\kappaChHodge(\cA)$, regardless of the detailed OPE
structure. The Heisenberg at level $k$, the affine
$\widehat{\mathfrak{sl}}_2$ at level $k$, and the Virasoro at central
charge $c$ all have the same functional form
$F_g^{\mathrm{sc}}(\cA)=\kappaChHodge(\cA)\lambda_g^{\mathrm{FP}}$
\textup{(}uniform-weight\textup{)}, differing only in the
value of $\kappaChHodge(\cA)$. For multi-weight families, this scalar
factorization is only unconditional at genus~$1$. By the
universality of the modular characteristic (Theorem~D), higher-degree
shadows encode the additional nonlinear data beyond this scalar lane.
\end{remark}

\subsubsection{The $\hat{A}$-genus generating function}

\begin{theorem}[Generating function for scalar shadow free energies; \ClaimStatusProvedElsewhere{} {\textup{Theorem~\ref{V1-thm:modular-characteristic}(ii)}}]
\label{thm:thqg-I-generating-function}
\index{A-hat genus@$\hat{A}$-genus!generating function|textbf}
For single-generator algebras \textup{(}and multi-generator
algebras satisfying the factorization condition of
Theorem~\textup{\ref{V1-thm:modular-characteristic}(i)(b))},
the formal generating series of scalar shadow free energies has the
closed-form expression:
\begin{equation}% label removed: eq:thqg-I-gf-ahat
\boxed{ \sum_{g=1}^{\infty} F_g^{\mathrm{sc}}(\cA)\,x^{2g} \;=\; \kappaChHodge(\cA) \cdot \left( \frac{x/2}{\sin(x/2)} - 1 \right). }
\end{equation}
Equivalently, with the convention
$F_0^{\mathrm{sc}}(\cA) = \kappaChHodge(\cA)$:
\begin{equation}% label removed: eq:thqg-I-gf-ahat-alt
\sum_{g=0}^{\infty} F_g^{\mathrm{sc}}(\cA)\,x^{2g} \;=\; \kappaChHodge(\cA) \cdot \frac{x/2}{\sin(x/2)}.
\end{equation}
\end{theorem}

\begin{proof}
We use the standard generating function for Bernoulli numbers. The generating function $\frac{t}{e^t - 1} = \sum_{n=0}^{\infty} \frac{B_n}{n!}\,t^n$ gives: The shadow free energy uses $\lambda_g^{\mathrm{FP}} = \frac{2^{2g-1}-1}{2^{2g-1}}\cdot\frac{|B_{2g}|}{(2g)!}$ with absolute values of Bernoulli numbers. The generating function is therefore
\begin{align}
\sum_{g=1}^{\infty} F_g^{\mathrm{sc}}(\cA)\,x^{2g} &\;=\; \kappaChHodge(\cA) \sum_{g=1}^{\infty} \frac{2^{2g-1}-1}{2^{2g-1}} \cdot \frac{|B_{2g}|}{(2g)!}\,x^{2g} \notag\\
&\;=\; \kappaChHodge(\cA) \left( \frac{x/2}{\sin(x/2)} - 1 \right), % label removed: eq:thqg-I-gf-derivation
\end{align}
where the second equality uses the identity
\begin{equation}% label removed: eq:thqg-I-cosecant-bernoulli
\frac{x/2}{\sin(x/2)} \;=\; 1 + \sum_{g=1}^{\infty} \frac{(2^{2g-1}-1)|B_{2g}|}{2^{2g-1}\cdot(2g)!}\,x^{2g},
\end{equation}
obtainable from the Bernoulli generating function by the substitution $t = ix$ and the relation $\sin(x/2) = (e^{ix/2} - e^{-ix/2})/(2i)$.

The Bernoulli generating function gives $\frac{t}{e^t - 1} = 1 - t/2 + \sum_{g=1}^\infty \frac{B_{2g}}{(2g)!} t^{2g}$. Setting $t = ix$:
\[
\frac{ix}{e^{ix} - 1} \;=\; 1 - \frac{ix}{2} + \sum_{g=1}^\infty \frac{B_{2g}}{(2g)!}(ix)^{2g} \;=\; 1 - \frac{ix}{2} + \sum_{g=1}^\infty \frac{(-1)^g B_{2g}}{(2g)!} x^{2g}.
\]
Since $(-1)^g B_{2g} = -|B_{2g}|$ for $g \geq 1$, we have
\[
\frac{ix}{e^{ix}-1} \;=\; 1 - \frac{ix}{2} - \sum_{g=1}^\infty \frac{|B_{2g}|}{(2g)!}x^{2g}.
\]
Combining with the identity $\frac{ix}{e^{ix}-1} + \frac{ix}{2} = \frac{ix/2}{\tanh(ix/2)} = \frac{x/2}{\tan(x/2)}$, and using $\frac{x/2}{\sin(x/2)} = \frac{x/2}{\tan(x/2)} + \frac{x/2\,(1-\cos(x/2))}{\sin(x/2)\cos(x/2)}$, we arrive at~\eqref{V1-eq:thqg-I-cosecant-bernoulli} by the standard manipulation of the partial fraction expansion $\csc(z) = 1/z + 2z\sum_{n=1}^\infty (-1)^n/(z^2 - n^2\pi^2)$.

More directly: the function $\frac{x/2}{\sin(x/2)}$ is even and analytic at $x = 0$ with value $1$. Its Taylor expansion $\sum_{g=0}^\infty a_{2g} x^{2g}$ has coefficients $a_{2g} = \frac{(2^{2g-1}-1)|B_{2g}|}{2^{2g-1}(2g)!}$, which is verified by comparing with the standard expansion of $x/\sin(x) = \sum_{g=0}^\infty \frac{(2^{2g}-2)\,|B_{2g}|}{(2g)!}\, x^{2g}$ and substituting $x \mapsto x/2$.
\end{proof}

\begin{remark}[Two generating functions]
% label removed: rem:thqg-I-two-gfs
On the proved scalar lane, there are two natural generating functions
in play. The \emph{algebraic} generating function
$\sum_g F_g^{\mathrm{sc}}(\cA)\,x^{2g}
= \kappaChHodge(\cA)\bigl(x/(2\sin(x/2)) - 1\bigr)$ uses the
real variable $x$ and converges for $|x| < 2\pi$. The \emph{physical}
generating function $\sum_g F_g^{\mathrm{sc}}\,\hbar^g$ uses the genus-counting
parameter~$\hbar$ and converges for $|\hbar| < 4\pi^2$. The relation
between the two is $\hbar = x^2$: the algebraic variable $x$ is the
``square root'' of the genus-counting parameter.
\end{remark}

\begin{remark}[Connection to the $\hat{A}$-genus]
% label removed: rem:thqg-I-ahat-connection
\index{A-hat genus@$\hat{A}$-genus!connection to shadow free energy}
On the proved scalar lane, the generating function $x/(2\sin(x/2))$ is
the square root of the Todd genus $\mathrm{td}(x) = x/(1-e^{-x})$ and
the Wick-rotated $\hat{A}$-genus $\hat{A}(x) = (x/2)/\sinh(x/2)$:
$x/(2\sin(x/2)) = \hat{A}(ix)$. The rotation passes from the Lorentzian
$\hat{A}$-genus (which controls index theory on spin manifolds) to the
Euclidean generating function (which controls genus integrals on moduli
spaces). The shadow free energy thus has a natural interpretation as a
``Wick-rotated $\hat{A}$-genus'' of the modular characteristic.
\end{remark}

\subsubsection{Absolute convergence from Bernoulli asymptotics}

\begin{theorem}[Absolute convergence of the scalar generating function; \ClaimStatusProvedHere]
\label{thm:thqg-I-absolute-convergence}
\index{absolute convergence!shadow free energy}
For any modular Koszul chiral algebra $\cA$ with generators of
uniform conformal weight, the
scalar genus-counting series
$\sum_{g=1}^{\infty} |F_g^{\mathrm{sc}}(\cA)|\,|\hbar|^g$
converges absolutely for $|\hbar| < 4\pi^2$. The generating function
$\sum_{g=1}^{\infty} F_g^{\mathrm{sc}}(\cA)\,\hbar^g$ extends to a meromorphic
function on $\mathbb{C}$ via the closed form
$\kappaChHodge(\cA) \cdot (\frac{\sqrt{\hbar}/2}{\sin(\sqrt{\hbar}/2)} - 1)$.
For multi-weight algebras, the unconditional statement is
the genus-$1$ scalar identity
$F_1^{\mathrm{sc}}(\cA)=\kappaChHodge(\cA)/24$;
at $g \geq 2$ the full free energy decomposes as
$F_g(\cA)=F_g^{\mathrm{sc}}(\cA)
+\delta F_g^{\mathrm{cross}}(\cA)$
\textup{(}Volume~I, Theorem~\textup{\ref*{V1-thm:multi-weight-genus-expansion}}\textup{)}.
\end{theorem}

\begin{proof}
On the uniform-weight lane,
$F_g^{\mathrm{sc}}(\cA)=\kappaChHodge(\cA)\lambda_g^{\mathrm{FP}}$ for all $g \geq 1$
\textup{(}Volume~I, Theorem~D\textup{)}. The asymptotic behavior of
the Faber--Pandharipande coefficients is controlled by Euler's formula for the Bernoulli numbers,
$|B_{2g}| = 2 \cdot \zeta(2g) \cdot (2g)! / (2\pi)^{2g}$, where
$\zeta(2g) \to 1$ as $g \to \infty$. Substituting into $\lambda_g^{\mathrm{FP}} = \frac{2^{2g-1}-1}{2^{2g-1}} \cdot \frac{|B_{2g}|}{(2g)!}$:
\begin{align}
|F_g^{\mathrm{sc}}(\cA)| &\;=\; |\kappaChHodge(\cA)| \cdot \frac{2^{2g-1}-1}{2^{2g-1}} \cdot \frac{|B_{2g}|}{(2g)!} \notag \\
&\;=\; |\kappaChHodge(\cA)| \cdot \underbrace{\frac{2^{2g-1}-1}{2^{2g-1}}}_{\to\,1} \cdot \underbrace{\frac{2\,\zeta(2g)}{(2\pi)^{2g}}}_{\sim\,2/(2\pi)^{2g}} \notag \\
&\;\sim\; \frac{2|\kappaChHodge(\cA)|}{(2\pi)^{2g}}. % label removed: eq:thqg-I-Fg-asymptotics
\end{align}
The two-sided bound for $g \geq 2$ is:
\[
\frac{2|\kappaChHodge(\cA)|}{(2\pi)^{2g}} \;\leq\; |F_g^{\mathrm{sc}}(\cA)| \;\leq\; \frac{2|\kappaChHodge(\cA)|}{(2\pi)^{2g}} \cdot \frac{\zeta(2g)}{1 - 2^{1-2g}} \;\leq\; \frac{4|\kappaChHodge(\cA)|}{(2\pi)^{2g}}\,.
\]
The ratio test now gives the convergence radius explicitly. The consecutive ratio is
\[
\frac{|F_{g+1}|}{|F_g|}
\;=\;
\frac{(2^{2g+1}-1)}{(2^{2g-1}-1)}
\cdot
\frac{|B_{2g+2}|}{|B_{2g}|}
\cdot
\frac{(2g)!}{(2g+2)!}
\cdot
\frac{2^{2g-1}}{2^{2g+1}}\,.
\]
By Euler's formula, $|B_{2g+2}|/|B_{2g}| = \zeta(2g+2)/\zeta(2g) \cdot (2g+2)(2g+1)/(2\pi)^2$.
Therefore $|F_{g+1}/F_g| = \zeta(2g+2)/\zeta(2g) \cdot \frac{2^{2g+1}-1}{2^{2g+1}} / \frac{2^{2g-1}-1}{2^{2g-1}} \cdot 1/(2\pi)^2 \to 1/(2\pi)^2$ as $g \to \infty$. The Cauchy--Hadamard radius of convergence is therefore
\[
R \;=\; \lim_{g \to \infty} |F_g/F_{g+1}| \;=\; (2\pi)^2 \;=\; 4\pi^2 \;\approx\; 39.48\,.
\]
This is not an artifact of a bound: the radius is \emph{exact}, determined by the first pole of the closed-form continuation $\kappaChHodge(\cA) \cdot (\frac{\sqrt{\hbar}/2}{\sin(\sqrt{\hbar}/2)} - 1)$. The poles occur at the zeros of $\sin(\sqrt{\hbar}/2)$,
i.e.\ at $\sqrt{\hbar}/2 = n\pi$, giving
$\hbar_n = 4\pi^2 n^2$ for $n = 1, 2, 3, \ldots$ The nearest pole to the origin is $\hbar_1 = 4\pi^2$, confirming $R = 4\pi^2$.

For the residues: writing $s(\hbar)=\sin(\sqrt{\hbar}/2)$, one
has $s'(\hbar_n)=(-1)^n/(8\pi n)$ at $\hbar_n = 4\pi^2 n^2$, so
\[
\Res_{\hbar = 4\pi^2 n^2}
\kappaChHodge(\cA) \cdot \frac{\sqrt{\hbar}/2}{\sin(\sqrt{\hbar}/2)}
\;=\; 8(-1)^n \kappaChHodge(\cA)n^2\pi^2.
\]
\end{proof}

\begin{remark}[Factorially divergent vs.\ convergent]
% label removed: rem:thqg-I-factorial-convergent
\index{factorial growth!absence of}
In standard perturbative quantum field theory, the genus-$g$
contribution typically grows as $(2g)!$, reflecting the factorial
number of Feynman diagrams. On the scalar lane, the
shadow free energy satisfies
$F_g^{\mathrm{sc}}(\cA) \sim 2\kappaChHodge(\cA)/(2\pi)^{2g}$, so the
series $\sum F_g^{\mathrm{sc}}\hbar^g$ is convergent, not merely asymptotic. The
reason is that the scalar shadow free energy there is an integral over
the compact moduli space $\overline{\mathcal{M}}_g$ of a smooth
cohomology class, not a sum over all Feynman diagrams.
\end{remark}

\subsubsection{Explicit values through genus 10}

\begin{proposition}[Shadow free energies through genus $10$ on the scalar lane; \ClaimStatusProvedHere]
% label removed: prop:thqg-I-Fg-values
\index{shadow free energy!explicit values}
If $\cA$ has generators of uniform conformal weight, then using the
Faber--Pandharipande coefficients
\textup{(}Proposition~\ref{prop:thqg-I-faber-pandharipande}\textup{)}:
\begin{align}
F_1^{\mathrm{sc}}(\cA) &= \frac{\kappaChHodge(\cA)}{24}, % label removed: eq:thqg-I-F1\\
F_2^{\mathrm{sc}}(\cA) &= \frac{7\kappaChHodge(\cA)}{5760}, % label removed: eq:thqg-I-F2\\
F_3^{\mathrm{sc}}(\cA) &= \frac{31\kappaChHodge(\cA)}{967680}, % label removed: eq:thqg-I-F3\\
F_4^{\mathrm{sc}}(\cA) &= \frac{127\kappaChHodge(\cA)}{154828800}, % label removed: eq:thqg-I-F4\\
F_5^{\mathrm{sc}}(\cA) &= \frac{73\kappaChHodge(\cA)}{3503554560}, % label removed: eq:thqg-I-F5\\
F_6^{\mathrm{sc}}(\cA) &= \frac{1414477\kappaChHodge(\cA)}{2678117105664000}, % label removed: eq:thqg-I-F6\\
F_7^{\mathrm{sc}}(\cA) &= \frac{8191\kappaChHodge(\cA)}{612141052723200}, % label removed: eq:thqg-I-F7\\
F_8^{\mathrm{sc}}(\cA) &= \frac{16931177\kappaChHodge(\cA)}{49950709902213120000}, % label removed: eq:thqg-I-F8\\
F_9^{\mathrm{sc}}(\cA) &= \frac{5749691557\kappaChHodge(\cA)}{669659197233029971968000}, % label removed: eq:thqg-I-F9\\
F_{10}^{\mathrm{sc}}(\cA) &= \frac{91546277357\kappaChHodge(\cA)}{420928638260761696665600000}. % label removed: eq:thqg-I-F10
\end{align}
For arbitrary modular Koszul algebras, the unconditional identity is
$F_1^{\mathrm{sc}}(\cA)=\kappaChHodge(\cA)/24$; higher-genus formulas require the stated
modular-Koszul hypotheses.
\end{proposition}

\begin{proof}
Direct substitution of the Bernoulli numbers into the
formula
$F_g^{\mathrm{sc}}(\cA)=\kappaChHodge(\cA) \frac{2^{2g-1}-1}{2^{2g-1}} \frac{|B_{2g}|}{(2g)!}$
\textup{(}uniform-weight\textup{)}.

At genus $g = 1$: $F_1^{\mathrm{sc}}(\cA) = \kappaChHodge(\cA) \cdot \frac{1}{2} \cdot \frac{1/6}{2} = \kappaChHodge(\cA)/24$.

At genus $g = 2$: $(2^3{-}1)/2^3 = 7/8$, $|B_4| = 1/30$, $4! = 24$. $F_2^{\mathrm{sc}}(\cA) = \kappaChHodge(\cA) \cdot \frac{7}{8} \cdot \frac{1/30}{24} = \frac{7\kappaChHodge(\cA)}{5760}$.

At genus $g = 3$: $|B_6| = 1/42$, $6! = 720$, $(2^5 - 1)/2^5 = 31/32$. $F_3^{\mathrm{sc}}(\cA) = \kappaChHodge(\cA) \cdot \frac{31}{32} \cdot \frac{1/42}{720} = \frac{31\kappaChHodge(\cA)}{32 \cdot 30240} = \frac{31\kappaChHodge(\cA)}{967680}$.

At genus $g = 4$: $|B_8| = 1/30$, $8! = 40320$, $(2^7 - 1)/2^7 = 127/128$. $F_4^{\mathrm{sc}}(\cA) = \kappaChHodge(\cA) \cdot \frac{127}{128} \cdot \frac{1/30}{40320} = \frac{127\kappaChHodge(\cA)}{128 \cdot 1209600} = \frac{127\kappaChHodge(\cA)}{154828800}$.

At genus $g = 5$: $|B_{10}| = 5/66$, $10! = 3628800$, $(2^9 - 1)/2^9 = 511/512$. $F_5^{\mathrm{sc}}(\cA) = \kappaChHodge(\cA) \cdot \frac{511}{512} \cdot \frac{5/66}{3628800} = \kappaChHodge(\cA) \cdot \frac{511 \cdot 5}{512 \cdot 66 \cdot 3628800} = \frac{2555\kappaChHodge(\cA)}{122624409600} = \frac{73\kappaChHodge(\cA)}{3503554560}$.

The remaining values are computed similarly.
\end{proof}

\begin{table}[ht]
\centering
\caption{Scalar shadow free energies $F_g^{\mathrm{sc}}(\cA)=\kappaChHodge(\cA)\lambda_g^{\mathrm{FP}}$ for three standard families, $g = 1, \ldots, 5$.}
% label removed: tab:thqg-I-Fg-specializations
\index{shadow free energy!table}
\renewcommand{\arraystretch}{1.5}
{\footnotesize
\begin{tabular}{|c|c|c|c|}
\hline
$\boldsymbol{g}$ & $\mathcal{H}_k$, $\kappaChHodge(\mathcal{H}_k) = k$ & $\widehat{\mathfrak{sl}}_{2,k}$, $\kappaChHodge(\widehat{\mathfrak{sl}}_{2,k}) = \frac{3(k+2)}{4}$ & $\mathrm{Vir}_c$, $\kappaChHodge(\mathrm{Vir}_c) = \frac{c}{2}$ \\
\hline
$1$ & $\dfrac{k}{24}$ & $\dfrac{k+2}{32}$ & $\dfrac{c}{48}$ \\[6pt]
\hline
$2$ & $\dfrac{7k}{5760}$ & $\dfrac{7(k+2)}{7680}$ & $\dfrac{7c}{11520}$ \\[6pt]
\hline
$3$ & $\dfrac{31k}{967680}$ & $\dfrac{31(k+2)}{1290240}$ & $\dfrac{31c}{1935360}$ \\[6pt]
\hline
$4$ & $\dfrac{127k}{154828800}$ & $\dfrac{127(k+2)}{206438400}$ & $\dfrac{127c}{309657600}$ \\[6pt]
\hline
$5$ & $\dfrac{73k}{3503554560}$ & $\dfrac{73(k+2)}{4671406080}$ & $\dfrac{73c}{7007109120}$ \\[6pt]
\hline
\end{tabular}}
\end{table}

\begin{remark}[Numerical magnitudes]
% label removed: rem:thqg-I-numerical-magnitudes
\index{shadow free energy!numerical estimates}
For the Heisenberg at $k = 1$: $F_1 \approx 0.0417$, $F_2 \approx 1.22 \times 10^{-3}$, $F_3 \approx 3.20 \times 10^{-5}$, $F_4 \approx 8.20 \times 10^{-7}$, $F_5 \approx 2.08 \times 10^{-8}$. The free energies decrease super-exponentially, with ratio $F_{g+1}/F_g \to 1/(2\pi)^2 \approx 0.0253$. By genus $10$, $F_{10} \approx 2.17 \times 10^{-16}$. The geometric decay with ratio $1/(4\pi^2) \approx 0.0253$ means that the genus expansion is an extremely rapidly converging series, even at $\hbar = 1$.
\end{remark}

\begin{proposition}[Cumulative partial sums; \ClaimStatusProvedHere]
% label removed: prop:thqg-I-partial-sums
\index{shadow free energy!partial sums}
For the Heisenberg at $k = 1$ and $\hbar = 1$, the partial sums of the genus expansion are:
\begin{align}
S_1 &= F_1 = 0.04167\,, \notag \\
S_2 &= S_1 + F_2 = 0.04288\,, \notag \\
S_3 &= S_2 + F_3 = 0.04291\,, \notag \\
S_4 &= S_3 + F_4 = 0.04291\,, \notag \\
S_5 &= S_4 + F_5 = 0.04291\,, % label removed: eq:thqg-I-partial-sums
\end{align}
converging to $\kappaChHodge(\mathcal H_1) \cdot (1/\mathrm{sinc}(1/2) - 1) = (1/(2\sin(1/2)) - 1) \cdot 1 \approx 0.04291$.
\end{proposition}

\begin{proof}
Direct numerical evaluation.
\end{proof}


% ======================================================================
% SUBSECTION 3: CONVERGENCE OF THE SCALAR SHADOW SERIES
% ======================================================================

\subsection{Scalar convergence and genuswise finiteness}
\label{subsec:thqg-I-convergence-grav}
\index{shadow partition series!scalar convergence|textbf}

\subsubsection{Definition of the scalar shadow genus series}

\begin{definition}[Scalar shadow genus series]
% label removed: def:thqg-I-grav-partition
\index{shadow partition series!scalar definition|textbf}
For a holographic modular Koszul datum
$\mathcal{H}(T) = (\cA, \cA^!, \mathcal{C}, r(z), \Theta_\cA,
\nabla^{\mathrm{hol}})$ and a formal genus-counting
parameter~$\hbar$, the \emph{scalar shadow genus series} is
\begin{equation}% label removed: eq:thqg-I-Z-grav-def
Z_{\mathrm{sh}}^{\mathrm{scal}}(T;\,\hbar)
\;:=\;
\sum_{g=0}^{\infty} \hbar^g
\int_{\overline{\mathcal{M}}_{g}}
\operatorname{Sh}_{g,0}(\Theta_\cA)
\;=\;
\sum_{g=0}^{\infty} F_g(\cA)\,\hbar^g\,.
\end{equation}
After a bulk comparison datum has been fixed, the same scalar series is
written
\begin{equation}% label removed: eq:thqg-I-Z-grav-scalar
Z_{\mathrm{grav}}^{\mathrm{scal}}(T;\,\hbar)
\;:=\;
Z_{\mathrm{sh}}^{\mathrm{scal}}(T;\,\hbar)\,.
\end{equation}
On the scalar lane the closed form
$\kappaChHodge(\cA) \cdot \frac{\sqrt{\hbar}/2}{\sin(\sqrt{\hbar}/2)}$
follows from Theorem~\ref{thm:thqg-I-generating-function}. For
arbitrary modular Koszul algebras, the unconditional scalar statement
is $F_1(\cA)=\kappaChHodge(\cA)/24$. The definition is the
degree-$0$ trace of the shadow Maurer--Cartan element; it is not a full
partition function over bulk metrics or all shadow degrees.
\end{definition}

\begin{remark}[Scalar versus full]
% label removed: rem:thqg-I-scalar-vs-full
The scalar shadow series captures only the scalar trace of the
corresponding cohomological modular-functor package. The degree-summed
shadow series also contains the higher-degree contributions from the
bar-intrinsic support packet: the spectral discriminant $\Delta_\cA$
\textup{(}degree~$2$ spectral channel\textup{)}, the cubic shadow
$\mathfrak{C}$ \textup{(}degree~$3$\textup{)}, the quartic resonance
class $\mathfrak{Q}$ \textup{(}degree~$4$\textup{)}, and all higher
shadows. For Gaussian algebras \textup{(}support shadow depth
\(r^\Theta_{\max}=2\)\textup{)}, the scalar and degree-summed shadow series
coincide.
\end{remark}

\begin{remark}[Genus-$0$ normalization]
% label removed: rem:thqg-I-genus-0
\index{shadow partition series!genus-0 normalization}
The genus-$0$ contribution $F_0(\cA) = \kappaChHodge(\cA)$ is a convention: it
supplies the completed scalar shadow series
$Z_{\mathrm{sh}}^{\mathrm{scal}}(T;\,\hbar)
= \kappaChHodge(\cA) + \sum_{g \geq 1} F_g \hbar^g$. This scalar series
is distinct from the degree-summed shadow series and from any physical
analytic partition function defined genuswise by sewing integrals. In
the semiclassical limit $\hbar \to 0$, only the genus-$0$ term survives.
\end{remark}

\subsubsection{Convergence from HS-sewing and Bernoulli asymptotics}

\begin{theorem}[Algebraic shadow finiteness; \ClaimStatusProvedHere]
\label{thm:thqg-I-perturbative-finiteness}
\index{perturbative finiteness!main theorem|textbf}
Let $\cA$ be a modular Koszul chiral algebra satisfying the HS-sewing condition \textup{(}Theorem~\textup{\ref{thm:general-hs-sewing})}. Then:
\begin{enumerate}[label=\textup{(\roman*)}]
\item \emph{Genus-by-genus finiteness.} For each $g \geq 1$, the genus-$g$ shadow amplitude $\int_{\overline{\mathcal{M}}_g} \operatorname{Sh}_{g,n}(\Theta_\cA)$ is finite at every degree~$n$.
\item \emph{Completed scalar shadow series on the proved scalar lane.} If $\cA$ has generators of uniform conformal weight, then the reduced series $\sum_{g=1}^{\infty} F_g(\cA)\,\hbar^g$ converges absolutely for $|\hbar| < 4\pi^2$ and equals
\[
\kappaChHodge(\cA)\left(
\frac{\sqrt{\hbar}/2}{\sin(\sqrt{\hbar}/2)}-1\right).
\]
The completed scalar series, with genus-$0$ term included, is
\[
Z_{\mathrm{sh}}^{\mathrm{scal}}(\cA;\,\hbar)
= \kappaChHodge(\cA)\cdot
\frac{\sqrt{\hbar}/2}{\sin(\sqrt{\hbar}/2)}.
\]
For arbitrary modular Koszul algebras, the universal scalar statement is the finite genus-$1$ term $F_1(\cA)=\kappaChHodge(\cA)/24$.
\item \emph{Degree-summed shadow series beyond the scalar lane.} At finite support shadow depth \(r^\Theta_{\max}(\cA)<\infty\), the degree-summed fixed-genus shadow amplitude is a finite degree sum. At infinite support shadow depth the bar-intrinsic MC element defines a class in the completed degree product; a numerical degree-summed trace requires an absolute shadow-degree estimate. If $\cA$ is Gaussian, then the degree-summed shadow series equals the scalar shadow series and therefore converges for $|\hbar| < 4\pi^2$ with meromorphic continuation. Beyond the Gaussian case these hypotheses give no lower bound for the degree-summed non-scalar genus radius.
\item \emph{No bar-complex counterterm.} No counterterm is introduced in the algebraic shadow complex at fixed genus and degree. The Maurer--Cartan equation $D_\cA\Theta_\cA + \tfrac{1}{2}[\Theta_\cA,\Theta_\cA] = 0$ (Theorem~\ref*{V1-thm:mc2-bar-intrinsic}) is an identity of the bar differential, and the Fulton--MacPherson compactification supplies the boundary strata on which the identity is evaluated.
\end{enumerate}
\end{theorem}

\begin{proof}
\emph{Part~(i).} Fix genus~$g$ and degree~$n$. The shadow representation $\operatorname{Sh}_{g,n}(\Theta_\cA) \in H^*(\overline{\mathcal{M}}_{g,n})$ is a cohomology class on a compact orbifold. Integration over $\overline{\mathcal{M}}_{g,n}$ is a continuous linear functional on this finite-dimensional space, hence bounded. The nontrivial content is that the class itself is well defined: this follows from the bar-intrinsic construction of $\Theta_\cA$ (Theorem~\ref*{V1-thm:mc2-bar-intrinsic}), which gives $\Theta_\cA := D_\cA - d_0 \in \mathrm{MC}(\gAmod)$ without convergence issues.

At the chain level, the integral is computed by the degree-cutoff lemma (Lemma~\ref{V1-lem:degree-cutoff}): the graph sum through degree~$r$ at genus~$g$ involves only finitely many graphs $\Gamma$ with $|\Gamma| \leq r$ edges and loop genus $\leq g$. Each graph amplitude $\ell_\Gamma = |\Aut(\Gamma)|^{-1} \int_{\overline{\mathcal{M}}_{g(\Gamma),n(\Gamma)}} \omega_\Gamma$ is a period integral on a compact moduli space, hence finite.

\emph{Part~(ii).} On the scalar lane this is
Theorem~\ref{thm:thqg-I-absolute-convergence}: the Bernoulli
asymptotics $|F_g(\cA)| \sim 2|\kappaChHodge(\cA)|/(2\pi)^{2g}$ give absolute
convergence for $|\hbar| < 4\pi^2$, and analytic continuation via the
closed form extends to all of~$\mathbb{C}$ as a meromorphic function.
For arbitrary modular Koszul algebras, the genus-$1$ clause is
Theorem~\ref{V1-thm:genus-universality}\textup{(iii)}.

\emph{Part~(iii).} For the degree-summed shadow series, each
fixed-genus contribution is graded by the support packet of the
bar-intrinsic MC element.
Finite support shadow depth gives a finite degree sum. Infinite support
shadow depth gives the
inverse-limit class of Theorem~\ref{V1-thm:recursive-existence} in the
weight-completed complex, not an absolute trace norm. In the Gaussian
case \(r^\Theta_{\max}=2\), all higher-degree corrections vanish, so the
degree-summed shadow series equals the scalar shadow series from
Part~(ii). A lower bound for the degree-summed non-scalar genus radius
is not a consequence of these hypotheses.

\emph{Part~(iv).} The absence of algebraic shadow counterterms is the combined consequence of three features:
\begin{enumerate}[label=\textup{(\alph*)}]
\item \emph{Compactness.} The Fulton--MacPherson compactification $\overline{\mathrm{FM}}_n(\Sigma_g)$ is a compact manifold with corners. All integrals are over compact domains.
\item \emph{Algebraic closure.} The Maurer--Cartan equation is an identity of the bar differential ($D_\cA^2 = 0$), not a perturbative expansion. There is no shadow counterterm to add because the equation holds at all genera without truncation.
\item \emph{Degree cutoff.} At each genus~$g$, only finitely many graph topologies contribute (Lemma~\ref{V1-lem:degree-cutoff}). The number of contributing graphs at genus $g$ and degree $r$ is bounded by a combinatorial function of $(g,r)$, and each graph amplitude is a finite integral.
\end{enumerate}
\end{proof}

\subsubsection{Comparison with standard QFT divergences}

\begin{remark}[Contrast with perturbative quantum gravity in $d = 4$]
% label removed: rem:thqg-I-contrast-4d
\index{UV divergence!standard quantum gravity}
In standard four-dimensional perturbative quantum gravity (Einstein--Hilbert action), the genus-$g$ contribution to the graviton scattering amplitude diverges for $g \geq 2$ as
\begin{equation}% label removed: eq:thqg-I-4d-divergence
\mathcal{A}_g \;\sim\; \Lambda_{\mathrm{UV}}^{2(g-1)} \int_{\mathcal{M}_g} |\Omega_g|^2 \cdot R^{2g},
\end{equation}
where $\Lambda_{\mathrm{UV}}$ is the UV cutoff and $R$ is the Riemann curvature. In the shadow setting, the algebraic analogue $\mathcal{A}_g^{\mathrm{sh}} = \int_{\overline{\mathcal{M}}_g} \operatorname{Sh}_{g,n}(\Theta_\cA)$ is finite at every genus because the integration domain is compact, the integrand is a cohomology class, and the Maurer--Cartan equation is a Ward identity on the bar complex. A physical scattering-amplitude statement requires the separate BV realisation and chain-level topologisation hypotheses.
\end{remark}

\begin{remark}[Contrast with string theory]
% label removed: rem:thqg-I-contrast-string
\index{string theory!UV finiteness}
In string theory, genus-$g$ amplitudes are finite individually (after integration over $\overline{\mathcal{M}}_g$), but the genus expansion $\sum_g g_s^{2g-2} \mathcal{A}_g$ is asymptotic: it diverges factorially as $\mathcal{A}_g \sim (2g)! \cdot \alpha'^{-g}$. On the uniform-weight scalar shadow lane, the diagonal genus expansion converges: $F_g(\cA) \sim \kappaChHodge(\cA)/(2\pi)^{2g}$ decays exponentially. The statement is about the scalar cohomological shadow; the multi-weight cross-channel correction series and the physical BV genus series are not identified with it here.
\end{remark}

\subsubsection{The role of Koszulness in finiteness}

\begin{proposition}[Koszulness and finiteness; \ClaimStatusProvedHere]
% label removed: prop:thqg-I-koszulness-finiteness
\index{Koszulness!role in finiteness}
For a modular Koszul chiral algebra $\cA$, the bar-intrinsic support
packet \(\Theta_\cA^{\leq r}\) has the following finiteness properties:
\begin{enumerate}[label=\textup{(\roman*)}]
\item At each degree $r$ and genus $g$, the graph sum
 \begin{equation}% label removed: eq:thqg-I-graph-sum-finite
 \Theta_\cA^{(g,r)} \;=\; \sum_{\substack{\Gamma \in \Gmod(g,r) \\ \Gamma \text{ connected}}} \frac{1}{|\Aut(\Gamma)|} \int_{\overline{\mathcal{M}}_{g(\Gamma),n(\Gamma)}} \omega_\Gamma
 \end{equation}
 is a finite sum over a finite set of graphs.
\item At fixed total degree $r$, the transferred cyclic minimal model contributes only arities $n\leq r+1$; if \(r^\Theta_{\max}(\cA)<\infty\), the nonzero support components lie in degrees \(n\leq r^\Theta_{\max}(\cA)\).
\item If \(r^\Theta_{\max}(\cA)<\infty\), then the bar-intrinsic support packet has no degree component above \(r^\Theta_{\max}(\cA)\).  This is stronger than eventual vanishing of the obstruction classes \(o_{r+1}\), which records only \(H^2\)-extendability.
\end{enumerate}
\end{proposition}

\begin{proof}
Part~(i) is the degree-cutoff lemma (Lemma~\ref{V1-lem:degree-cutoff}). The set $\Gmod(g,r)$ of connected modular graphs with loop genus $g$ and total degree $\leq r$ is finite: the number of vertices $v$, edges $e$, and external legs $n$ satisfy $v - e + n = 1 - g$ (the Euler characteristic), $2e + n \leq rv$ (each vertex has at most $r$ half-edges), and $g = e - v + 1 + \sum_{v'} g_{v'}$ where $g_{v'}$ is the genus decoration at each vertex. These constraints bound $v, e, n$ in terms of $g$ and $r$.

Part~(ii) is the arity filtration in the HPL tree formula. A tree contributing to total degree $r$ has at most $r+1$ leaves; finite support shadow depth gives the additional support cutoff \(r^\Theta_{\max}(\cA)\) in the chosen bar-intrinsic packet. Chiral Koszulness supplies bar concentration, not vanishing of the transferred tower.

Part~(iii) is the definition of support depth: \(\Theta_{\cA,n}=0\) for every \(n>r^\Theta_{\max}(\cA)\).  The obstruction equation \(o_{n+1}=0\) only says that a degree-\(n\) truncation can be extended; it does not kill the corresponding \(H^1\)-lift coordinate.
\end{proof}

\begin{proposition}[Effective graph count; \ClaimStatusProvedHere]
\label{prop:thqg-I-graph-count}
\index{graph sum!effective count}
The number of connected modular graphs contributing to $\Theta_\cA^{(g,r)}$ is bounded by
\begin{equation}% label removed: eq:thqg-I-graph-bound
|\Gmod(g,r)| \;\leq\; C(g,r) \;:=\; \sum_{v=1}^{3g-3+r} \binom{v(r-1)/2 + g - 1}{g} \cdot r^v,
\end{equation}
which is finite for fixed $(g,r)$. As $r$ varies at fixed $g$, this
bound is a combinatorial bound, not a polynomial degree-growth
estimate.
\end{proposition}

\begin{proof}
A connected graph with $v$ vertices and loop genus $g$ has $e = v + g - 1$ edges. Each vertex has valence at most $r$ (by the degree cutoff), so $2e + n \leq rv$, giving $n \leq rv - 2(v+g-1) = v(r-2) - 2g + 2$. The condition $n \geq 0$ requires $v \geq (2g-2)/(r-2)$ (for $r \geq 3$). The number of graphs with $v$ labeled vertices and $e$ edges is at most $\binom{v(v-1)/2}{e}$ (choosing edges from all possible pairs); the symmetry group reduces this, but the crude bound suffices. The factor $r^v$ accounts for the vertex labeling (choice of operation at each vertex from $\{m_2, m_3, \ldots, m_r\}$).
\end{proof}

\subsubsection{The structure of finiteness at each genus}

\begin{proposition}[Genus-$g$ amplitude decomposition; \ClaimStatusProvedHere]
% label removed: prop:thqg-I-genus-g-decomposition
\index{genus-$g$ amplitude!decomposition}
The genus-$g$ shadow amplitude decomposes as
\begin{equation}% label removed: eq:thqg-I-genus-g-decomp
\operatorname{Tr}_g(\Theta_\cA)
\;:=\;
F_g^{\mathrm{scal}}(\cA)
\;+\;
\sum_{2\le r\le r^\Theta_{\max}(\cA)} F_g^{(r)}(\cA),
\end{equation}
where $F_g^{\mathrm{scal}}$ is the scalar (degree-$0$) contribution and
$F_g^{(r)}$ is the degree-$r$ correction. On the scalar
lane, $F_g^{\mathrm{scal}}(\cA) = \kappaChHodge(\cA) \cdot \lambda_g^{\mathrm{FP}}$ for
all $g \geq 1$; for arbitrary modular Koszul algebras, the
unconditional scalar statement is $F_1^{\mathrm{scal}}(\cA)=\kappaChHodge(\cA)/24$
\textup{(}$g = 1$ is unconditional at all weights\textup{)}. % For
a Gaussian algebra (\(r^\Theta_{\max}=2\)), the sum is empty and only the
scalar term survives. For a Lie-type algebra (\(r^\Theta_{\max}=3\), canonical affine branch), there
is a single cubic correction $F_g^{(3)}$. For a contact-type algebra
\(r^\Theta_{\max}=4\), canonical contact branch), the support
corrections stop at degree~$4$. For a mixed
algebra (\(r^\Theta_{\max}=\infty\)), the right hand side is the completed
degree expansion; its numerical trace is an absolutely convergent
series precisely when the scalar functional satisfies
$\sum_{r \geq 2}|F_g^{(r)}(\cA)|<\infty$.
\end{proposition}

\begin{proof}
The bar-intrinsic support packet decomposes by degree:
$\Theta_\cA = \sum_{r \geq 0} \Theta_\cA^{(r)}$ where
$\Theta_\cA^{(r)}$ is the degree-$r$ component. The genus-$g$
projection at degree $0$ gives $F_g^{\mathrm{scal}}$; on the
scalar lane this equals
$\kappaChHodge(\cA)\lambda_g^{\mathrm{FP}}$, while for arbitrary modular Koszul
algebras the universal scalar statement is the genus-$1$ identity
$F_1^{\mathrm{scal}}(\cA)=\kappaChHodge(\cA)/24$. At degree $r \geq 2$, the projection
$\int_{\overline{\mathcal{M}}_g} \operatorname{Sh}_{g,r}(\Theta_\cA)$
involves a graph sum with vertex labels from the degree-$r$ operations,
and each term is finite by the degree-cutoff argument. The vanishing of
\(F_g^{(r)}\) for \(r>r^\Theta_{\max}(\cA)\) follows from the
support-depth definition in the selected bar-intrinsic packet.
\end{proof}

\begin{remark}[Pixton ideal generation from the MC element]
\label{rem:vol2-pixton-from-mc}
\index{Pixton relations!from MC element}
On the semisimple locus of the shadow CohFT, the MC-descended
tautological relations generate the full Pixton ideal
$I_g \subset R^*(\overline{\cM}_{g,n})$
\textup{(}Volume~I, Theorem~\textup{\ref*{V1-thm:pixton-from-mc-semisimple})}.
By Givental--Teleman reconstruction, semisimplicity
of the shadow CohFT $\Omega^\cA_{g,n}$ enables Teleman's classification,
and the MC-tautological descent produces the full set of Pixton relations
as an output of the genus recursion. Thus MC descent gives not only
finite shadow amplitudes but also the full tautological constraint ideal on
$\overline{\cM}_g$.
\end{remark}

% ======================================================================
% SUBSECTION 4: THE HEISENBERG PARTITION FUNCTION AS PROTOTYPE
% ======================================================================

\subsection{The Heisenberg partition function as prototype}
\label{subsec:thqg-I-heisenberg-prototype}
\index{Heisenberg!partition function|textbf}

The Heisenberg algebra is the archetype: Gaussian support shadow depth
\(r^\Theta_{\max}=2\), closed-form solvability via second quantization,
and explicit Fredholm determinant formulas at all genera.

\subsubsection{The Heisenberg at genus $g$}

\begin{theorem}[Heisenberg determinant-line shadow; \ClaimStatusProvedHere]
\label{thm:heisenberg-sewing}
\index{Heisenberg!genus-$g$ partition function}
For the rank-$1$ Heisenberg algebra $\mathcal{H}_k$ on a genus-$g$ surface $\Sigma_g$ with period matrix $\Omega \in \mathbb{H}_g$ \textup{(}the Siegel upper half-space\textup{)}, the genus-$g$ scalar shadow determined by the double-pole coefficient $k$ is the Fredholm determinant
\begin{equation}% label removed: eq:thqg-I-heisenberg-Zg
\boxed{ Z_g(\mathcal{H}_k;\,\Omega) \;=\; \bigl(\det\operatorname{Im}\Omega\bigr)^{-k/2} \cdot \bigl(\det{}'_\zeta \Delta_{\Sigma_g}\bigr)^{-k/2}, }
\end{equation}
where $\det{}'_\zeta$ denotes the zeta-regularized determinant of the scalar Laplacian $\Delta_{\Sigma_g}$ on $\Sigma_g$ \textup{(}excluding the zero mode\textup{)}. The shadow free energy extracted from this is
\begin{equation}% label removed: eq:thqg-I-heisenberg-Fg
F_g(\mathcal{H}_k) \;=\; k \cdot \frac{2^{2g-1}-1}{2^{2g-1}} \cdot \frac{|B_{2g}|}{(2g)!}\,,
\end{equation}
confirming the general formula~\eqref{V1-eq:thqg-I-Fg-formula} at $\kappaChHodge(\mathcal{H}_k) = k$.
\end{theorem}

\begin{proof}
By Wick's theorem, all amplitudes of the Heisenberg algebra factorize into propagator pairings. The propagator on a genus-$g$ surface is the Green's function
\begin{equation}% label removed: eq:thqg-I-green-genus-g
G(z,w;\,\Omega) \;=\; \log|E(z,w;\,\Omega)| - \frac{(\operatorname{Im}\int_w^z \omega)^T (\operatorname{Im}\Omega)^{-1} (\operatorname{Im}\int_w^z \omega)}{2},
\end{equation}
where $E(z,w;\,\Omega)$ is the prime form and $\omega = (\omega_1, \ldots, \omega_g)^T$ is the normalized basis of holomorphic differentials. The prime form is a section of $K^{-1/2} \boxtimes K^{-1/2}$ (not $K^{+1/2}$).

The scalar shadow is the determinant-line power of the kinetic
operator:
\[
Z_g \;=\; (\det{}'_\zeta \bar{\partial}^*\bar{\partial})^{-k/2} \;=\; (\det\operatorname{Im}\Omega)^{-k/2}\, (\det{}'_\zeta \Delta)^{-k/2},
\]
where the factorization into period-matrix and Laplacian contributions is the Belavin--Knizhnik factorization~\cite{BK86}. The Polyakov formula identifies $\det{}'_\zeta \Delta$ with the norm of the Mumford determinant line on $\overline{\mathcal{M}}_g$; its curvature is the Hodge class $\lambda_1=c_1(\det\mathbb{E})$. The scalar genus-$g$ number is obtained by the marked Faber--Pandharipande pairing
\[
\lambda_g^{\mathrm{FP}}
=\int_{\overline{\mathcal{M}}_{g,1}}\psi_1^{2g-2}\lambda_g.
\]
On the Heisenberg scalar lane this gives
\[
F_g(\mathcal{H}_k) = k\,\lambda_g^{\mathrm{FP}}.
\]
\end{proof}

\begin{remark}[The Belavin--Knizhnik factorization]
% label removed: rem:thqg-I-belavin-knizhnik
\index{Belavin--Knizhnik!factorization}
The factorization $\det{}'_\zeta(\bar\partial^*\bar\partial) = \det(\operatorname{Im}\Omega) \cdot \det{}'_\zeta(\Delta)$ separates the ``holomorphic'' contribution (the period matrix, which depends on the complex structure of $\Sigma_g$) from the ``analytic'' contribution (the Laplacian determinant, which depends on the conformal class). The period matrix factor $(\det\operatorname{Im}\Omega)^{-k/2}$ is the modular-covariant piece that transforms under $\mathrm{Sp}(2g,\mathbb{Z})$, while the Laplacian factor is modular-invariant. Together they give the conformal block for the Heisenberg algebra on $\Sigma_g$.
\end{remark}

\subsubsection{The Fredholm determinant formula}

\begin{proposition}[Fredholm determinant representation; \ClaimStatusProvedElsewhere{} {\textup{Theorem~\ref{V1-thm:heisenberg-one-particle-sewing}}}]
% label removed: prop:thqg-I-fredholm
\index{Fredholm determinant!Heisenberg|textbf}
The genus-$g$ scalar shadow admits a Fredholm determinant representation:
\begin{equation}% label removed: eq:thqg-I-fredholm-det
Z_g(\mathcal{H}_k;\,\Omega) \;=\; \det(1 - K_g)^{-k},
\end{equation}
where $K_g$ is the trace-class sewing operator on the one-particle Bergman space $A^2(\partial D)$, with eigenvalues $q_1^{n_1} \cdots q_{3g-3}^{n_{3g-3}}$ indexed by the sewing moduli $(q_1, \ldots, q_{3g-3})$. The trace-class norm satisfies
\begin{equation}% label removed: eq:thqg-I-Kg-trace-norm
\|K_g\|_1 \;=\; \prod_{i=1}^{3g-3} \frac{q_i}{1-q_i} \;<\; \infty \quad\text{for all } q_i \in (0,1).
\end{equation}
\end{proposition}

\begin{proof}
The pair-of-pants amplitude for the Heisenberg is the second quantization $\Gamma(R)$ of the one-particle Bergman restriction $R \colon A^2(D_{\mathrm{out}}) \to A^2(D_{\mathrm{in},1}) \otimes A^2(D_{\mathrm{in},2})$ (Theorem~\ref{V1-thm:heisenberg-one-particle-sewing}). The one-particle restriction has matrix entries $R_{n,m} = \frac{k}{2\pi i} \oint\oint z^{-n-1}w^{-m-1}G_{\mathrm{pants}}(z,w)\,dz\,dw$, satisfying $|R_{n,m}| \leq C\,q^{(n+m)/2}$ from the collar length $t = -\log q$. Hence $R$ is Hilbert--Schmidt and $\Gamma(R)$ is trace class.

The composition of $3g - 3$ sewing operators gives the genus-$g$ operator $K_g = R_1 \circ \cdots \circ R_{3g-3}$ on the one-particle space. For one free boson, second quantization gives $\Tr\,\Gamma(K_g)=\det(1-K_g)^{-1}$; the scalar shadow uses the determinant-line exponent $k$. The trace-class norm bound follows from $\|K_g\|_1 \leq \prod_{i=1}^{3g-3} \|R_i\|_1 = \prod_{i=1}^{3g-3} q_i/(1-q_i)$.
\end{proof}

\begin{remark}[The Bergman space]
% label removed: rem:thqg-I-bergman
\index{Bergman space}
The one-particle Bergman space $A^2(\partial D)$ consists of holomorphic functions on the annular collar around a sewing circle, square-integrable with respect to the Bergman kernel. An orthonormal basis is $\{z^n / \sqrt{2\pi}\}_{n \geq 1}$, and the one-particle sewing operator in this basis has matrix elements $K_{nm} = \delta_{nm} q^n$. The Fredholm determinant is therefore $\det(1 - K) = \prod_{n=1}^\infty (1 - q^n)$, giving the Dedekind eta function upon including the zero-point energy.
\end{remark}

\subsubsection{Comparison with the Dedekind eta function}

\begin{example}[Genus $1$: the Dedekind eta function]
% label removed: ex:thqg-I-genus1-eta
\index{Dedekind eta!Heisenberg partition function}
At genus $g = 1$, the period matrix is $\Omega = \tau \in \mathbb{H}$ and the sewing parameter is $q = e^{2\pi i\tau}$. The rank-$1$ holomorphic oscillator determinant is
\begin{equation}% label removed: eq:thqg-I-genus1-fredholm
Z_{1,\mathrm{osc}}^{\mathrm{hol}}(\mathcal{H}_k;\,\tau)
\;=\;
\prod_{n=1}^{\infty} (1 - q^n)^{-1}
\;=\;
q^{1/24}\eta(\tau)^{-1}.
\end{equation}
The $k$-weighted scalar determinant-line shadow is obtained by taking
the determinant-line power with exponent $k$, adjoining the anti-holomorphic
copy, and inserting the zero-mode factor:
\begin{equation}% label removed: eq:thqg-I-genus1-scalar-eta
Z_{1}^{\mathrm{scal}}(\mathcal{H}_k;\,\tau)
\;=\;
(\operatorname{Im}\tau)^{-k/2}\,\bigl|\eta(\tau)\bigr|^{-2k}.
\end{equation}
Here $\eta(\tau) = q^{1/24}\prod_{n=1}^{\infty}(1 - q^n)$ is the Dedekind eta function. The shadow free energy at genus $1$ is $F_1(\mathcal{H}_k) = k/24$, matching~\eqref{V1-eq:thqg-I-F1} with $\kappaChHodge(\mathcal H_k) = k$.

The modular characteristic is the coefficient of the determinant-line
curvature: the $q$-expansion of
$\log \eta(\tau) = \frac{2\pi i\tau}{24} - \sum_{n=1}^\infty \log(1-q^n)$
begins with $\frac{2\pi i \tau}{24}$, and the genus-$1$ shadow
free energy is the constant term of $-k\log|\eta(\tau)|$, namely
$k/24$.
\end{example}

\begin{example}[Genus $2$: explicit formula]
% label removed: ex:thqg-I-genus2
\index{Heisenberg!genus-2 partition function}
At genus $g = 2$, the period matrix $\Omega = \begin{psmallmatrix} \tau_1 & \tau_{12} \\ \tau_{12} & \tau_2 \end{psmallmatrix} \in \mathbb{H}_2$ and the scalar shadow is
\begin{equation}% label removed: eq:thqg-I-genus2
Z_2(\mathcal{H}_k;\,\Omega) \;=\; \bigl(\det\operatorname{Im}\Omega\bigr)^{-k/2} \cdot \bigl(\det{}'_\zeta \Delta_{\Sigma_2}\bigr)^{-k/2}.
\end{equation}
The shadow free energy is $F_2(\mathcal{H}_k) = 7k/5760$, confirming~\eqref{V1-eq:thqg-I-F2}.

The Fredholm determinant at genus $2$ involves the sewing of three pairs of pants along three circles. In the Schottky uniformization, $\Sigma_2 = (\mathbb{P}^1 \setminus \{6 \text{ discs}\}) / \langle \gamma_1, \gamma_2, \gamma_3 \rangle$ with three sewing parameters $q_1, q_2, q_3$. The one-particle sewing operator is $K_2 = R_1 \circ R_2 \circ R_3$, and $\det(1 - K_2) = \prod_{n_1,n_2,n_3 \geq 1} (1 - q_1^{n_1} q_2^{n_2} q_3^{n_3})$. The shadow free energy is obtained by integrating over the Siegel modular variety $\mathcal{M}_2 = \mathbb{H}_2 / \mathrm{Sp}(4,\mathbb{Z})$.
\end{example}

\begin{example}[Genus $3$]
% label removed: ex:thqg-I-genus3-heisenberg
\index{Heisenberg!genus-3 partition function}
At genus $g = 3$, the surface $\Sigma_3$ has $3g - 3 = 6$ sewing parameters. The scalar shadow Fredholm determinant is $\det(1 - K_3)^{-k}$ where $K_3$ is a trace-class operator on the one-particle Bergman space. The shadow free energy is $F_3(\mathcal{H}_k) = 31k/967680$. Numerically at $k = 1$: $F_3 \approx 3.20 \times 10^{-5}$, which is already three orders of magnitude smaller than $F_1 = 1/24 \approx 0.042$.
\end{example}

\subsubsection{One-particle Bergman reduction}

\begin{remark}[One-particle reduction principle]
% label removed: rem:thqg-I-one-particle
\index{Bergman space!one-particle reduction|textbf}
The Heisenberg scalar shadow factorizes completely into one-particle data. The one-particle sewing operator $T_q$ on the Bergman space $A^2(\partial D)$ has eigenvalues $q^n$ ($n = 1, 2, 3, \ldots$), and the determinant-line shadow is
\[
Z_g(\mathcal{H}_k) \;=\; \det(1 - T_{q_1} \circ \cdots \circ T_{q_{3g-3}})^{-k} \;=\; \exp\!\left( k \sum_{m=1}^{\infty} \frac{1}{m} \Tr\bigl(T_{q_1} \cdots T_{q_{3g-3}}\bigr)^m \right).
\]
The shadow free energy $F_g$ is the constant term of $\log Z_g$ in the expansion in sewing moduli. The one-particle reduction is specific to the Heisenberg (Gaussian support shadow depth \(r^\Theta_{\max}=2\)): for non-Gaussian algebras, the multi-particle sector contributes data not determined by the one-particle space.
\end{remark}

\begin{proposition}[Exponential representation; \ClaimStatusProvedHere]
% label removed: prop:thqg-I-exponential-rep
\index{Fredholm determinant!exponential representation}
The Fredholm determinant admits the exponential representation
\begin{equation}% label removed: eq:thqg-I-exponential-fredholm
\log \det(1 - K_g)^{-k} \;=\; k \sum_{m=1}^\infty \frac{1}{m} \Tr(K_g^m) \;=\; k \sum_{m=1}^\infty \frac{1}{m} \sum_{\mathbf{n} \in \mathbb{Z}_{\geq 1}^{3g-3}} \prod_{i=1}^{3g-3} q_i^{m n_i},
\end{equation}
which converges absolutely for all $q_i \in (0,1)$. The shadow free energy $F_g$ is extracted as the coefficient of the leading term in the $\mathbf{q}$-expansion, integrated over $\overline{\mathcal{M}}_g$ with the Weil--Petersson measure.
\end{proposition}

\begin{proof}
The Fredholm determinant identity $\log\det(1 - K) = \Tr\log(1 - K) = -\sum_{m=1}^\infty \frac{1}{m}\Tr(K^m)$ holds for trace-class $K$ with $\|K\|_1 < 1$. For $q_i \in (0,1)$, the eigenvalues of $K_g$ are $\prod_{i=1}^{3g-3} q_i^{n_i}$ with $n_i \geq 1$, all less than $1$. The trace $\Tr(K_g^m) = \sum_{\mathbf{n}} \prod q_i^{mn_i}$ converges absolutely because $\sum_{\mathbf{n}} \prod q_i^{n_i} = \prod_{i=1}^{3g-3} q_i/(1-q_i) < \infty$.
\end{proof}

\begin{remark}[Rank-$d$ Heisenberg]
% label removed: rem:thqg-I-rank-d-heisenberg
\index{Heisenberg!rank-$d$}
For the rank-$d$ Heisenberg algebra $\mathcal{H}_k^{\otimes d}$ (i.e., $d$ independent free bosons at level $k$), the partition function is the $d$-th power of the rank-$1$ answer:
\begin{equation}% label removed: eq:thqg-I-rank-d
Z_g(\mathcal{H}_k^{\otimes d};\,\Omega) \;=\; Z_g(\mathcal{H}_k;\,\Omega)^d \;=\; \bigl(\det\operatorname{Im}\Omega\bigr)^{-dk/2} \cdot \bigl(\det{}'_\zeta\Delta_{\Sigma_g}\bigr)^{-dk/2}.
\end{equation}
The modular characteristic is $\kappaChHodge(\mathcal{H}_k^{\otimes d}) = dk$, and the shadow free energy is $F_g(\mathcal{H}_k^{\otimes d}) = dk \cdot \lambda_g^{\mathrm{FP}}$. This is consistent with the independent sum factorization (Proposition~\ref{V1-prop:independent-sum-factorization}): for $\cA = \cA_1 \oplus \cA_2$ with vanishing mixed OPE, $\kappaChHodge(\cA) = \kappaChHodge(\cA_1) + \kappaChHodge(\cA_2)$.
\end{remark}


% ======================================================================
% SUBSECTION 5: FINITENESS FOR THE STANDARD LANDSCAPE
% ======================================================================

\subsection{Finiteness for the standard landscape}
\label{subsec:thqg-I-standard-landscape}
\index{standard landscape!finiteness|textbf}

\subsubsection{Verification of hypotheses (H1)--(H2)}

The HS-sewing criterion (Theorem~\ref{thm:general-hs-sewing}) requires two conditions:
\begin{enumerate}[label=\textbf{(H\arabic*)}]
\item \label{item:thqg-I-H1} \emph{Subexponential sector growth:} $\log \dim H_n = o(n)$.
\item \label{item:thqg-I-H2} \emph{Polynomial OPE growth:} $|C^{c,k}_{a,i;\,b,j}| \leq K(a+b+c+1)^N$.
\end{enumerate}
The Hilbert norm is part of the sewing datum; \ref{item:thqg-I-H2}
is not invariant under an arbitrary rescaling of homogeneous bases.

\begin{proposition}[Standard families satisfy HS-sewing; \ClaimStatusProvedElsewhere{} {\textup{Theorem~\ref{thm:general-hs-sewing}}}]
\label{prop:thqg-I-standard-hs}
\index{HS-sewing!standard families}
Every algebra in the following list satisfies \ref{item:thqg-I-H1} and \ref{item:thqg-I-H2}:
\begin{enumerate}[label=\textup{(\alph*)}]
\item Heisenberg $\mathcal{H}_k$ (any rank, any level $k \neq 0$).
\item Affine Kac--Moody $\widehat{\fg}_k$ ($\fg$ simple, $k \neq -h^\vee$).
\item Virasoro $\mathrm{Vir}_c$ (any $c \in \mathbb{C}$).
\item $\beta\gamma$ system (any central charge).
\item Principal $\mathcal{W}$-algebras $\mathcal{W}_k(\fg)$ ($\fg$ simple, generic $k$).
\item Lattice VOAs $V_\Lambda$ ($\Lambda$ even integral).
\item Free fermions $\psi_1, \ldots, \psi_n$ (any rank).
\end{enumerate}
\end{proposition}

\begin{proof}
For each family, we verify \ref{item:thqg-I-H1} and \ref{item:thqg-I-H2} by direct computation. The key estimates were computed in Computations~\ref{comp:thqg-I-hs-heisenberg}--\ref{comp:thqg-I-hs-fermion}.

\emph{(a) Heisenberg.} $\dim H_n = p(n)$, $\log p(n) \sim \pi\sqrt{2n/3} = o(n)$. OPE degree $N = 1$.

\emph{(b) Affine Kac--Moody.} $\dim H_n$ has subexponential growth (Kac--Peterson~\cite{KP84}). OPE degree $N = 1$.

\emph{(c) Virasoro.} $\dim H_n = p(n)$, $\log p(n) = o(n)$. OPE degree $N = 2$.

\emph{(d) $\beta\gamma$.} $\dim H_n = p(n)^2$, $\log \dim H_n = o(n)$. OPE degree $N = 1$.

\emph{(e) $\mathcal{W}_N$.} Sector growth subexponential with exponent $O(n^{(N-1)/N})$. OPE degree $N_{\max} = 2N$.

\emph{(f) Lattice VOA.} $\dim H_n = |\Lambda^*/\Lambda|\cdot p_r(n)$; $\log p_r(n) = O(\sqrt{n})$. OPE degree $\leq r$.

\emph{(g) Free fermions.} $\dim H_n$ grows at most as $2^n$ (number of subsets of $n$ fermion modes), but the conformal-weight grading gives polynomial growth. OPE degree $N = 0$.
\end{proof}

\subsubsection{Detailed verification for each family}

\begin{proposition}[Heisenberg: detailed HS verification; \ClaimStatusProvedHere]
\label{prop:thqg-I-heisenberg-detail}
\index{Heisenberg!HS verification}
For the rank-$1$ Heisenberg $\mathcal{H}_k$ at level $k \neq 0$:
\begin{enumerate}[label=\textup{(\roman*)}]
\item The fundamental OPE coefficient is
$C^{n-1,\lambda}_{n,\mu;\,0,\mathrm{vac}}
= k\,\delta_{|\lambda|-|\mu|,1}$ for the mode
$a_{-1}|0\rangle$. In the sewing Wick basis the complete
multiplication tensor satisfies
$|C^{c,\lambda}_{a,\mu;\,b,\nu}|
\leq C_k(a+b+c+1)$.
\item The HS norm at total weight $n$ satisfies
 \begin{equation}% label removed: eq:thqg-I-heisenberg-exact-hs
 \sum_{\substack{a+b+c = n}} \|m^c_{a,b}\|_{\HS}^2
 \;\leq\;
 C_k^2(n+1)^2
 \sum_{\substack{a+b+c=n}} p(a)p(b)p(c),
 \end{equation}
and this bound gives the HS-sewing condition after multiplication by
$q^n$ and summation over $n$.
\end{enumerate}
\end{proposition}

\begin{proof}
The Heisenberg OPE is computed by Wick's theorem: each contracted pair
of modes gives $k\,m\,\delta_{m+n,0}$. The sewing Wick basis absorbs the
multinomial factors in the Gaussian norm, leaving a linear mode-number
bound at total weight $a+b+c$. Squaring the coefficient bound and
summing over the three partition bases gives the displayed estimate.
\end{proof}

\begin{proposition}[Affine Kac--Moody: detailed HS verification; \ClaimStatusProvedHere]
% label removed: prop:thqg-I-affine-detail
\index{affine Kac--Moody!HS verification}
For $\widehat{\fg}_k$ with $\fg$ simple of dimension $d$, rank $r$, and dual Coxeter number $h^\vee$, at non-critical level $k \neq -h^\vee$:
\begin{enumerate}[label=\textup{(\roman*)}]
\item The Kac--Peterson character formula gives $\chi_k(\tau) = \sum_{n \geq 0} (\dim H_n) q^n = q^{-c_{\text{eff}}/24} \cdot \Theta_{k+h^\vee,\fg}(\tau) / \eta(\tau)^d$, where $\Theta_{k+h^\vee,\fg}$ is the affine theta function. The dimension $\dim H_n$ is dominated by $\eta(\tau)^{-d}$, giving $\log \dim H_n \sim \pi\sqrt{2dn/3}$.
\item The OPE of currents $J^a_m J^b_n = k m \delta^{ab}\delta_{m+n,0} + f^{ab}{}_c J^c_{m+n}$ has structure constants bounded by $\max(|k|, \|f\|_\infty) \cdot n$ for mode number $n$, which is polynomial of degree $N = 1$.
\item The HS sewing sum satisfies
 \begin{equation}% label removed: eq:thqg-I-affine-hs-refined
 \sum_{a,b,c \geq 0} q^{a+b+c} \|m^c_{a,b}\|_{\HS}^2 \;\leq\; (|k| + \|f\|_\infty)^2 \cdot \prod_{n=1}^\infty (1-q^n)^{-2d} \cdot P(q),
 \end{equation}
 where $P(q) = \sum_{n \geq 0} n^2 (\dim H_n)^2 q^n$ is a convergent power series for $q \in (0,1)$.
\end{enumerate}
\end{proposition}

\begin{proof}
(i) is the Kac--Peterson asymptotic~\cite{KP84}. (ii) follows from the affine Lie algebra commutation relations. (iii) combines the estimates from (i) and (ii) with the general HS bound~\eqref{V1-eq:thqg-I-hs-criterion-bound}.
\end{proof}

\begin{proposition}[Virasoro: detailed HS verification; \ClaimStatusProvedHere]
% label removed: prop:thqg-I-virasoro-detail
\index{Virasoro!HS verification}
For $\mathrm{Vir}_c$ at central charge $c$:
\begin{enumerate}[label=\textup{(\roman*)}]
\item The Verma module at weight $h$ has $\dim H_n = p(n)$ (partitions of $n$, corresponding to states $L_{-n_1}L_{-n_2}\cdots L_{-n_k}|h\rangle$ with $n_1 + \cdots + n_k = n$).
\item The OPE $T(z)T(w) \sim c/(2(z-w)^4) + 2T(w)/(z-w)^2 + \partial T(w)/(z-w)$ has a quartic pole. The resulting structure constants satisfy $|C^c_{a,b}| \leq (|c|/2 + 2) \cdot (a+b+1)^2$, which is polynomial of degree $N = 2$.
\item The refined HS bound:
 \begin{equation}% label removed: eq:thqg-I-virasoro-refined-hs
 \sum_{a,b,c \geq 0} q^{a+b+c} \|m^c_{a,b}\|_{\HS}^2 \;\leq\; (|c|/2 + 2)^2 \sum_{n \geq 0} n^4 p(n)^3 q^n.
 \end{equation}
 The series $\sum n^4 p(n)^3 q^n$ converges for all $q \in (0,1)$ because $p(n)^3 \leq e^{3\pi\sqrt{2n/3} + O(\log n)}$ grows subexponentially while $q^n$ decays exponentially. Concretely, for $q = 0.5$: the $n = 100$ term is $100^4 \cdot p(100)^3 \cdot 0.5^{100} \approx 10^8 \cdot (1.9 \times 10^8)^3 \cdot 10^{-30} \approx 10^{2}$, and the series converges.
\end{enumerate}
\end{proposition}

\begin{proof}
(i) is standard Virasoro representation theory. (ii) follows from the mode algebra $[L_m, L_n] = (m-n)L_{m+n} + \frac{c}{12}m(m^2-1)\delta_{m+n,0}$, which gives $|C| \leq |c|/12 \cdot m^3 + |m-n| \cdot |L_{m+n}|$; the polynomial bound is obtained by tracking the maximal power of conformal weight in iterated commutators. (iii) follows from substitution into~\eqref{V1-eq:thqg-I-hs-criterion-bound}.
\end{proof}

\subsubsection{Convergence radii and analytic continuation}

\begin{proposition}[Convergence radius of the scalar shadow closed form; \ClaimStatusProvedHere]
% label removed: prop:thqg-I-convergence-radii
\index{convergence radius!scalar shadow closed form}
The scalar shadow closed form
\[
Z_{\mathrm{sh}}^{\mathrm{scal}}(\cA;\,\hbar)
= \kappaChHodge(\cA)\cdot
\frac{\sqrt{\hbar}/2}{\sin(\sqrt{\hbar}/2)}
\]
has:
\begin{enumerate}[label=\textup{(\roman*)}]
\item \emph{Radius of convergence in $\hbar$:} $R = 4\pi^2 \approx 39.48$.
\item \emph{First singularity:} A pole of order $1$ at $\hbar = 4\pi^2$, corresponding to the first zero of $\sin(\sqrt{\hbar}/2)$ at $\sqrt{\hbar} = 2\pi$.
\item \emph{Analytic continuation:} $Z^{\mathrm{scal}}$ extends to a meromorphic function on $\mathbb{C}$ with poles at $\hbar = 4\pi^2 n^2$, $n \in \mathbb{Z}_{> 0}$.
\end{enumerate}
\end{proposition}

\begin{proof}
The function $f(\hbar) = \frac{\sqrt{\hbar}/2}{\sin(\sqrt{\hbar}/2)}$ is analytic at $\hbar = 0$ (removable singularity: $f(0) = 1$). The root test gives $\limsup_{g \to \infty} |a_g|^{1/g} = 1/(2\pi)^2 = 1/(4\pi^2)$, so the radius of convergence is $4\pi^2$. The poles at $\hbar = 4\pi^2 n^2$ are simple with residues $8(-1)^n\kappaChHodge(\cA)n^2\pi^2$.
\end{proof}

\begin{remark}[Scalar pole scales]
% label removed: rem:thqg-I-hagedorn
\index{Hagedorn singularity}
The poles at $\hbar = 4\pi^2 n^2$ are poles of the scalar meromorphic
shadow. At $\hbar = 4\pi^2$, the scalar Taylor series reaches its
first pole and is replaced by the meromorphic continuation. No
physical phase transition or full non-scalar singularity is implied
without the additional BV saddle analysis.
\end{remark}

\subsubsection{Admissible levels and singular fibers}

\begin{remark}[Admissible levels]
% label removed: rem:thqg-I-admissible
\index{admissible level!finiteness}
At admissible levels $k = -h^\vee + p/q$ (with $p,q$ coprime positive integers and $p \geq h^\vee$), the universal vertex algebra $V_k(\fg)$ has null vectors and the simple quotient $L_k(\fg)$ differs from $V_k(\fg)$. The HS-sewing criterion applies to $V_k(\fg)$ unconditionally (by Proposition~\ref{prop:thqg-I-standard-hs}). For the simple quotient $L_k(\fg)$, the Kac--Wakimoto character formula~\cite{KW88} still supplies the subexponential sector-growth half of the HS criterion, and the admissible-level bar analysis gives weightwise finite-page control on the non-degenerate lane together with reduced/nilradical obstruction surfaces on the degenerate lane. What is \emph{not} proved here is the full HS-sewing criterion for $L_k(\fg)$ itself: polynomial OPE growth and compatible Hilbert--Schmidt norm control on the quotient are additional analytic input not established on this surface.

The modular characteristic is $\kappaChHodge(V_k(\fg)) = td/(2h^\vee)$ with $t = k + h^\vee$, which vanishes at critical level $k = -h^\vee$. At critical level, the scalar shadow free energies all vanish: $F_g = 0$ for all $g \geq 1$. The Feigin--Frenkel center emerges as the replacement for the scalar genus package.
\end{remark}

\begin{remark}[Singular fibers]
% label removed: rem:thqg-I-singular-fibers
\index{singular fiber!modular characteristic}
As the level $k$ varies, the modular characteristic $\kappaChHodge(k) = td/(2h^\vee)$ with $t = k + h^\vee$ traces a linear function of $k$. The critical level $k = -h^\vee$ is the unique singular fiber where $\kappaChHodge$ vanishes. Near the singular fiber, $\kappaChHodge \to 0$ and the scalar shadow free energies $F_g \to 0$ uniformly: $|F_g(k)| \leq C |k + h^\vee| / (2\pi)^{2g}$ for all $g$. The scalar class collapses to $\Theta_\cA^{\min} = 0$ at $k = -h^\vee$; this does not determine the higher-degree tower.
\end{remark}

\begin{remark}[Critical level and the Feigin--Frenkel center]
% label removed: rem:thqg-I-ff-center
\index{Feigin--Frenkel center}
At critical level $k = -h^\vee$, the generic Sugawara tensor $T(z) = \frac{1}{2(k+h^\vee)} \sum_a :J^a(z) J^a(z):$ has a pole, hence the Sugawara construction is \emph{undefined}. The Feigin--Frenkel theorem identifies the center $Z(V_{-h^\vee}(\fg)) \cong \mathrm{Fun}\,\mathrm{Op}_{\fg^\vee}(\mathbb{D})$ with the algebra of functions on the space of $\fg^\vee$-opers on the formal disc. The \emph{scalar} shadow free energy vanishes at this level because $\kappaChHodge(V_{-h^\vee}(\fg)) = 0$ (forcing $m_0 = 0$, hence the scalar lane is uncurved); the higher-degree components of the MC element are not determined by $\kappaChHodge=0$ alone, and $F_1 = 0$ does not imply $F_g = 0$ at higher degree. The Feigin--Frenkel center provides a replacement for the scalar genus package: the ``classical'' genus-$g$ invariant is the space of $\fg^\vee$-opers on $\Sigma_g$, which is finite-dimensional and does not require the HS-sewing machinery.
\end{remark}

\subsubsection{Summary table}

\begin{table}[ht]
\centering
\caption{Perturbative finiteness: standard landscape}
\label{tab:thqg-I-finiteness-summary}
\index{perturbative finiteness!summary table}
\renewcommand{\arraystretch}{1.4}
{\footnotesize
\begin{tabular}{|l|c|c|c|c|c|}
\hline
\textbf{Family $\cA$} & $\boldsymbol{\dim H_n}$ & \textbf{OPE $N$} & $\boldsymbol{\kappaChHodge(\cA)}$ & $\boldsymbol{r^\Theta_{\max}}$ & \textbf{HS Status} \\
\hline
$\mathcal{H}_k$ & $p(n)$ & $1$ & $k$ & $2$ & \textbf{Proved} \\
\hline
$\widehat{\fg}_k$ & $\sim e^{\pi\sqrt{2dn/3}}$ & $1$ & $\frac{(k+h^\vee)d}{2h^\vee}$ & $3$ & \textbf{Proved} \\
\hline
$\mathrm{Vir}_c$ & $p(n)$ & $2$ & $c/2$ & $\infty$ & \textbf{Proved} \\
\hline
$\beta\gamma$ & $p(n)^2$ & $1$ & $1$ & $4$ & \textbf{Proved} \\
\hline
$\mathcal{W}_N$ & $N$-colored & $2N$ & $c(H_N - 1)$ & $\infty$ & \textbf{Proved} \\
\hline
$V_\Lambda$ & $|\Lambda^*/\Lambda|\cdot p_r(n)$ & $r$ & $\rank(\Lambda)$ & $2$ & \textbf{Proved} \\
\hline
Free fermions & polynomial & $0$ & $n/4$ & $2$ & \textbf{Proved} \\
\hline
\end{tabular}}
\end{table}

\begin{remark}[Exhaustiveness]
% label removed: rem:thqg-I-exhaustiveness
\index{standard landscape!exhaustiveness}
The seven families in Table~\ref{tab:thqg-I-finiteness-summary} exhaust the ``standard landscape'' of finitely strongly generated chiral algebras with simple Lie symmetry. Finiteness extends to:
\begin{enumerate}[label=\textup{(\alph*)}]
\item \emph{Tensor products:} $\cA \otimes \cB$ satisfies HS-sewing if both $\cA$ and $\cB$ do (the HS norm of the tensor product is the product of HS norms).
\item \emph{Cosets:} $\cA / \cB$ (the coset of $\cB \subset \cA$) satisfies HS-sewing if $\cA$ does (the sector dimensions can only decrease).
\item \emph{Orbifolds:} $\cA^G$ (orbifold by a finite group $G$) satisfies HS-sewing if $\cA$ does (the twisted sectors contribute additional terms, but all with subexponential growth).
\item \emph{Non-principal $\mathcal{W}$-algebras:} $\mathcal{W}_k(\fg, f)$ for nilpotent $f$ satisfies HS-sewing at generic $k$ (the DS reduction does not increase OPE degree).
\end{enumerate}
\end{remark}


% ======================================================================
% SUBSECTION 6: IMPLICATIONS FOR 3D QUANTUM GRAVITY
% ======================================================================

\subsection{Implications for 3d quantum gravity}
\label{subsec:thqg-I-3d-gravity}
\index{3d quantum gravity|textbf}
\index{BTZ black hole!scalar shadow comparison}

Three-dimensional gravity with negative cosmological constant $\Lambda = -1/\ell^2$ is the physical arena where the perturbative finiteness theorem has the most direct consequences. The boundary theory is a 2d CFT with central charge determined by the Brown--Henneaux formula $c = 3\ell/(2G_N)$.

\subsubsection{The Brown--Henneaux identification}

\begin{setup}[Brown--Henneaux central charge]
\label{setup:thqg-I-brown-henneaux}
\index{Brown--Henneaux!central charge}
The asymptotic symmetry algebra of AdS$_3$ gravity at spatial infinity is two copies of the Virasoro algebra with central charge
\begin{equation}% label removed: eq:thqg-I-brown-henneaux
c \;=\; \frac{3\ell}{2G_N}\,,
\end{equation}
where $\ell$ is the AdS radius and $G_N$ is the three-dimensional Newton constant~\cite{BH86}. In the twisted holographic setting, $\cA = \mathrm{Vir}_c$ is the boundary chiral algebra, with modular characteristic $\kappaChHodge(\cA) = c/2 = 3\ell/(4G_N)$.
\end{setup}

\begin{remark}[Semiclassical limit]
% label removed: rem:thqg-I-semiclassical
\index{semiclassical limit}
The semiclassical limit is $G_N \to 0$ with $\ell$ fixed, which gives
$c \to \infty$ and $\kappaChHodge(\mathrm{Vir}_c)=c/2 \to \infty$. On the Virasoro scalar
lane one has $F_g(\mathrm{Vir}_c) = \kappaChHodge(\mathrm{Vir}_c) \cdot \lambda_g^{\mathrm{FP}}$
\textup{(}uniform-weight\textup{)} for $g \geq
1$, so each scalar coefficient grows linearly with~$\kappaChHodge(\mathrm{Vir}_c)$ while the
ratios $F_g/F_0 = \lambda_g^{\mathrm{FP}}$ are independent of
$\kappaChHodge(\mathrm{Vir}_c)$. If one specializes the formal genus parameter to
$\hbar = 1/c \sim G_N/\ell$, then the reduced scalar series
$\sum_{g \geq 1} F_g \hbar^g$ is an honest convergent expansion in
$1/c$, while the full scalar sewing series still carries the
tree-level term $F_0=\kappaChHodge(\mathrm{Vir}_c)=c/2$.
\end{remark}

\subsubsection{BTZ-normalized scalar shadow series}

\begin{theorem}[BTZ-normalized scalar shadow series; \ClaimStatusProvedHere]
% label removed: thm:thqg-I-btz
\index{BTZ black hole!shadow free energy}
Under the Brown--Henneaux identification, the genus-$1$ scalar shadow
term is
\begin{equation}% label removed: eq:thqg-I-btz-F1
F_1^{\mathrm{BTZ}} \;=\; F_1(\mathrm{Vir}_c) \;=\; \frac{c}{48} \;=\; \frac{\ell}{32G_N}\,.
\end{equation}
On the proved scalar lane \textup{(}here $\mathrm{Vir}_c$\textup{)}, the
associated scalar shadow series is
\begin{equation}% label removed: eq:thqg-I-btz-full
Z_{\mathrm{BTZ}}^{\mathrm{pert}}(\hbar) \;=\; \frac{c}{2} \cdot \frac{\sqrt{\hbar}/2}{\sin(\sqrt{\hbar}/2)} \;=\; \frac{3\ell}{4G_N} \cdot \frac{\sqrt{\hbar}/2}{\sin(\sqrt{\hbar}/2)}\,.
\end{equation}
No comparison theorem identifying this scalar series with the full
analytic BTZ partition function is proved here.
\end{theorem}

\begin{proof}
By the Brown--Henneaux identification
(Setup~\ref{setup:thqg-I-brown-henneaux}), the boundary algebra is
$\mathrm{Vir}_c$ with $\kappaChHodge(\mathrm{Vir}_c) = c/2$. The genus-$1$ scalar value is
$F_1(\mathrm{Vir}_c) = \kappaChHodge(\mathrm{Vir}_c)/24 = c/48$. Since $\mathrm{Vir}_c$ lies on the proved
scalar lane, the scalar generating function is
\eqref{V1-eq:thqg-I-gf-ahat} with $\kappaChHodge(\mathrm{Vir}_c) = c/2$. Substituting
$c = 3\ell/(2G_N)$ gives $\kappaChHodge(\mathrm{Vir}_c) = 3\ell/(4G_N)$. The argument stops at the
scalar shadow series; no comparison map to the full analytic BTZ
partition function is constructed here.
\end{proof}

\begin{remark}[Comparison with known results]
% label removed: rem:thqg-I-comparison-btz
\index{BTZ black hole!comparison}
The genus-$1$ correction $F_1 = c/48$ has been computed by multiple
methods: (a)~metric one-loop determinant
(Giombi--Maloney--Yin~\cite{GMY08}):
the graviton determinant on the solid torus gives
$\prod_{n=2}^{\infty} (1 - q^n)^{-1}$; (b)~Chern--Simons theory
(Witten~\cite{Witten88}): 3d gravity as
$\mathrm{SL}(2,\mathbb{R}) \times \mathrm{SL}(2,\mathbb{R})$ CS theory,
with $c/24$ for the full theory and $c/48$ for the holomorphic half;
(c)~modular bootstrap (Maloney--Witten~\cite{MW10}). These comparisons
support the genus-$1$ normalization used above. Beyond genus~$1$, the
shadow method provides the scalar series shown above on the Virasoro
scalar lane, but its identification with the full analytic BTZ
partition function is additional dictionary data not proved here.
\end{remark}

\subsubsection{Higher-genus corrections}

\begin{proposition}[Higher-genus BTZ-normalized scalar coefficients; \ClaimStatusProvedHere]
\label{prop:thqg-I-btz-higher-genus}
\index{BTZ black hole!higher-genus corrections}
The genus-$g$ scalar shadow coefficient in the BTZ normalization is
\begin{equation}% label removed: eq:thqg-I-btz-Fg
F_g^{\mathrm{BTZ}} \;=\; \frac{c}{2} \cdot \frac{2^{2g-1}-1}{2^{2g-1}} \cdot \frac{|B_{2g}|}{(2g)!} \;=\; \frac{3\ell}{4G_N} \cdot \lambda_g^{\mathrm{FP}}\,.
\end{equation}
The first five corrections are:
\begin{align}
F_1^{\mathrm{BTZ}} &= \frac{c}{48}\,, % label removed: eq:thqg-I-btz-F1-val \\
F_2^{\mathrm{BTZ}} &= \frac{7c}{11520}\,, % label removed: eq:thqg-I-btz-F2-val \\
F_3^{\mathrm{BTZ}} &= \frac{31c}{1935360}\,, % label removed: eq:thqg-I-btz-F3-val \\
F_4^{\mathrm{BTZ}} &= \frac{127c}{309657600}\,, % label removed: eq:thqg-I-btz-F4-val \\
F_5^{\mathrm{BTZ}} &= \frac{73c}{7007109120}\,. % label removed: eq:thqg-I-btz-F5-val
\end{align}
The ratio $F_{g+1}/F_g \to 1/(4\pi^2)$ as $g \to \infty$, consistent with the first scalar pole at $\hbar = 4\pi^2$.
\end{proposition}

\begin{proof}
Direct from $\kappaChHodge(\mathrm{Vir}_c) = c/2$ and the universal formula (Proposition~\ref{prop:thqg-I-Fg-closed-form}).
\end{proof}

\begin{remark}[Numerical values for $c = 24$ (monster CFT)]
% label removed: rem:thqg-I-monster
\index{monster CFT}
For the monster CFT with $c = 24$ (the Frenkel--Lepowsky--Meurman moonshine module $V^\natural$), $\kappaChHodge(V^\natural) = 12$, since $\dim V_1^\natural = 0$ forces the modular characteristic to be determined by the Virasoro sector alone:
\begin{align*}
F_1 &= 12/24 = 1/2, \\
F_2 &= 7 \cdot 12/5760 = 84/5760 = 7/480, \\
F_3 &= 31 \cdot 12/967680 = 372/967680 = 31/80640, \\
F_4 &= 127 \cdot 12/154828800 = 1524/154828800 = 127/12902400.
\end{align*}
The one-loop contribution $F_1 = 1/2$ is order $1$, and each
subsequent genus is suppressed by the Bernoulli factor. The scalar
sewing series at $\hbar = 1$ is
\[
Z_{\mathrm{sh}}^{\mathrm{scal}}(V^\natural;\,1)
= 12 \cdot \frac{1/2}{\sin(1/2)}
\approx 12.515.
\]
\end{remark}

\begin{computation}[BTZ-normalized scalar shadow through five loops]
\label{comp:btz-five-loop}%
\index{BTZ black hole!five-loop scalar shadow|textbf}%
The BTZ-normalized scalar shadow expansion through genus~$5$, on the
proved scalar lane
$(\mathrm{Vir}_c$, $\kappaChHodge(\mathrm{Vir}_c)=c/2)$:
\begin{equation}\label{eq:btz-five-loop}
S_{\mathrm{BTZ}}^{\mathrm{scal}}
\;=\;
S_{\mathrm{BH}}
\;+\;
\sum_{g=1}^{5} F_g^{\mathrm{BTZ}} \cdot
\kappaChHodge(\mathrm{Vir}_c)^{2g-2},
\end{equation}
with $F_g^{\mathrm{BTZ}}$ as in
Proposition~\textup{\ref{prop:thqg-I-btz-higher-genus}}.
The explicit values are $F_1 = c/48$,
$F_2 = 7c/11520$,
$F_3 = 31c/1935360$,
$F_4 = 127c/309657600$,
$F_5 = 73c/7007109120$.
Series convergence is verified by the ratio test:
the consecutive ratios
\begin{align*}
F_2/F_1 &= 7/240, &
F_3/F_2 &= 31/1176, \\
F_4/F_3 &= 127/4960, &
F_5/F_4 &= 73 \cdot 309657600\,/\,(127 \cdot 7007109120)
\;\approx\; 0.0254
\end{align*}
converge to $1/(4\pi^2) \approx 0.0253$, the reciprocal of
the convergence radius $\hbar_{\mathrm{crit}} = 4\pi^2$.
At genus~$5$ the ratio is $|F_5/F_4| = 2555/100584
\approx 0.0254$, well within the convergence disk.
The scalar shadow series
$Z_{\mathrm{BTZ}}^{\mathrm{pert}}(\hbar)
= \kappaChHodge(\mathrm{Vir}_c) \cdot
\frac{\sqrt{\hbar}/2}{\sin(\sqrt{\hbar}/2)}$
has convergence radius $|\hbar| < 4\pi^2$, and the five-loop
truncation captures $>99.5\%$ of the exact scalar sum at
$\hbar = 1$.
This is not a derivation of the physical BTZ entropy; the comparison
with $S_{\mathrm{BH}}$ requires the Cardy/vacuum-dominance dictionary
and the saddle analysis fenced in Remark~\ref{rem:thqg-I-comparison-btz}.
\end{computation}

\subsubsection{The quartic contact invariant in sub-subleading corrections}

\begin{remark}[Quartic contact invariant in 3d gravity]
% label removed: rem:thqg-I-quartic-3d
\index{quartic contact invariant!3d gravity}
The scalar channel captures only the degree-$0$ information. The first non-scalar correction comes from the quartic contact invariant
\begin{equation}% label removed: eq:thqg-I-quartic-contact
Q^{\mathrm{contact}}_{\mathrm{Vir}} \;=\; \frac{10}{c(5c+22)}\,,
\end{equation}
which is the degree-$4$, genus-$0$ component of the support packet. For $c = 3\ell/(2G_N)$, this gives $Q^{\mathrm{contact}} = 10/(c(5c+22)) = 2/c^2 + O(1/c^3)$, a shadow channel whose Brown--Henneaux scaling begins at order $G_N^2/\ell^2$. Interpreting it as a two-loop graviton correction requires a separate metric comparison datum. The genus-$1$ Hessian correction is
\begin{equation}% label removed: eq:thqg-I-genus1-hessian
\delta_{H,\mathrm{Vir}}^{(1)} \;=\; \frac{120}{c^2(5c+22)}\,x^2\,.
\end{equation}
These corrections are specific to the mixed support-shadow class
\(r^\Theta_{\max}=\infty\) of the Virasoro algebra.
\end{remark}

\begin{remark}[Support depth and gravitational comparison complexity]
% label removed: rem:thqg-I-depth-complexity
\index{support shadow depth!gravitational comparison complexity}
The support shadow depth classification ($G/L/C/M$ for Gaussian/Lie/Contact/Mixed) stratifies the bar-intrinsic support packet; after a physical comparison datum is supplied, it can stratify nonlinear gravitational comparison channels:
\begin{enumerate}[label=\textup{(\alph*)}]
\item \emph{Gaussian (\(r^\Theta_{\max}=2\)):} Heisenberg, lattice VOAs, free fermions. Only the quadratic/scalar channel survives. On the proved analytic locus, the scalar shadow series is represented by a Fredholm determinant.
\item \emph{Lie-type (\(r^\Theta_{\max}=3\)):} Affine Kac--Moody algebras. The degree-summed shadow trace receives trivalent-graph corrections.
\item \emph{Contact-type (\(r^\Theta_{\max}=4\)):} $\beta\gamma$ system. Quartic interactions but no higher support components on the canonical contact branch.
\item \emph{Mixed (\(r^\Theta_{\max}=\infty\)):} Virasoro, $\mathcal{W}_N$. The full support packet is infinite, and the perturbative finiteness theorem is needed in full strength. The quintic obstruction \(o^{(5)}_{\mathrm{Vir}}\neq0\) is the first obstruction witness to the genuinely infinite tower.
\end{enumerate}
The Brown--Henneaux boundary theory for pure 3d gravity is $\mathrm{Vir}_c$, which falls in the mixed class.
\end{remark}

\subsubsection{Large-$c$ expansion and gravitational coupling}

\begin{proposition}[Leading Brown--Henneaux expansion of the reduced scalar shadow series; \ClaimStatusProvedHere]
% label removed: prop:thqg-I-1c-expansion
\index{Brown--Henneaux!scalar shadow large-$c$ expansion}
In the semiclassical specialization $\hbar = 1/c$, the reduced scalar
series \textup{(}obtained by subtracting the genus-$0$ term
$F_0=c/2$\textup{)} has the convergent expansion
\begin{equation}% label removed: eq:thqg-I-1c-expansion
Z_{\mathrm{sh}}^{\mathrm{scal}}(\mathrm{Vir}_c;\,1/c) - \frac{c}{2}
\;=\;
\frac{1}{48} + \frac{7}{11520\,c} + \frac{31}{1935360\,c^2}
\;+\; \frac{127}{309657600\,c^3} + O(c^{-4})\,.
\end{equation}
The leading term $1/48$ is the genus-$1$ scalar shadow coefficient of
the reduced series, and each subsequent term is suppressed by one
additional Brown--Henneaux power $G_N/\ell \sim 1/c$.
\end{proposition}

\begin{proof}
From $F_g(\mathrm{Vir}_c) = (c/2)\lambda_g^{\mathrm{FP}}$, the
genus-$g$ contribution to the reduced scalar series at $\hbar = 1/c$ is
\[
F_g \hbar^g = \frac{c}{2}\lambda_g^{\mathrm{FP}}\,c^{-g}
= \frac{\lambda_g^{\mathrm{FP}}}{2c^{g-1}}.
\]
Therefore
\[
Z_{\mathrm{sh}}^{\mathrm{scal}}(\mathrm{Vir}_c;1/c) - \frac{c}{2}
= \sum_{g \geq 1} \frac{\lambda_g^{\mathrm{FP}}}{2c^{g-1}}.
\]
At $g = 1$, $\lambda_1^{\mathrm{FP}} = 1/24$, giving $1/48$; at
$g = 2$, one obtains $7/(11520c)$; and the remaining coefficients are
computed similarly.
\end{proof}

\begin{proposition}[Brown--Henneaux normalization of the scalar shadow closed form; \ClaimStatusProvedHere]
% label removed: prop:thqg-I-newton-expansion
\index{Brown--Henneaux!scalar shadow Newton expansion}
Substituting $c = 3\ell/(2G_N)$ and $\hbar = 2G_N/(3\ell)$ in the scalar
shadow closed form gives
\begin{equation}% label removed: eq:thqg-I-newton-expansion
Z_{\mathrm{sh}}^{\mathrm{scal}}(\mathrm{Vir}_c;\,2G_N/(3\ell)) \;=\; \frac{3\ell}{4G_N} \cdot \frac{\sqrt{2G_N/(3\ell)}/2}{\sin(\sqrt{2G_N/(3\ell)}/2)}\,.
\end{equation}
For $G_N/\ell \ll 1$, the Brown--Henneaux-normalized scalar shadow
expansion is
\begin{equation}% label removed: eq:thqg-I-G-expansion
Z_{\mathrm{sh}}^{\mathrm{scal}}
\;=\; \frac{3\ell}{4G_N}\left(1 + \frac{G_N}{36\ell} + \frac{7G_N^2}{12960\ell^2} + O(G_N^3/\ell^3)\right).
\end{equation}
\end{proposition}

\begin{proof}
Direct Taylor expansion of $x/(2\sin(x/2))$ at $x = 0$ with $x^2 = 2G_N/(3\ell)$.
\end{proof}

\subsubsection{Towards the non-perturbative regime}

\begin{conjecture}[Non-perturbative comparison datum; \ClaimStatusConjectured]
\label{conj:thqg-I-non-perturbative}
\index{non-perturbative!scalar shadow comparison}
The meromorphic scalar shadow closed form
\[
Z_{\mathrm{sh}}^{\mathrm{scal}}(\cA;\,\hbar)
=\kappaChHodge(\cA)\cdot
\frac{\sqrt{\hbar}/2}{\sin(\sqrt{\hbar}/2)}
\]
is expected to be the scalar projection of a non-perturbative physical
comparison datum only after a bulk completion, contour, measure,
unitarity/positivity structure, and saddle sector have been supplied.
Such a datum should satisfy:
\begin{enumerate}[label=\textup{(\roman*)}]
\item Pole-controlled scales: any completion compatible with the proved scalar lane reflects the distinguished action scales $4\pi^2 n^2$ coming from the pole set of the meromorphic continuation.
\item Exponential sectors: if genuinely non-perturbative contributions are present in the physical completion, their leading scales are expected to be compatible with $e^{-4\pi^2 n^2/\hbar}$; identifying them with non-handlebody saddle points is heuristic.
\item Positivity datum: positivity for real $\hbar > 0$ belongs to the completed physical partition function after the unitarity/measure datum is fixed; it is not a consequence of the scalar shadow closed form alone.
\end{enumerate}
\end{conjecture}

\begin{evidence}[Evidence for the resurgent structure]
% label removed: evid:thqg-I-resurgent
The Taylor coefficients $F_g \sim 2\kappaChHodge(\cA)/(2\pi)^{2g}$ show that the scalar genus series already converges absolutely for $|\hbar| < 4\pi^2$ and continues meromorphically beyond that disk. The Borel transform is
\begin{equation}% label removed: eq:thqg-I-borel-transform
\hat{Z}_{\mathrm{sh}}(\xi) \;=\; \sum_{g=1}^\infty \frac{F_g}{\Gamma(g)} \xi^{g-1} \;=\; \kappaChHodge(\cA) \sum_{g=1}^\infty \frac{\lambda_g^{\mathrm{FP}}}{(g-1)!} \xi^{g-1}.
\end{equation}
Since $\lambda_g^{\mathrm{FP}} \sim 2/(2\pi)^{2g}$, the ratio $\lambda_g^{\mathrm{FP}}/\Gamma(g) \sim 2/((2\pi)^{2g}(g-1)!)$, and the series converges for all $\xi \in \mathbb{C}$: the Borel transform is entire, so the usual asymptotic-resurgence argument does not apply here. What is visible in the proved scalar lane is instead the meromorphic pole set $\hbar = 4\pi^2 n^2$ of the original function. These pole scales are compatible with, but do not prove, heuristic exponential sectors of order $e^{-4\pi^2 n^2/\hbar}$ in a fuller non-perturbative completion.
\end{evidence}

\begin{remark}[BTZ black holes as non-perturbative saddles]
\label{rem:thqg-I-btz-saddles}
\index{BTZ black hole!non-perturbative}
In 3d gravity, the BTZ black holes of mass $M = n^2/(8G_N\ell)$ ($n \geq 1$) are natural candidates for non-perturbative saddle points. Their Euclidean action is $I_n = 4\pi^2 n^2 / \hbar$ (with $\hbar = 2G_N/(3\ell)$). The coincidence between these action scales and the pole scales of the scalar meromorphic continuation suggests, but does not prove, that BTZ sectors contribute terms of the form $c_n e^{-I_n/\hbar}$. No Stokes-ray or Borel-plane identification is established by the convergent scalar genus series itself.
\end{remark}

%% ===================================== Stokes phenomena on the Steinberg object
\providecommand{\SteinbergB}{\mathfrak{S}_b}

\subsubsection{Stokes phenomena on the Steinberg object}
% label removed: sssec:thqg-stokes-steinberg
\index{Stokes phenomena!Steinberg object}
\index{Steinberg object!Stokes filtration}

The pole scales $\hbar = 4\pi^2 n^2$ of
$Z_{\mathrm{sh}}^{\mathrm{scal}}(\cA;\,\hbar)
= \kappaChHodge(\cA)\cdot
\frac{\sqrt{\hbar}/2}{\sin(\sqrt{\hbar}/2)}$
define a possible Stokes-filtered organization on the
Steinberg object~$\SteinbergB$. On the proved scalar lane, however,
the Borel transform
\[
B[Z_{\mathrm{sh}}^{\mathrm{scal}}](\zeta) \;=\; \sum_{g \geq 1} \frac{F_g}{\Gamma(g)}\,\zeta^{g-1}
\]
is entire, so these scales are not singularities of the scalar Borel
transform itself. Rather, they are candidate action scales for a
fuller non-perturbative completion, heuristically associated with the
$n$-th BTZ saddle points (Remark~\ref{rem:thqg-I-btz-saddles}).
The Stokes-filtered Steinberg construction below is a conjectural
organization of possible exponential sectors,
not a theorem extracted from the proved scalar lane.

\begin{construction}[The Stokes-filtered Steinberg]
% label removed: constr:thqg-stokes-filtered-steinberg
\index{Steinberg object!Stokes-filtered}
Define the \emph{Stokes-filtered Steinberg} $\SteinbergB^{\mathrm{Stokes}}$ as follows. For each $n \geq 1$, let $\SteinbergB^{\mathrm{BTZ},n}$ denote the Steinberg object associated to the $n$-th BTZ saddle, i.e.\ the analogue of~$\SteinbergB$ constructed from the gauge equivalence data of the $n$-th non-handlebody saddle point with Euclidean action $I_n = 4\pi^2 n^2 / \hbar$. The perturbative Steinberg~$\SteinbergB$ is $\SteinbergB^{\mathrm{BTZ},0}$ (the vacuum saddle). The Stokes-filtered Steinberg is the datum
\[
\SteinbergB^{\mathrm{Stokes}} \;=\; \bigl(\SteinbergB,\;\{\SteinbergB^{\mathrm{BTZ},n}\}_{n \geq 1},\;\{S_n\}_{n \geq 1}\bigr)
\]
where $S_n \colon H_*^{\mathrm{BM}}(\SteinbergB) \to H_*^{\mathrm{BM}}(\SteinbergB^{\mathrm{BTZ},n})$ is a conjectural Stokes map in a fuller non-perturbative completion, indexed by the action scale $\zeta_n = 4\pi^2 n^2$.
\end{construction}

\begin{conjecture}[Resurgent structure from Stokes-filtered BM homology; \ClaimStatusConjectured]
% label removed: prop:thqg-stokes-bm-resurgence
\index{resurgence!Steinberg BM homology}
\index{BM homology!Stokes-filtered}
A conjectural physical completion whose scalar projection is
$Z_{\mathrm{sh}}^{\mathrm{scal}}$, if it admits exponential sectors at
the pole scales $\zeta_n = 4\pi^2 n^2$, is organized by the
Stokes-filtered Borel--Moore homology of the Steinberg:
\begin{equation}% label removed: eq:thqg-stokes-bm-decomposition
H_*^{\mathrm{BM},\,\mathrm{Stokes}}(\SteinbergB) \;=\; H_*^{\mathrm{BM}}(\SteinbergB) \;\oplus\; \bigoplus_{n \geq 1} H_*^{\mathrm{BM}}(\SteinbergB^{\mathrm{BTZ},n})\,,
\end{equation}
where $\SteinbergB^{\mathrm{BTZ},n}$ is the Steinberg object for the $n$-th BTZ saddle. The Stokes automorphism
\[
\mathfrak{S}_n \colon H_*^{\mathrm{BM}}(\SteinbergB) \;\longrightarrow\; H_*^{\mathrm{BM}}(\SteinbergB)
\]
would exchange the perturbative and $n$-th BTZ sectors. In such a
completion one expects exponential contributions of the form
\[
\mathfrak{S}_n([\SteinbergB]) \cdot e^{-\zeta_n / \hbar},
\]
where $[\SteinbergB] \in H_*^{\mathrm{BM}}(\SteinbergB)$ is the
fundamental class. The proved scalar lane does not identify actual
Borel-plane singularities or lateral-resummation discontinuities.
\end{conjecture}

\begin{remark}[Positivity requires physical comparison data]
% label removed: rem:thqg-stokes-positivity
\index{positivity!scalar shadow comparison}
\index{BM homology!fundamental class positivity}
Positivity of a physical completion for real $\hbar > 0$ (away from
the scalar pole set $\hbar = 4\pi^2 n^2$) is not a property of the
meromorphic scalar shadow alone. It would require a comparison with a
positive geometric pairing, for example through the Steinberg
presentation. The perturbative contribution would be
$\langle [\SteinbergB], [\SteinbergB] \rangle_{\mathrm{BM}}$, the
self-pairing of the BM fundamental class $[\SteinbergB]$, which is
positive since $\SteinbergB$ carries a canonical orientation as a
complex-algebraic variety. In a fuller non-perturbative completion, the
pole scales $4\pi^2 n^2$ could mark where BTZ strata
$\SteinbergB^{\mathrm{BTZ},n}$ become relevant. Current proofs do not
identify these scales with literal Borel-plane Stokes jumps.
\end{remark}

\begin{remark}[Pure 3d gravity as a model]
% label removed: rem:thqg-I-pure-3d
\index{3d gravity!pure}
Pure 3d gravity with $\Lambda < 0$ is the simplest non-trivial gravitational theory. Its boundary algebra is $\mathrm{Vir}_c$ with $c = 3\ell/(2G_N)$, which falls in the mixed class. The perturbative finiteness theorem provides a complete, explicit, genus-by-genus calculation of the scalar-shadow partition series:
\[
Z_{\mathrm{sh}}^{\mathrm{scal}}(\mathrm{Vir}_c;\,\hbar)
\;=\; \frac{3\ell}{4G_N} \cdot
\frac{\sqrt{\hbar}/2}{\sin(\sqrt{\hbar}/2)}\,.
\]
This is the scalar-shadow partition series determined by
\(\kappaChHodge(\mathrm{Vir}_c)=c/2\), not the full gravitational
partition function over bulk geometries.
The non-scalar corrections from $Q^{\mathrm{contact}} = 10/(c(5c+22))$ and the infinite support packet encode higher comparison channels, suppressed by $1/c^2$ relative to the scalar contribution before any metric interpretation is installed.
\end{remark}

\begin{remark}[Comparison with matrix models]
% label removed: rem:thqg-I-matrix-models
\index{matrix model!scalar shadow comparison}
In 2d topological gravity (Witten--Kontsevich), the genus-$g$ free energy involves Bernoulli numbers with a different prefactor: the Witten--Kontsevich generating function is $\sum_g F_g^{\mathrm{WK}} (2g-2)! x^{2g-2}$, with $F_g^{\mathrm{WK}} = \chi(\mathcal{M}_g) / (2g-2)!$ involving the orbifold Euler characteristic. The relation to the shadow generating function is $\frac{x/2}{\sin(x/2)} = \left.\frac{x/2}{\sinh(x/2)}\right|_{x \mapsto ix}$, reflecting the Wick rotation from Lorentzian to Euclidean signature. In the matrix model, the $\sinh$ function arises from the Gaussian integral, while in the shadow theory, the $\sin$ function arises from the Hodge integral; the two are related by analytic continuation $x \mapsto ix$.
\end{remark}

\begin{remark}[Higher-spin gravity]
% label removed: rem:thqg-I-higher-spin
\index{higher-spin gravity}
For higher-spin gravity in AdS$_3$ with $\mathcal{W}_N$ symmetry, the boundary algebra is $\mathcal{W}_k(\mathfrak{sl}_N)$ at central charge $c = (N-1) - N(N^2-1)(k+N-1)^2/(k+N)$. The modular characteristic is $\kappaChHodge(\mathcal{W}_N) = c \cdot (H_N - 1)$ where $H_N = \sum_{j=1}^N 1/j$ is the $N$-th harmonic number ($\kappaChHodge=c/2$ for $N=2$, $\kappaChHodge=5c/6$ for $N=3$). The scalar shadow closed form is
\[
Z_{\mathrm{sh}}^{W_N,\mathrm{scal}}(\hbar) \;=\; \kappaChHodge(\mathcal{W}_N) \cdot \frac{\sqrt{\hbar}/2}{\sin(\sqrt{\hbar}/2)}\,,
\]
which has the same scalar functional form as the Virasoro case but
with a different normalization. The scalar convergence radius and pole
set are identical. The degree-summed non-scalar correction tower is different:
$\mathcal{W}_N$ has support shadow depth \(r^\Theta_{\max}=\infty\)
and the shadow tower includes contributions from all $W^{(s)}$ fields
($s = 2, \ldots, N$).
\end{remark}

\subsubsection{The holographic dictionary}

\begin{remark}[The holographic dictionary]
% label removed: rem:thqg-I-holographic-dictionary
\index{holographic dictionary!perturbative finiteness}
The algebraic shadow finiteness theorem
(Theorem~\ref{thm:thqg-I-perturbative-finiteness}) enters the
holographic dictionary through the scalar-shadow lane:
\begin{center}
\renewcommand{\arraystretch}{1.3}
\begin{tabular}{l|l}
\textbf{Boundary algebra} & \textbf{Comparison channel} \\
\hline
$\cA$ (chiral algebra) & boundary conditions \\
$\kappaChHodge(\cA)$ (modular characteristic) & cosmological constant \\
$F_g(\cA)$ (shadow free energy) & genus-$g$ scalar vacuum amplitude \\
HS-sewing condition & scalar sewing finiteness \\
$\Theta_\cA$ (MC element) & classical solution + perturbations \\
$D_\cA^2 = 0$ & Ward identity \\
\(r^\Theta_{\max}(\cA)\) (support shadow depth) & gravitational complexity \\
$4\pi^2$ (scalar convergence radius) & first scalar pole scale \\
$Q^{\mathrm{contact}}$ (quartic invariant) & quartic/contact comparison channel \\
$\cA^!$ (Koszul dual) & dual boundary conditions \\
\end{tabular}
\end{center}
The passage from boundary data to a physical bulk partition function
requires a comparison datum. Algebraically, given the OPE data of
$\cA$, the modular characteristic $\kappaChHodge(\cA)$ is computable,
and the genuswise shadow amplitudes are finite. On the proved scalar
lane, this is supplemented by the closed scalar shadow series
$Z_{\mathrm{sh}}^{\mathrm{scal}}
= \kappaChHodge(\cA)\cdot x/(2\sin(x/2))$. Beyond that lane the
unconditional scalar data are fixed-genus finiteness and the universal
genus-$1$ term. The algebraic shadow complex requires no bar-complex
counterterm at fixed genus and degree; BV renormalisation and anomaly
matching belong to the separate physical realisation.
\end{remark}

\begin{remark}[Summary of the finiteness mechanism]
% label removed: rem:thqg-I-finiteness-summary
\index{perturbative finiteness!mechanism summary}
The perturbative finiteness of twisted quantum gravity rests on a chain of three implications:
\begin{equation}% label removed: eq:thqg-I-chain
\underbrace{D_\cA^2 = 0}_{\text{bar complex}} \;\Longrightarrow\; \underbrace{\Theta_\cA \in \mathrm{MC}(\gAmod)}_{\text{all-genera MC}} \;\Longrightarrow\; \underbrace{\int_{\overline{\mathcal{M}}_g}\operatorname{Sh}_{g,0}(\Theta_\cA) < \infty}_{\text{genuswise finite amplitudes}}.
\end{equation}
The first implication is the bar-intrinsic construction
\textup{(}Theorem~\ref*{V1-thm:mc2-bar-intrinsic}\textup{)}: the
element $\Theta_\cA := D_\cA - d_0$ is automatically Maurer--Cartan
because the bar differential squares to zero. The second implication
is the projection onto genus~$g$: the shadow free energy is the
integral of a cohomology class over a compact moduli space, hence
finite. On the proved scalar lane this genuswise finiteness is
supplemented by the closed formula
$F_g^{\mathrm{sc}}(\cA)=\kappaChHodge(\cA)\lambda_g^{\mathrm{FP}}$
\textup{(}uniform-weight\textup{)}; without the scalar-lane
hypothesis the universal scalar identity is
$F_1^{\mathrm{sc}}(\cA)=\kappaChHodge(\cA)/24$. The
mechanism is algebraic: the fixed-degree bar complex has no
counterterm and no asymptotic truncation.
\end{remark}


% ======================================================================
% SUBSECTION 7: THE HS-SEWING CONDITION AND THE MC RECURSION
% ======================================================================

\subsection{HS-sewing and the MC recursion}
% label removed: subsec:thqg-I-hs-mc-recursion
\index{HS-sewing!MC recursion}

The HS-sewing condition and the MC recursion interact at two levels: the HS condition provides the \emph{analytic} input (convergence of sewing amplitudes), while the MC recursion provides the \emph{algebraic} input (unique determination of higher-genus data from lower-genus data). Together, they give a genus-by-genus algorithm for computing the gravitational partition function: the analytic half (HS-sewing convergence) is proved; the algebraic comparison between the sewing output and the bar-intrinsic MC element (the BV/BRST identification) remains conjectural at each genus.

\subsubsection{The sewing amplitude as a graph sum}

\begin{proposition}[Graph-sum representation of sewing amplitudes; \ClaimStatusProvedHere]
% label removed: prop:thqg-I-graph-sum-sewing
\index{sewing amplitude!graph-sum representation}
Let $\cA$ be a modular Koszul chiral algebra satisfying HS-sewing. The genus-$g$ sewing amplitude on a surface $\Sigma_g$ with pants decomposition $\mathcal{P}$ has the graph-sum representation
\begin{equation}% label removed: eq:thqg-I-graph-sewing
Z_{\mathcal{P}}(\cA;\,\mathbf{q}) \;=\; \sum_{\Gamma \in \Gmod(g)} \frac{1}{|\Aut(\Gamma)|}\,\ell_\Gamma(\mathbf{q}),
\end{equation}
where the sum runs over connected modular graphs $\Gamma$ of total genus $g$, and $\ell_\Gamma(\mathbf{q})$ is the graph amplitude: the product of vertex labels (transferred cyclic operations $\ell_n^{\mathrm{tr}}$) and edge propagators ($P_\cA(q_e) = \sum_{n \geq 0} q_e^n \pi_n$, the weighted projection onto the $n$-th conformal weight).
\end{proposition}

\begin{proof}
The pants decomposition $\mathcal{P}$ determines a dual graph $\Gamma_\mathcal{P}$ with one vertex for each pair of pants and one edge for each sewing circle. The vertex labels are the three-point functions of $\cA$ (the OPE coefficients $C^c_{a,b}$), and the edge labels are the sewing propagators. The sewing amplitude is the contraction of all vertex and edge tensors:
\[
Z_{\mathcal{P}} \;=\; \sum_{\text{weight assignments}} \prod_{\text{vertices}} C^{c_v}_{a_v,b_v} \prod_{\text{edges}} q_e^{n_e} \delta_{n_e = c_{v'} = a_{v''}},
\]
where the sum over weight assignments runs over all ways to assign conformal weights to the internal edges. This is the graph amplitude $\ell_{\Gamma_\mathcal{P}}(\mathbf{q})$.

The independence from the pants decomposition (sewing consistency) means that $Z$ depends only on the genus $g$ and the moduli of $\Sigma_g$, not on the choice of $\mathcal{P}$. Different pants decompositions give different graphs $\Gamma_\mathcal{P}$, but the sum of all amplitudes is the same.
\end{proof}

\begin{proposition}[Fixed-degree graph convergence; \ClaimStatusProvedHere]
% label removed: prop:thqg-I-graph-sum-convergence
\index{graph sum!convergence}
Under HS-sewing, the degree-$r$ graph sum
\[
Z_{\mathcal{P},r}(\cA;\,\mathbf{q})
:=
\sum_{\Gamma \in \Gmod(g,r)}
\frac{1}{|\Aut(\Gamma)|}\,\ell_\Gamma(\mathbf{q})
\]
converges absolutely for each fixed pair $(g,r)$:
\begin{equation}% label removed: eq:thqg-I-graph-sum-degree
|Z_{\mathcal{P},r}(\cA;\,\mathbf{q})| \;\leq\; \sum_{\Gamma \in \Gmod(g,r)} \frac{|\ell_\Gamma(\mathbf{q})|}{|\Aut(\Gamma)|} \;<\; \infty.
\end{equation}
The full graph expression is the completed series
$\widehat{\sum}_{r \geq 0} Z_{\mathcal{P},r}$; absolute convergence
of $\sum_{r \geq 0} Z_{\mathcal{P},r}$ is an additional
shadow-degree summability estimate.
\end{proposition}

\begin{proof}
At each degree $r$, the number of graphs $|\Gmod(g,r)|$ is finite
(Proposition~\ref{prop:thqg-I-graph-count}). Each graph amplitude
$|\ell_\Gamma(\mathbf{q})|$ is bounded by the product of vertex norms
and propagator norms:
$|\ell_\Gamma| \leq \prod_{v} \|\ell_{n_v}^{\mathrm{tr}}\|
\cdot \prod_{e} q_e^{n_e}$. The HS-sewing condition bounds the sum
over internal weight assignments. Thus the degree-$r$ sum is finite.
The inverse-limit theorem
(Theorem~\ref{V1-thm:recursive-existence}) supplies the completed
degree tower. It does not supply a norm estimate for
$\sum_r |Z_{\mathcal{P},r}|$.
\end{proof}

\subsubsection{Interplay between analytic and algebraic finiteness}

\begin{remark}[Two finiteness mechanisms]
% label removed: rem:thqg-I-two-mechanisms
\index{finiteness!two mechanisms}
Two independent mechanisms enter the perturbative finiteness theorem:
\begin{enumerate}[label=\textup{(\roman*)}]
\item \emph{Algebraic finiteness} (from the MC equation). The bar-intrinsic construction $\Theta_\cA = D_\cA - d_0$ gives an MC element at all genera. The genus-$g$ component $\Theta^{(g)}$ is uniquely determined by the lower-genus data through the contracting homotopy. This determines the \emph{cohomology class} $[\operatorname{Sh}_{g,0}(\Theta_\cA)] \in H^*(\overline{\mathcal{M}}_g)$ without reference to convergence.
\item \emph{Analytic finiteness} (from HS-sewing). The HS condition guarantees that the sewing construction converges and that the genus-$g$ partition function, as a number and not only as a cohomology class, is finite.
\end{enumerate}
The algebraic mechanism gives existence and uniqueness of the scalar
shadow data as a formal quantity. The analytic mechanism gives
convergence and finiteness of the sewing integral as a number. The two
outputs remain distinct until a comparison map is constructed.
\end{remark}

\begin{remark}[Third finiteness: physical UV finiteness]
\label{rem:thqg-I-three-mechanisms}
\index{finiteness!three mechanisms}
The third mechanism is physical UV finiteness: a Costello--Gwilliam
renormalisation-group fixed point for the $3$d HT-QFT realisation of
\(\Phi_{\mathrm{hol}}(\cA,\omega,T_\omega,\Xi_\cA)\). Its input is a
chain-level $E_3$-topological factorisation algebra on the realized
bulk
\(\Obs^{\mathrm{bulk}}(\cF_{\cA,\Xi})\), together with the
bulk--Hochschild comparison
\(\chi_{\mathrm{HT},\cA}\colon
Z^{\mathrm{der}}_{\mathrm{ch}}(\cA)\to
\Obs^{\mathrm{bulk}}(\cF_{\cA,\Xi})\); the theorem above is only a
scalar shadow theorem. Costello--Li and Costello--Gaiotto supply the
required BV models in their stated gauge-theoretic lanes; class~M uses
the cohomological and weight-completed topologisation surfaces, and the
original-complex chain-level statement remains a separate hypothesis
outside those ambients.
\end{remark}

\begin{proposition}[Parallel analytic and algebraic finiteness; \ClaimStatusProvedHere]
% label removed: prop:thqg-I-consistency
\index{finiteness!consistency}
For a modular Koszul chiral algebra $\cA$ satisfying HS-sewing:
\begin{enumerate}[label=\textup{(\roman*)}]
\item For each fixed genus $g \geq 1$, the analytic sewing partition
function
\[
Z_g(\cA) = \int_{\overline{\mathcal{M}}_g}
Z_{\mathcal{P}}(\cA;\,\mathbf{q})\,\mu_{\mathrm{WP}}
\]
is finite.
\item Independently, the universal scalar algebraic statement is
$F_1^{\mathrm{sc}}(\cA)=\kappaChHodge(\cA)/24$ for every modular Koszul algebra; on the
scalar lane one further has
$F_g^{\mathrm{sc}}(\cA)=\kappaChHodge(\cA)\lambda_g^{\mathrm{FP}}$ for all $g \geq 1$.
\item On the scalar lane, the scalar genus series
\[
\sum_{g=1}^{\infty} F_g^{\mathrm{sc}}(\cA)\,\hbar^g
\]
converges absolutely for $|\hbar| < 4\pi^2$ and agrees there with
\[
\kappaChHodge(\cA)\cdot\left(\frac{\sqrt{\hbar}/2}{\sin(\sqrt{\hbar}/2)} - 1\right).
\]
Equivalently, after adjoining the genus-$0$ term
$F_0(\cA)=\kappaChHodge(\cA)$, the completed series agrees with
$\kappaChHodge(\cA)\cdot\frac{\sqrt{\hbar}/2}{\sin(\sqrt{\hbar}/2)}$.
\item The analytic sewing integrals in \textup{(i)} and the algebraic
scalar coefficients in \textup{(ii)} are distinct without a comparison
map.
\end{enumerate}
\end{proposition}

\begin{proof}
Part~(i) is Corollary~\ref{cor:thqg-I-genus-g-partition}.

Part~(ii) is Theorem~\ref{V1-thm:genus-universality}\textup{(iii)} on the
scalar lane together with the universal genus-$1$
statement of Theorem~\ref{V1-thm:genus1-universal-curvature}.

Part~(iii) is Theorem~\ref{thm:thqg-I-absolute-convergence}.

Part~(iv) follows by inspection of the preceding arguments: they
establish analytic finiteness and algebraic scalar control separately,
but do not construct a comparison map between them.
\end{proof}

\subsubsection{The degree-summed shadow partition series}

\begin{theorem}[Completed degree-summed shadow partition series and Gaussian scalar convergence; \ClaimStatusProvedHere]
\label{thm:thqg-I-full-convergence}
\index{shadow partition series!degree-summed convergence}
For a modular Koszul chiral algebra $\cA$ of support shadow depth
\(r^\Theta_{\max}(\cA)\)
satisfying HS-sewing, the completed degree-summed shadow partition
series satisfies:
\begin{enumerate}[label=\textup{(\roman*)}]
\item If \(r^\Theta_{\max}(\cA)<\infty\) (Gaussian, Lie, or contact
 class on the canonical standard-family branches), then
\begin{equation}% label removed: eq:thqg-I-full-finite-depth
Z_g^{\mathrm{sh,deg}}(T)
\;:=\;
\sum_{0\le r\le r^\Theta_{\max}(\cA)}
\int_{\overline{\mathcal{M}}_g}
\operatorname{Sh}_{g,r}(\Theta_\cA)
\end{equation}
is a finite sum in degree at each genus.
\item If \(r^\Theta_{\max}(\cA)=\infty\) (mixed class), then the degree tower
defines an element of the completed product
\[
\prod_{r \geq 0}
H^*(\overline{\mathcal{M}}_{g,n(g,r)})
\]
by Theorem~\ref{V1-thm:recursive-existence}. A numerical
degree-summed genus-$g$ shadow amplitude is obtained from it exactly
when the shadow trace
satisfies
\[
\sum_{r \geq 0}
\left|
\int_{\overline{\mathcal{M}}_{g,n(g,r)}}
\operatorname{Sh}_{g,r}(\Theta_\cA)
\right| < \infty .
\]
\item If \(r^\Theta_{\max}(\cA)=2\) (Gaussian class), then the degree-summed shadow
partition series equals the scalar shadow series:
\begin{equation}% label removed: eq:thqg-I-leading-term
Z_{\mathrm{sh}}^{\mathrm{deg}}(T;\,\hbar)
\;=\;
Z_{\mathrm{sh}}^{\mathrm{scal}}(T;\,\hbar),
\end{equation}
and therefore converges absolutely for $|\hbar| < 4\pi^2$ with
meromorphic continuation given by the closed form of
Theorem~\ref{thm:thqg-I-absolute-convergence}.
\end{enumerate}
\end{theorem}

\begin{proof}
Part~(i) is immediate from finite support shadow depth: at each genus \(g\),
only degrees \(r\le r^\Theta_{\max}(\cA)\) contribute, and each contribution is
finite.

Part~(ii) is the meaning of convergence in the completed weight
filtration: Theorem~\ref{V1-thm:recursive-existence} produces the
inverse limit of the degree truncations. Applying the genus-$g$ scalar
trace gives a number only under the displayed absolute summability
condition.

Part~(iii) is the Gaussian case: if \(r^\Theta_{\max}=2\), all higher-degree
corrections vanish, so the degree-summed shadow partition series
equals the scalar shadow series and inherits its convergence domain and
meromorphic continuation.
\end{proof}

\subsubsection{Modular characteristics for all standard families}

\begin{table}[ht]
\centering
\caption{Modular characteristics, support shadow depths, and support-depth classes for the standard landscape.}
% label removed: tab:thqg-I-modular-characteristics
\index{modular characteristic!standard landscape}
\renewcommand{\arraystretch}{1.4}
{\footnotesize
\begin{tabular}{|l|c|c|c|c|c|}
\hline
\textbf{Family $\cA$} & $\boldsymbol{\kappaChHodge(\cA)}$ & $\boldsymbol{r^\Theta_{\max}}$ & \textbf{Class} & $\boldsymbol{F_1(\cA)}$ & $\boldsymbol{F_2(\cA)}$ \\
\hline
$\mathcal{H}_k$ (rank $d$) & $dk$ & $2$ & G & $dk/24$ & $7dk/5760$ \\
\hline
$\widehat{\mathfrak{sl}}_{2,k}$ & $\frac{3(k+2)}{4}$ & $3$ & L & $\frac{k+2}{32}$ & $\frac{7(k+2)}{7680}$ \\
\hline
$\widehat{\mathfrak{sl}}_{N,k}$ & $\frac{(k+N)(N^2-1)}{2N}$ & $3$ & L & $\frac{(k+N)(N^2-1)}{48N}$ & $\frac{7(k+N)(N^2-1)}{11520N}$ \\
\hline
$\widehat{\mathfrak{so}}_{2N,k}$ & $\frac{(k+2N-2)N(2N-1)}{2(2N-2)}$ & $3$ & L &; &; \\
\hline
$\mathrm{Vir}_c$ & $c/2$ & $\infty$ & M & $c/48$ & $7c/11520$ \\
\hline
$\beta\gamma$ & $1$ & $4$ & C & $1/24$ & $7/5760$ \\
\hline
$\mathcal{W}_3(k)$ & $5c_3/6$ & $\infty$ & M &; &; \\
\hline
$\mathcal{W}_N(k)$ & $c_N(H_N - 1)$ & $\infty$ & M &; &; \\
\hline
$V_\Lambda$ & $\rank(\Lambda)$ & $2$ & G & $\rank(\Lambda)/24$ & $7\rank(\Lambda)/5760$ \\
\hline
$V_\Lambda^{N,q}$ (twisted) & $\rank(\Lambda)$ & $2$ & G & $\rank(\Lambda)/24$ & $7\rank(\Lambda)/5760$ \\
\hline
Free fermions (rank $n$) & $n/4$ & $2$ & G & $n/96$ & $7n/23040$ \\
\hline
$\widehat{\mathfrak{sl}}_{2,k}^{\otimes d}$ & $\frac{3d(k+2)}{4}$ & $3$ & L & $\frac{d(k+2)}{32}$ & $\frac{7d(k+2)}{7680}$ \\
\hline
\end{tabular}}
\end{table}

\begin{remark}[Lattice curvature-braiding orthogonality]
% label removed: rem:thqg-I-lattice-ortho
\index{lattice VOA!curvature-braiding orthogonality}
For the lattice VOA $V_\Lambda^{N,q}$ with cocycle twist $(N,q)$, the modular characteristic $\kappaChHodge(V_\Lambda^{N,q}) = \rank(\Lambda)$ is independent of the cocycle (Theorem~\ref{V1-thm:lattice:curvature-braiding-orthogonal}). The curvature $\kappaChHodge(V_\Lambda^{N,q})$ is a genus-$1$ invariant, while the cocycle twist acts only on the braiding (a genus-$0$ structure). The two are orthogonal in the modular cyclic deformation complex.
\end{remark}

\subsubsection{The HS-sewing condition as a stability criterion}

\begin{proposition}[HS-sewing and the strong filtration; \ClaimStatusProvedHere]
% label removed: prop:thqg-I-hs-filtration
\index{HS-sewing!strong filtration}
The HS-sewing condition~\eqref{V1-eq:thqg-I-hs-sewing-condition} is equivalent to the strong filtration axiom
\begin{equation}% label removed: eq:thqg-I-strong-filtration
m_q(F^i \otimes F^j) \;\subset\; F^{i+j}, \qquad F^n := \bigoplus_{k \geq n} H_k,
\end{equation}
holding at the level of Hilbert--Schmidt operators. The HS norm of the multiplication restricted to $F^i \otimes F^j$ decays as $O(q^{i+j})$, giving a submultiplicative bound.
\end{proposition}

\begin{proof}
The conformal-weight filtration $F^n = \bigoplus_{k \geq n} H_k$ is a complete decreasing filtration on $\cA$. The OPE preserves the conformal weight: $m^c_{a,b} = 0$ unless $c \leq a + b$ (by the conformal-weight bound on OPE). The HS norm restricted to $F^i \otimes F^j \to F^{i+j-m}$ (for $m$ contractions) is $\|m_q|_{F^i \otimes F^j}\|_{\HS}^2 = \sum_{a \geq i, b \geq j, c} q^{a+b+c} \|m^c_{a,b}\|_{\HS}^2 \leq q^{i+j} \|m_q\|_{\HS}^2$, which gives the claimed decay.
\end{proof}

\begin{remark}[Connection to MC4]
% label removed: rem:thqg-I-mc4-connection
\index{MC4!connection to HS-sewing}
The strong filtration axiom~\eqref{V1-eq:thqg-I-strong-filtration} is the analytic incarnation of the algebraic strong filtration that drives the completion-tower theorem (MC4, Theorem~\ref{V1-thm:completed-bar-cobar-strong}). In the algebraic setting, the strong filtration $\mu_r(F^{i_1}, \ldots, F^{i_r}) \subset F^{i_1 + \cdots + i_r}$ ensures that the bar-cobar construction passes to the inverse limit. In the analytic setting, the HS strong filtration ensures that the sewing construction converges. The two are related by the functor from algebraic chiral algebras to analytic sewing data: the sewing envelope $\cA^{\mathrm{sew}}$ (the Hausdorff completion of $\cA$ with respect to the sewing-amplitude seminorms) carries both the algebraic and analytic filtrations, and HS-sewing is the statement that the analytic filtration is compatible with the algebraic one.
\end{remark}

\subsubsection{Effective bounds for the genus expansion}

\begin{proposition}[Effective bound on the genus-$g$ scalar coefficient; \ClaimStatusProvedHere]
% label removed: prop:thqg-I-effective-bound
\index{shadow free energy!effective bound}
For any modular Koszul chiral algebra $\cA$ on the uniform-weight lane, with modular characteristic $\kappaChHodge(\cA)$ and
HS-sewing parameter $q_0 \in (0,1)$, the genus-$g$ shadow scalar
coefficient satisfies the effective bound
\begin{equation}% label removed: eq:thqg-I-effective-bound
|F_g(\cA)| \;\leq\; \frac{2|\kappaChHodge(\cA)|}{(2\pi)^{2g}} \cdot \left(1 + \frac{1}{2^{2g-1} - 1}\right) \;\leq\; \frac{4|\kappaChHodge(\cA)|}{(2\pi)^{2g}} \quad\text{for all } g \geq 1.
\end{equation}
The bound is tight: $|F_g(\cA)| \sim 2|\kappaChHodge(\cA)|/(2\pi)^{2g}$ as $g \to \infty$.
\end{proposition}

\begin{proof}
From $|F_g(\cA)| = |\kappaChHodge(\cA)| \cdot \frac{2^{2g-1}-1}{2^{2g-1}} \cdot \frac{|B_{2g}|}{(2g)!}$ and Euler's formula $|B_{2g}| = \frac{2(2g)!}{(2\pi)^{2g}} \cdot \zeta(2g)$:
\[
|F_g(\cA)| \;=\; |\kappaChHodge(\cA)| \cdot \frac{2^{2g-1}-1}{2^{2g-1}} \cdot \frac{2\zeta(2g)}{(2\pi)^{2g}} \;\leq\; |\kappaChHodge(\cA)| \cdot 1 \cdot \frac{2 \cdot \pi^2/6}{(2\pi)^{2g}} \;\leq\; \frac{4|\kappaChHodge(\cA)|}{(2\pi)^{2g}}
\]
for $g \geq 1$ (using $\zeta(2g) \leq \zeta(2) = \pi^2/6$ and $(2^{2g-1}-1)/2^{2g-1} \leq 1$).

The asymptotic tightness follows from $\zeta(2g) \to 1$ and $(2^{2g-1}-1)/2^{2g-1} \to 1$ as $g \to \infty$.
\end{proof}

\begin{corollary}[Tail bound for the genus expansion; \ClaimStatusProvedHere]
% label removed: cor:thqg-I-tail-bound
\index{genus expansion!tail bound}
The tail of the genus expansion satisfies
\begin{equation}% label removed: eq:thqg-I-tail
\left|\sum_{g=G+1}^{\infty} F_g(\cA)\,\hbar^g\right| \;\leq\; \frac{4|\kappaChHodge(\cA)|}{(2\pi)^{2G+2} - |\hbar|} \cdot |\hbar|^{G+1} \quad\text{for } |\hbar| < (2\pi)^2.
\end{equation}
In particular, at $\hbar = 1$ and $\kappaChHodge(\cA) = 1$, the tail beyond genus $G = 5$ is bounded by $4/(4\pi^2 - 1) \cdot (4\pi^2)^{-5} \approx 4 \times 10^{-9}$.
\end{corollary}

\begin{proof}
By the effective bound: $\sum_{g > G} |F_g(\cA)| |\hbar|^g \leq 4|\kappaChHodge(\cA)| \sum_{g > G} (|\hbar|/(2\pi)^2)^g = 4|\kappaChHodge(\cA)| \cdot (|\hbar|/(2\pi)^2)^{G+1} / (1 - |\hbar|/(2\pi)^2)$, and the result follows by simplification.
\end{proof}

\subsubsection{Scalar-shadow comparison with selected 3-manifold saddles}

\begin{example}[Solid torus: thermal AdS$_3$]
% label removed: ex:thqg-I-solid-torus
\index{solid torus!thermal AdS}
The solid torus $D^2 \times S^1$ with conformal boundary
$T^2 = S^1 \times S^1$ and modular parameter~$\tau$ is the thermal
AdS$_3$ saddle in the usual comparison.  The boundary-CFT character is
\[
Z_{\mathrm{char}}(\cA;\tau)
= \Tr_{\mathcal{H}_q(\cA)} q^{L_0-c/24},
\qquad q=e^{2\pi i\tau}.
\]
For $\cA=\mathrm{Vir}_c$, the vacuum-character factor
\[
q^{-c/24}\prod_{n=2}^{\infty}(1-q^n)^{-1}
\]
agrees with the holomorphic one-loop graviton determinant of
\cite{GMY08} after the thermal-AdS boundary-mode normalization has
been fixed.  The scalar shadow value
$F_1=c/48=\kappaChHodge(\mathrm{Vir}_c)/24$ is the Hodge trace; it is
not extracted from the nonconstant oscillator product.
\end{example}

\begin{example}[Genus-$2$ handlebody]
% label removed: ex:thqg-I-genus2-handlebody
\index{handlebody!genus-2}
The genus-$2$ handlebody $H_2$ has conformal boundary~$\Sigma_2$ and
is a candidate bulk filling after a handlebody saddle prescription has
been chosen.  A physical handlebody partition function requires a
conformal-block pairing, a measure or sum over intermediate labels,
and a metric normalization of the handlebody action.  With such data
one can compare the boundary block norm
\[
Z^{\mathrm{block}}_{H_2}(\Omega)
=
\sum_{\alpha} |\mathcal F_\alpha(\Omega)|^2
\]
with a large-\(c\) vacuum-block saddle ansatz of the form
\[
|\det(\operatorname{Im}\Omega)|^{c/2}
\exp\!\bigl(-c\,S_{H_2}^{\mathrm{met}}(\Omega)\bigr).
\]
Vacuum dominance and the metric action \(S_{H_2}^{\mathrm{met}}\) are
part of the physical comparison datum, not consequences of the scalar
shadow theorem.  The scalar value
\[
F_2^{\mathrm{scal}}(\mathrm{Vir}_c)
=\frac{7c}{11520}
\]
is the connected genus-$2$ Hodge-trace shadow coefficient on the
uniform-weight scalar lane.  It is not a computed two-loop handlebody
determinant and not a physical handlebody partition function.
\end{example}

\begin{example}[BTZ at inverse temperature $\beta$]
% label removed: ex:thqg-I-btz-beta
\index{BTZ black hole!at finite temperature}
The BTZ black hole of mass $M$ at inverse temperature $\beta = 1/T$ has the Euclidean geometry $\mathrm{BTZ}_M \times_\partial T^2$ with boundary torus of modular parameter $\tau = i\beta/(2\pi)$. If the Brown--Henneaux/Cardy dictionary and saddle comparison are supplied, the physical partition function has a sector of the form $Z_{\mathrm{BTZ}}(\beta) = e^{-\beta M} \cdot Z_{\mathrm{pert}}(\tau)$. At the Hawking--Page temperature $\beta_{\mathrm{HP}} = 2\pi \ell$, the BTZ and thermal saddles have equal classical action. The scalar shadow free energies computed here are candidate perturbative corrections to that comparison; the leading scalar coefficient is $F_1(\mathrm{Vir}_c) = c/48 = \ell/(32G)$. A shift of the physical transition temperature requires the separate saddle/measure calculation.
\end{example}

\subsubsection{The two-dimensional convergence theorem}

\begin{theorem}[Completed degree-summed shadow two-parameter convergence; \ClaimStatusProvedHere]
\label{thm:thqg-I-2d-convergence}
\index{convergence!two-dimensional}
\index{shadow partition series!two-parameter convergence}
For a modular Koszul chiral algebra $\cA$ satisfying HS-sewing:
\begin{enumerate}[label=\textup{(\roman*)}]
\item For each fixed genus $g$, the fixed-genus degree-summed shadow
expression
\begin{equation}% label removed: eq:thqg-I-2d-convergence
Z_g^{\mathrm{sh,deg}}(T;\,t) \;:=\; \sum_{r=0}^\infty t^r
\int_{\overline{\mathcal{M}}_{g,n(g,r)}} \operatorname{Sh}_{g,r}(\Theta_\cA)
\end{equation}
is a polynomial in~$t$ when $\cA$ has finite support shadow depth. At infinite
support shadow depth it is the completed shadow-degree series. It converges
absolutely on $|t|\leq 1$ under the estimate
\[
\sum_{r=0}^\infty
\left|
\int_{\overline{\mathcal{M}}_{g,n(g,r)}}
\operatorname{Sh}_{g,r}(\Theta_\cA)
\right| < \infty .
\]
More generally, a bound
\[
\left|
\int_{\overline{\mathcal{M}}_{g,n(g,r)}}
\operatorname{Sh}_{g,r}(\Theta_\cA)
\right|
\leq C_g R^{-r} r^{-1-\epsilon}
\]
gives absolute convergence for $|t|<R$.
\item If $\cA$ is Gaussian \textup{(}\(r^\Theta_{\max}=2\)\textup{)}, then the
degree-summed shadow two-parameter series
\[
Z_{\mathrm{sh}}^{\mathrm{deg}}(T;\,\hbar,t)
\;:=\;
\sum_{g=0}^\infty \hbar^g Z_g^{\mathrm{sh,deg}}(T;\,t)
\]
is independent of~$t$ and converges absolutely for $|\hbar|<4\pi^2$,
with meromorphic continuation given by the scalar shadow closed form.
\item Apart from such degree estimates and the Gaussian scalar lane, a
bidisk for the completed degree-summed shadow genus expansion does not
follow from HS-sewing or from recursive existence.
\end{enumerate}
\end{theorem}

\begin{proof}
Part~(i) is the distinction between the completed tower and a normed
shadow trace. Finite support shadow depth makes the sum finite in~$r$. At
infinite support shadow depth the inverse limit is formal in the degree
filtration; the displayed summability hypotheses are precisely the
estimates that make the degree-weighted shadow trace into an analytic
function of~$t$.

Part~(ii) is the Gaussian case of
Theorem~\ref{thm:thqg-I-full-convergence}\textup{(iii)}: only the scalar
shadow degree survives, so the $t$-dependence disappears and the scalar
genus series inherits the radius $4\pi^2$ and meromorphic continuation
from Theorem~\ref{thm:thqg-I-absolute-convergence}.

Part~(iii) follows because HS-sewing controls internal conformal-weight
sums at fixed degree, while
Theorem~\ref{V1-thm:recursive-existence} controls the completed
degree filtration. Neither assertion is a uniform estimate in
$(g,r)$.
\end{proof}

\subsubsection{Factored bounds for the double series}

The mixed class has a numerical bidisk exactly when the full double
series
\begin{equation}\label{eq:double-series-def}
Z(\hbar,t) \;:=\; \sum_{g=0}^\infty \sum_{r=0}^\infty
Z_{g,r}\,\hbar^{2g}\,t^r,
\qquad
Z_{g,r} \;:=\; \int_{\overline{\mathcal{M}}_{g,n(g,r)}}
\operatorname{Sh}_{g,r}(\Theta_\cA),
\end{equation}
converge absolutely on a bidisk
$\{|\hbar| < R_\hbar\} \times \{|t| < R_t\}$? On the scalar lane
($r = 0$), the radius is $R_\hbar = 4\pi^2$. On the shadow tower at
fixed genus~$0$, the radius is $R_t = 1/\rho$ for
$c > c^* \approx 6.125$. The simultaneous estimate is the
factored bound
\begin{equation}\label{eq:factored-bound-conj}
|Z_{g,r}| \;\leq\; C \cdot (2\pi)^{-2g} \cdot \rho^r \cdot r^{-5/2},
\end{equation}
which would imply absolute convergence on
$\{|\hbar| < 4\pi^2\} \times \{|t| < 1/\rho\}$ for $c > c^*$.
The estimates below isolate the contact contribution, the fixed-degree
conditional estimate, and the graph-growth obstruction.

\begin{proposition}[Contact contribution from tautological decay; \ClaimStatusProvedHere]
\label{prop:contact-contribution-bound}
\index{factored bound!contact contribution}
\index{stable graph!single-vertex}
For the Virasoro algebra $\mathrm{Vir}_c$ with $c>c^*$, fix the
shadow degree $r$. Suppose the single-vertex contact class factors as
\[
Z_{g,r}^{\mathrm{contact}}
=
S_r\,I_{g,r}^{\mathrm{ct}},
\]
where $S_r$ is the degree-$r$ coefficient of
$\sqrt{Q_{\mathrm{Vir}}}$ and the tautological integral satisfies
\begin{equation}\label{eq:contact-tautological-decay}
|I_{g,r}^{\mathrm{ct}}|
\leq
B_r(2\pi)^{-2g}
\qquad(g\geq 1).
\end{equation}
Then
\begin{equation}\label{eq:contact-bound}
\bigl|Z_{g,r}^{\mathrm{contact}}\bigr|
\leq
C_0B_r\,\rho^r r^{-5/2}(2\pi)^{-2g}.
\end{equation}
If $B_r\leq B$ uniformly in $r$, the contact sector satisfies the
factored bound~\eqref{eq:factored-bound-conj}.
\end{proposition}

\begin{proof}
The shadow-radius estimate gives
$|S_r|\leq C_0\rho^r r^{-5/2}$. Multiplication by
\eqref{eq:contact-tautological-decay} gives
\eqref{eq:contact-bound}. The uniform assertion is the same
inequality with $B_r$ replaced by~$B$.
\end{proof}

\begin{proposition}[Multi-vertex graph bound; \ClaimStatusProvedHere]
\label{prop:multi-vertex-bound}
\index{factored bound!multi-vertex contribution}
\index{stable graph!multi-vertex bound}
For a stable graph $\Gamma$ of genus $g$ with $r$ external
legs, $v$ vertices, and $e$ internal edges ($e = v + g - 1$
for connected graphs), the graph amplitude satisfies
\begin{equation}\label{eq:multi-vertex-amplitude-bound}
|A_\Gamma| \;\leq\;
\frac{1}{|\Aut(\Gamma)|}
\prod_{j=1}^v |S_{r_j}|
\cdot \prod_{i=1}^e \frac{1}{|\kappaChHodge(\mathrm{Vir}_c)|}
\cdot
\prod_{j=1}^v I_{g_j, n_j}\,,
\end{equation}
where $r_j$ is the degree at vertex $j$,
$g_j$ is the genus label, $n_j = r_j + 2e_j$ is the
total valence \textup{(}$e_j$ edges incident to vertex $j$\textup{)},
$\kappaChHodge(\mathrm{Vir}_c) = c/2$, and
$I_{g_j,n_j} := \sum_{|d|=3g_j-3+n_j}
\langle \tau_{d_1} \cdots \tau_{d_{n_j}} \rangle_{g_j}$
is the total $\psi$-class intersection volume at vertex $j$.
\end{proposition}

\begin{proof}
Each vertex of $\Gamma$ carries a shadow coefficient $S_{r_j}$
from the $r_j$ external legs and a $\psi$-class integral
$I_{g_j, n_j}$ from the moduli space
$\overline{\mathcal{M}}_{g_j, n_j}$ at that vertex (here $n_j$
is the total valence: external legs plus half-edges from
internal edges). Each internal edge contracts two half-edges
via the propagator $P = 1/\kappaChHodge(\mathrm{Vir}_c)$ (the inverse Hessian in the
scalar sector). The amplitude factorises over vertices and
edges, and the automorphism symmetry factor divides.

The bound follows from the triangle inequality applied to the
integral representation: the moduli integral at each vertex is
bounded by the sum of absolute values of its $\psi$-class
intersection numbers, and the edge propagators contribute
$1/|\kappaChHodge(\mathrm{Vir}_c)|$ each.
\end{proof}

\begin{proposition}[Stable-graph count at fixed degree; \ClaimStatusProvedHere]
\label{prop:graph-count-polynomial}
\index{stable graph!counting bound}
For fixed $(g,r)$ the set $\mathrm{Stab}(g,r)$ is finite. A crude
raw count is
\begin{equation}\label{eq:graph-count-fixed-r}
|\mathrm{Stab}(g,r)|
\;\leq\;
\sum_{v=1}^{2g-2+r}
(2g+r)^{2(3g+r-3)}
\binom{g+v}{v}
(2g+r)^r .
\end{equation}
For fixed $r$ this is $\exp(O(g\log g))$. The automorphism-weighted
mass
\[
M_{g,r}:=
\sum_{\Gamma\in \mathrm{Stab}(g,r)}
\frac{1}{|\Aut(\Gamma)|}
\]
is finite, but the Harer--Zagier alternating Euler characteristic is
not a positive upper bound for $M_{g,r}$.
\end{proposition}

\begin{proof}
A connected stable graph $\Gamma$ with genus $g$, $r$ external
legs, $v$ vertices, and $e$ edges satisfies
$g = \sum_j g_j + b_1(\Gamma)$ where $b_1 = e - v + 1$ is the
first Betti number and $g_j$ are the vertex genus labels. The
stability condition $2g_j - 2 + n_j > 0$ at each vertex (where
$n_j$ is the valence) constrains the combinatorics.

The key bound is on the number of vertices: $v \leq 2g - 2 + r$
(from $\sum_j (2g_j - 2 + n_j) = 2g - 2 + r$ with each summand
$\geq 1$). The number of edges is at most $3g-3+r$. For fixed
$v$, a graph with at most this many edges is bounded by
$(v^2)^{3g-3+r}$, hence by
$(2g+r)^{2(3g+r-3)}$. The genus labels are bounded by weak
compositions of genus among the vertices, hence by
$\binom{g+v}{v}$, and the $r$ legs have at most
$v^r\leq (2g+r)^r$ assignments. Summing over $v$ gives
\eqref{eq:graph-count-fixed-r}. The estimate is
$\exp(O(g\log g))$ for fixed $r$.

The positive weighted mass $M_{g,r}$ is finite because the raw set is
finite. The Harer--Zagier formula computes an alternating orbifold
Euler characteristic. Alternating cancellation is not an upper bound
for $\sum_\Gamma |\Aut(\Gamma)|^{-1}$.
\end{proof}

\begin{theorem}[Fixed-degree factored bound under weighted graph summability; \ClaimStatusProvedHere]
\label{thm:factored-bound-fixed-r}
\index{factored bound!fixed degree}
\index{convergence!factored bound}
For the Virasoro algebra $\mathrm{Vir}_c$ with $c > c^*$
\textup{(}so that $\rho < 1$ and the shadow obstruction tower converges\textup{)},
fix $r$ and assume the contact tautological decay
\eqref{eq:contact-tautological-decay}. Then the single-vertex
degree-$r$ contact amplitude satisfies
\[
\bigl|Z_{g,r}^{\mathrm{contact}}\bigr|
\;\leq\;
C_r^{\mathrm{ct}}\,(2\pi)^{-2g}\rho^r .
\]
If the non-contact stable-graph sector satisfies the weighted
summability estimate
\begin{equation}\label{eq:weighted-graph-summability}
\sum_{\substack{\Gamma\in \mathrm{Stab}(g,r)\\ v(\Gamma)\geq 2}}
\frac{1}{|\Aut(\Gamma)|}
\prod_j |S_{r_j}|
|\kappaChHodge(\mathrm{Vir}_c)|^{-e(\Gamma)}
\prod_j I_{g_j,n_j}
\;\leq\;
C_r^{\mathrm{mv}}\,(2\pi)^{-2g}\rho^r ,
\end{equation}
then the full genus-$g$ degree-$r$ amplitude satisfies
\begin{equation}\label{eq:factored-bound-fixed-r}
|Z_{g,r}| \;\leq\; C_r \cdot (2\pi)^{-2g} \cdot \rho^r
\end{equation}
for all $g \geq 1$, with $C_r=C_r^{\mathrm{ct}}+C_r^{\mathrm{mv}}$.
\end{theorem}

\begin{proof}
The contact estimate is Proposition~\ref{prop:contact-contribution-bound}.
The full amplitude is the sum of its contact part and the non-contact
stable-graph sector. Proposition~\ref{prop:multi-vertex-bound} gives
the absolute value of the latter by the left hand side of
\eqref{eq:weighted-graph-summability}. The displayed summability
estimate therefore gives
\[
|Z_{g,r}|
\leq
C_r^{\mathrm{ct}}(2\pi)^{-2g}\rho^r
+
C_r^{\mathrm{mv}}(2\pi)^{-2g}\rho^r ,
\]
which is \eqref{eq:factored-bound-fixed-r}.
\end{proof}

\begin{corollary}[Primitive factored bound for $r \leq 4$ under tautological decay; \ClaimStatusProvedHere]
\label{cor:factored-bound-r4}
\index{factored bound!low degree}
For $\mathrm{Vir}_c$ with $c > c^*$, assume
\eqref{eq:contact-tautological-decay} in degrees $2$ and~$4$. Then
the scalar term and the single-vertex low-degree primitives satisfy
the fixed-degree factored bound, where $Z_{g,2}^{\mathrm{prim}}$
denotes the single-vertex spectral primitive:
\begin{enumerate}[label=\textup{(\roman*)}]
\item $r = 0$ \textup{(}scalar, uniform-weight\textup{)}: $Z_{g,0} = F_g =
 \kappaChHodge(\mathrm{Vir}_c) \cdot \lambda_g^{\mathrm{FP}} \leq 4|\kappaChHodge(\mathrm{Vir}_c)| \cdot
 (2\pi)^{-2g}$ \textup{(}the proved effective
 bound~\eqref{V1-eq:thqg-I-effective-bound}\textup{)}.
\item $r = 2$ \textup{(}spectral primitive\textup{)}:
$|Z_{g,2}^{\mathrm{prim}}| \leq
 C_2 \cdot (2\pi)^{-2g} \cdot \rho^2$.
\item $r = 4$ \textup{(}quartic contact primitive\textup{)}:
$|Z_{g,4}^{\mathrm{contact}}|
 \leq C_4 \cdot (2\pi)^{-2g} \cdot \rho^4$, where the dominant
 term is
 $Q^{\mathrm{contact}}_{\mathrm{Vir}}
 \cdot \int_{\overline{\mathcal{M}}_{g,4}} [\psi\text{-class}]$
 with $Q^{\mathrm{contact}} = 10/(c(5c+22))$.
\end{enumerate}
The corresponding full low-degree amplitudes satisfy the same
inequalities under the weighted graph-summability estimate
\eqref{eq:weighted-graph-summability} for $r=2,4$.
\end{corollary}

\begin{proof}
At $r=0$ the assertion is the effective scalar bound. The
$r=2$ and $r=4$ estimates are the single-vertex cases of
Proposition~\ref{prop:contact-contribution-bound}. The last assertion
is Theorem~\ref{thm:factored-bound-fixed-r} applied in degrees
$2$ and~$4$.
\end{proof}

\begin{remark}[The obstacle to full bidisk convergence]
\label{rem:bidisk-obstacle}
\index{factored bound!obstacle to bidisk convergence}
Theorem~\ref{thm:factored-bound-fixed-r} gives
$|Z_{g,r}| \leq C_r (2\pi)^{-2g} \rho^r$ for each
\emph{fixed} $r$ only after the weighted stable-graph summability
estimate \eqref{eq:weighted-graph-summability}; the contact summand
has this bound under the tautological decay
\eqref{eq:contact-tautological-decay}. For full bidisk
convergence, one needs uniform control
$C_r \leq C \cdot r^{-5/2}$ as $r \to \infty$. The obstruction is the
growth of the graph count and the vertex intersection volumes: at
fixed $g$ and $r \to \infty$, the displayed graph bound is
combinatorial rather than polynomial
(Proposition~\ref{prop:thqg-I-graph-count}), and the vertex
intersection volumes $I_{g_j, n_j}$ grow factorially in the valence
$n_j$ by the Stirling-type bound on $\psi$-class integrals.

The conjectured $r^{-5/2}$ suppression would follow if the
shadow coefficients $S_r$ (which carry the algebraic data of
the Virasoro OPE) cancel the combinatorial growth of the
moduli-space integrals. This cancellation is visible in the
shadow obstruction tower at genus~$0$, where $\sqrt{Q_{\mathrm{Vir}}}$
produces the $r^{-5/2}$ decay as a consequence of the
branch-point structure of the quadratic polynomial
$Q_{\mathrm{Vir}}(t) = c^2 + 12ct + (180c + 872)/(5c+22)
\cdot t^2$. Uniform persistence of this cancellation over
stable graphs and genera is exactly the weighted summability
estimate.

The constant $C_r$ in
Theorem~\ref{thm:factored-bound-fixed-r} absorbs two
contributions: the multi-vertex weighted mass and the vertex
intersection volumes at fixed degree. The $r$-sum diverges if
$C_r$ grows faster than
$\rho^{-r} r^{5/2}$; then summing over $r$ would diverge.
The vertex contributions at multi-vertex graphs involve convolutions
of the shadow obstruction tower with itself: a
vertex of degree $r_j$ contributes
$|S_{r_j}| \leq C_0 \rho^{r_j} r_j^{-5/2}$, and the sum
over all vertex-degree assignments is a convolution of shadow
coefficients which converges if $\rho < 1$. This convolution estimate
does not control the positive stable-graph mass by itself.

\emph{Quantitative threshold.} The condition for the full
shadow-radius convergence is $\rho(c) < 1$, which is
$c > c^* \approx 6.125$, the same threshold as the genus-$0$ shadow
tower convergence. It is necessary for the bidisk bound and becomes
sufficient only together with the weighted graph-summability estimate.
For $c$ sufficiently large (so that
$\rho \ll 1$ and $|\kappaChHodge(\mathrm{Vir}_c)| = c/2 \gg 1$), the propagator
suppression $1/|\kappaChHodge(\mathrm{Vir}_c)|$ per edge and the shadow decay
$\rho^{r_j}$ per vertex provide ample margin.

The analytic assertions separate as follows:
\begin{enumerate}[label=\textup{(\alph*)}]
\item The scalar genus series has radius $4\pi^2$ on the
uniform-weight scalar lane.
\item The shadow obstruction tower exists as a completed inverse
limit by Theorem~\ref{V1-thm:recursive-existence}.
\item The contact and primitive low-degree summands satisfy factored
bounds under tautological decay; the full fixed-degree bound is the
implication \eqref{eq:contact-tautological-decay} and
\eqref{eq:weighted-graph-summability}
$\Rightarrow$ \eqref{eq:factored-bound-fixed-r}.
\item The full bidisk is equivalent to a uniform factored estimate in
both variables.
\end{enumerate}
\end{remark}

\begin{remark}[Absence of shadow counterterms: detailed mechanism]
% label removed: rem:thqg-I-no-uv-detailed
\index{UV finiteness!detailed mechanism}
Five structures remove counterterms from the algebraic shadow complex
at fixed genus and degree:
\begin{enumerate}[label=\textup{(\arabic*)}]
\item \emph{Compact integration domain.} All shadow integrals are over the Deligne--Mumford compactification $\overline{\mathcal{M}}_{g,n}$ or the Fulton--MacPherson compactification $\overline{\mathrm{FM}}_n(\Sigma_g)$.
\item \emph{Resolved diagonal.} The shadow representation $\operatorname{Sh}_{g,n}(\Theta_\cA)$ is a cohomology class on $\overline{\mathcal{M}}_{g,n}$, and the Fulton--MacPherson boundary records the coincident-point strata rather than discarding them.
\item \emph{Closed MC equation.} The identity $D_\cA^2 = 0$ holds at all genera without truncation. No counterterm is created by cutting the bar complex at finite order.
\item \emph{Finite graph sum.} At each genus $g$ and degree $r$, only finitely many graphs contribute.
\item \emph{Scalar diagonal decay.} On the uniform-weight scalar lane the diagonal free energies satisfy $F_g(\cA) \sim \kappaChHodge(\cA)/(2\pi)^{2g}$. Multi-weight cross-channel series and physical Borel singularities are separate analytic questions.
\end{enumerate}
\end{remark}

\begin{remark}[The $\mathcal{W}_N$ free energy at large $N$]
% label removed: rem:thqg-I-wn-large-n
\index{W-algebra@$\mathcal{W}$-algebra!large-$N$ limit}
For $\mathcal{W}_N(\mathfrak{sl}_N)$ in the 't~Hooft limit $N \to \infty$ with $\lambda = N/(k+N)$ fixed, the central charge is $c_N = (N-1) - N(N^2-1)(k+N-1)^2/(k+N) \sim -N^2\lambda$ for large $N$. The modular characteristic scales as $\kappaChHodge(\mathcal W_N) \sim N^2$, and the scalar genus-$g$ free energy as $F_g \sim N^2/(2\pi)^{2g}$. The scalar shadow closed form in this limit is
\[
Z_{\mathrm{sh}}^{\mathcal{W}_N,\mathrm{scal}}(\hbar) \;\approx\; N^2 \cdot \frac{\sqrt{\hbar}/2}{\sin(\sqrt{\hbar}/2)},
\]
which has the same scalar functional form as the finite-$N$ case but
with $\kappaChHodge$ growing quadratically in $N$. The scalar convergence
radius $4\pi^2$ is independent of $N$; the full non-scalar large-$N$
series also depends on the growth of the $\mathcal W_N$ shadow tower.
\end{remark}


% ======================================================================
% SUBSECTION 8: THE SELBERG ZETA CONNECTION AND SPECTRAL DETERMINANTS
% ======================================================================

\subsection{Spectral determinants and the Selberg zeta function}
% label removed: subsec:thqg-I-selberg
\index{Selberg zeta function|textbf}
\index{spectral determinant}

The Heisenberg scalar determinant is a spectral determinant on the
constant-curvature surface. Its Selberg form is obtained from the
leading coefficient of the Selberg zeta function at the zero forced by
the constant Laplace eigenfunction.

\subsubsection{The Selberg zeta function}

\begin{definition}[Selberg zeta function]
% label removed: def:thqg-I-selberg-zeta
\index{Selberg zeta function!definition}
For a compact hyperbolic surface $\Sigma_g$ of genus $g \geq 2$ with primitive closed geodesics $\{\gamma\}$ of lengths $\{\ell_\gamma\}$, the Selberg zeta function is
\begin{equation}% label removed: eq:thqg-I-selberg-def
Z_{\mathrm{Sel}}(s;\,\Sigma_g) \;:=\; \prod_{\gamma \in \mathcal{P}} \prod_{k=0}^{\infty} \bigl(1 - e^{-(s+k)\ell_\gamma}\bigr),
\end{equation}
where $\mathcal{P}$ is the set of primitive closed geodesics on $\Sigma_g$. The product converges absolutely for $\Re(s) > 1$ and extends meromorphically to $\mathbb{C}$.
\end{definition}

\begin{theorem}[Determinant formula; \ClaimStatusProvedElsewhere]
\label{thm:thqg-I-determinant-formula}
\index{Selberg zeta function!determinant formula}
The Selberg zeta function has a simple zero at $s=1$. Put
\[
Z_{\mathrm{Sel}}^*(1;\Sigma_g)
:=
\lim_{s\to 1}\frac{Z_{\mathrm{Sel}}(s;\Sigma_g)}{s-1}.
\]
The zeta-regularized determinant of the scalar Laplacian on $\Sigma_g$
is related to the Selberg zeta function by
\begin{equation}% label removed: eq:thqg-I-det-selberg
\det{}'_\zeta \Delta_{\Sigma_g}
\;=\;
C_g^{\mathrm{DP}}\,
Z_{\mathrm{Sel}}^*(1;\,\Sigma_g),
\end{equation}
where $C_g^{\mathrm{DP}}>0$ is the D'Hoker--Phong normalization
constant in curvature $-1$ normalization.
This is the D'Hoker--Phong formula~\cite{DP86}, building on the Selberg trace formula.
\end{theorem}

\begin{proof}
The Selberg trace formula expresses the spectral zeta function
$\zeta_\Delta(s) = \sum_{\lambda_n > 0} \lambda_n^{-s}$ in terms of
closed geodesics. Analytic continuation to $s=0$ gives
$\zeta_\Delta'(0)=-\log\det{}'_\zeta\Delta$. On the Selberg side the
zero eigenvalue of $\Delta$ is the simple zero of
$Z_{\mathrm{Sel}}(s)$ at $s=1$; hence the determinant is proportional
to the leading coefficient $Z_{\mathrm{Sel}}^*(1)$, not to the value
$Z_{\mathrm{Sel}}(1)$.
\end{proof}

\begin{corollary}[Heisenberg partition function via Selberg zeta; \ClaimStatusProvedHere]
\label{cor:thqg-I-heisenberg-selberg}
\index{Heisenberg!Selberg zeta representation}
The Heisenberg scalar shadow on a genus-$g$ hyperbolic surface admits
the representation
\begin{equation}% label removed: eq:thqg-I-heisenberg-selberg
Z_g^{\mathrm{scal}}(\mathcal{H}_k;\,\Sigma_g)
\;=\;
(\det\operatorname{Im}\Omega)^{-k/2}
\cdot
\bigl(C_g^{\mathrm{DP}}Z_{\mathrm{Sel}}^*(1;\Sigma_g)\bigr)^{-k/2}.
\end{equation}
For $\Re(s)>1$ the logarithm has the absolutely convergent
closed-geodesic expansion
\begin{equation}% label removed: eq:thqg-I-selberg-expansion
\log Z_{\mathrm{Sel}}(s)^{-k/2}
\;=\;
\frac{k}{2} \sum_{\gamma \in \mathcal{P}} \sum_{m=1}^{\infty}
\frac{1}{m}\cdot
\frac{e^{-ms\ell_\gamma}}{1-e^{-m\ell_\gamma}} .
\end{equation}
At $s=1$ the determinant uses analytic continuation and the leading
coefficient $Z_{\mathrm{Sel}}^*(1)$.
\end{corollary}

\begin{proof}
Substitute the determinant formula
(Theorem~\ref{thm:thqg-I-determinant-formula}) into the Heisenberg
scalar shadow (Theorem~\ref{prop:thqg-I-heisenberg-detail}). For
$\Re(s)>1$, the Selberg expansion follows from
$\log\prod_{k\geq0}(1-xq^k)
=-\sum_{m\geq1}x^m/(m(1-q^m))$ with
$x=e^{-s\ell_\gamma}$ and $q=e^{-\ell_\gamma}$. The formula at
$s=1$ is the analytically continued leading coefficient.
\end{proof}

\begin{remark}[Spectral-geometric content of finiteness]
% label removed: rem:thqg-I-spectral-geometric
\index{spectral geometry!finiteness}
The Heisenberg spectral determinant is finite because the non-zero
Laplace spectrum on a compact hyperbolic surface is discrete and
positive. The eigenvalues $\lambda_n>0$ for $n\geq1$, the spectral
zeta function converges for $\Re(s)>1$ by Weyl asymptotics, and
zeta regularization defines $\det{}'_\zeta\Delta$ by analytic
continuation at $s=0$. Compactness, the omitted zero mode, and the
Selberg leading coefficient are the spectral inputs.
\end{remark}

\subsubsection{The Polyakov formula and the Mumford form}

\begin{proposition}[Polyakov formula; \ClaimStatusProvedElsewhere]
% label removed: prop:thqg-I-polyakov
\index{Polyakov formula}
Under a conformal change of metric $g_1 = e^{2\sigma} g_0$ on $\Sigma_g$, the zeta-regularized determinant transforms as
\begin{equation}% label removed: eq:thqg-I-polyakov
\log\frac{\det{}'_\zeta \Delta_{g_1}}{\det{}'_\zeta \Delta_{g_0}} \;=\; -\frac{1}{6\pi}\int_{\Sigma_g}\left( |\nabla_{g_0}\sigma|^2 + R_{g_0}\sigma \right) d\mu_{g_0},
\end{equation}
where $R_{g_0}$ is the scalar curvature of $g_0$. The right side is the Liouville action.
\end{proposition}

\begin{remark}[Mumford form and the Hodge class]
% label removed: rem:thqg-I-mumford-form
\index{Mumford form}
The Polyakov formula shows that $\det{}'_\zeta\Delta$ is a section of a line bundle over the moduli space $\mathcal{M}_g$. The Mumford isomorphism identifies this line bundle with $\lambda^{\otimes 13}$, where $\lambda = \det\mathbb{E}$ is the determinant of the Hodge bundle. The first Chern class $c_1(\lambda) = \lambda_1$ is the Hodge class. The top class $\lambda_g=c_g(\mathbb{E})$ enters the Faber--Pandharipande number only on $\overline{\mathcal{M}}_{g,1}$ together with the degree-compensating insertion $\psi_1^{2g-2}$. The Mumford form $\mu_g$ is the canonical section of $\lambda^{\otimes 13}$ on $\mathcal{M}_g$, and the Belavin--Knizhnik factorization expresses the bosonic string integrand as $|\mu_g|^{-2} \cdot (\det\operatorname{Im}\Omega)^{-13}$.

For the Heisenberg at level $k$, the integrand is $|\mu_g|^{-2k/13} \cdot (\det\operatorname{Im}\Omega)^{-k}$, and the scalar shadow free energy is $k\lambda_g^{\mathrm{FP}}$.
\end{remark}

\subsubsection{Comparison with the bosonic string}

\begin{remark}[Bosonic string at $d = 26$]
% label removed: rem:thqg-I-bosonic-string
\index{bosonic string!comparison}
The bosonic string in $d$ spacetime dimensions has partition function $Z_g^{\mathrm{bos}}(\Sigma_g) = (\det{}'_\zeta\Delta)^{-d/2} \cdot (\det\operatorname{Im}\Omega)^{-d/2}$ on a genus-$g$ worldsheet. At $d = 26$, the conformal anomaly cancels and the integrand is modular invariant. The shadow free energy of the $d$-dimensional free-boson scalar lane is $F_g(V_{\mathrm{bos},d}) = d \cdot \lambda_g^{\mathrm{FP}}$ (since $\kappaChHodge(V_{\mathrm{bos},d})=d$).

The scalar shadow generating function differs from the bosonic string in two ways: (i)~the scalar genus expansion converges on the proved uniform-weight lane, while the string coefficients $F_g^{\mathrm{string}} \sim (2g-2)! \cdot \alpha'^{-g}$ grow factorially; (ii)~the shadow obstruction tower provides additional degree-dependent corrections beyond the scalar channel, which are absent in the free boson theory.
\end{remark}

\subsubsection{Analytic continuation of the scalar shadow closed form}

\begin{proposition}[Pole structure of the meromorphic continuation; \ClaimStatusProvedHere]
% label removed: prop:thqg-I-pole-structure
\index{scalar shadow closed form!pole structure}
The meromorphic scalar shadow closed form
\[
Z_{\mathrm{sh}}^{\mathrm{scal}}(\cA;\,\hbar)
= \kappaChHodge(\cA)\cdot
\frac{\sqrt{\hbar}/2}{\sin(\sqrt{\hbar}/2)}
\]
has:
\begin{enumerate}[label=\textup{(\roman*)}]
\item Simple poles at $\hbar_n = 4\pi^2 n^2$ for $n = 1, 2, 3, \ldots$
\item Residues $\Res_{\hbar = \hbar_n} Z_{\mathrm{sh}}^{\mathrm{scal}} = (-1)^{n} \cdot 8\kappaChHodge(\cA)n^2\pi^2$.
\item If $\kappaChHodge(\cA)\neq0$, no zeros on $\mathbb{C}$; if
$\kappaChHodge(\cA)=0$, the scalar shadow closed form is identically
zero.
\item On every vertical strip away from the pole set, $|Z_{\mathrm{sh}}^{\mathrm{scal}}(\sigma + i t)| \leq C(\sigma) \cdot e^{|t|^{1/2}/2}$ for $|\sigma| \leq R$ and $|t| \to \infty$.
\end{enumerate}
\end{proposition}

\begin{proof}
(i)--(ii): The zeros of $\sin(\sqrt{\hbar}/2)$ occur at $\sqrt{\hbar}/2 = n\pi$, i.e., $\hbar = 4\pi^2 n^2$. The residue at $\hbar_n$ is computed from the Laurent expansion:
\begin{align*}
\frac{\sqrt{\hbar}/2}{\sin(\sqrt{\hbar}/2)} &= \frac{n\pi + \epsilon/(8n\pi) + O(\epsilon^2)}{\sin(n\pi + \epsilon/(8n\pi) + \cdots)} \\
&= \frac{n\pi + O(\epsilon)}{(-1)^n \epsilon/(8n\pi) + O(\epsilon^2)} \\
&= \frac{(-1)^{-n} \cdot 8n^2\pi^2}{\epsilon} + O(1),
\end{align*}
where $\epsilon = \hbar - 4\pi^2 n^2$ and
$\sqrt{4\pi^2 n^2 + \epsilon} = 2n\pi + \epsilon/(4n\pi) + O(\epsilon^2)$.
Hence $\Res = \kappaChHodge(\cA) \cdot (-1)^{n} \cdot 8n^2\pi^2$.

(iii): At $\hbar = -4\pi^2 n^2$, $\sin(\sqrt{-4\pi^2 n^2}/2) = \sin(i n\pi) = i\sinh(n\pi) \neq 0$, and $\sqrt{-4\pi^2 n^2}/2 = i n\pi$, so $Z_{\mathrm{sh}}^{\mathrm{scal}} = \kappaChHodge(\cA) \cdot in\pi / (i\sinh(n\pi)) = \kappaChHodge(\cA) \cdot n\pi/\sinh(n\pi)$. At $\hbar=0$ the numerator and denominator vanish to first order, and the singularity is removable with value $\kappaChHodge(\cA)$. Thus $Z_{\mathrm{sh}}^{\mathrm{scal}}$ has no zeros when $\kappaChHodge(\cA)\neq0$.

(iv): For $|\hbar| \to \infty$ in the direction $\arg(\hbar) = \theta \neq 0, \pi$, $|\sin(\sqrt{\hbar}/2)| \geq C \cdot e^{|\operatorname{Im}\sqrt{\hbar}|/2}$, and $|\sqrt{\hbar}/2| \leq |\hbar|^{1/2}/2$, giving the bound.
\end{proof}

\begin{remark}[Partial fraction expansion]
% label removed: rem:thqg-I-partial-fraction
\index{scalar shadow closed form!partial fraction}
The meromorphic function
$Z_{\mathrm{sh}}^{\mathrm{scal}}(\cA;\,\hbar)
= \kappaChHodge(\cA) \cdot
\sqrt{\hbar}/(2\sin(\sqrt{\hbar}/2))$ admits the partial fraction expansion
\begin{equation}% label removed: eq:thqg-I-partial-fraction
Z_{\mathrm{sh}}^{\mathrm{scal}}(\cA;\,\hbar) \;=\; \kappaChHodge(\cA) \left( 1 + 2\sum_{n=1}^{\infty} \frac{(-1)^{n} \hbar}{\hbar - 4\pi^2 n^2} \right),
\end{equation}
which converges for $\hbar \notin \{4\pi^2 n^2 : n \geq 1\}$. This representation makes the pole structure manifest and provides an efficient numerical evaluation for $\hbar$ away from the poles. The partial fraction expansion is equivalent to the Mittag-Leffler theorem applied to the meromorphic function $\sqrt{\hbar}/(2\sin(\sqrt{\hbar}/2))$.
\end{remark}

\subsubsection{The entropy of the shadow free energy}

\begin{definition}[Shadow entropy]
% label removed: def:thqg-I-shadow-entropy
\index{shadow entropy}
Assume $\kappaChHodge(\cA)\neq0$ and choose a branch of the logarithm
on a simply connected domain avoiding the scalar pole set
$\{4\pi^2n^2:n\geq1\}$. The \emph{shadow entropy} of a modular Koszul
chiral algebra $\cA$ on this domain is
\begin{equation}% label removed: eq:thqg-I-shadow-entropy
S_{\mathrm{sh}}(\cA;\,\hbar)
\;:=\;
\log Z_{\mathrm{sh}}^{\mathrm{scal}}(\cA;\,\hbar)
\;=\;
\log\kappaChHodge(\cA)
+ \log\frac{\sqrt{\hbar}/2}{\sin(\sqrt{\hbar}/2)}\,.
\end{equation}
If $\kappaChHodge(\cA)>0$ and $\hbar$ approaches the first scalar pole
from below on the real axis, $S_{\mathrm{sh}} \to +\infty$
\textup{(}logarithmic divergence\textup{)}:
\begin{equation}% label removed: eq:thqg-I-entropy-divergence
S_{\mathrm{sh}}(\cA;\,4\pi^2 - \epsilon) \;\sim\; -\log\epsilon + \log(4\pi^2\kappaChHodge(\cA)) \quad\text{as } \epsilon \to 0^+.
\end{equation}
\end{definition}

\begin{remark}[Bekenstein--Hawking interpretation]
% label removed: rem:thqg-I-bekenstein-hawking
\index{Bekenstein--Hawking entropy}
In the Brown--Henneaux comparison context with
$\kappaChHodge(\mathrm{Vir}_c)=3\ell/(4G_N)$ and
$\hbar = 2G_N/(3\ell)$, the shadow entropy at
$\hbar = 2G_N/(3\ell)$ is
\[
S_{\mathrm{sh}} \;=\; \log\frac{3\ell}{4G_N} + \log\frac{\sqrt{G_N/(6\ell)}}{\sin(\sqrt{G_N/(6\ell)})} \;\approx\; \log\frac{3\ell}{4G_N} + \frac{G_N}{72\ell} + O(G_N^2/\ell^2),
\]
where the leading term $\log(3\ell/(4G_N))$ is the logarithm of the
Brown--Henneaux scalar normalization, and the corrections are of order
$G_N/\ell$. The Bekenstein--Hawking entropy of the BTZ black hole of mass
$M$ is $S_{\mathrm{BH}} = \pi r_+/(2G_N) =
2\pi\sqrt{M\ell/(2G_N)}$, which is parametrically larger than
$S_{\mathrm{sh}}$ for macroscopic black holes ($M\ell/G_N \gg 1$). The
shadow entropy is a logarithm of the scalar shadow series, not the black
hole entropy.
\end{remark}

\subsubsection{Summary of quantitative finiteness results}

\begin{table}[ht]
\centering
\caption{Quantitative finiteness data for the standard landscape at genus $g = 1, 2, 3$. Higher-genus scalar entries and radii are shown only where proved.}
% label removed: tab:thqg-I-quantitative
\index{perturbative finiteness!quantitative data}
\renewcommand{\arraystretch}{1.4}
{\footnotesize
\begin{tabular}{|l|c|c|c|c|c|}
\hline
\textbf{Family $\cA$} & $\boldsymbol{\kappaChHodge}$ & $\boldsymbol{F_1}$ & $\boldsymbol{F_2}$ & $\boldsymbol{F_3}$ & $\boldsymbol{R_{\mathrm{sh}}^{\mathrm{scal}}}$ \\
\hline
$\mathcal{H}_1$ & $1$ & $1/24$ & $7/5760$ & $31/967680$ & $4\pi^2$ \\
\hline
$\widehat{\mathfrak{sl}}_{2,1}$ & $9/4$ & $3/32$ & $21/7680$ & $93/1290240$ & $4\pi^2$ \\
\hline
$\mathrm{Vir}_{26}$ & $13$ & $13/48$ & $91/11520$ & $403/1935360$ & $4\pi^2$ \\
\hline
$\beta\gamma$ & $1$ & $1/24$ &; &; &; \\
\hline
$V_{E_8}$ & $8$ & $1/3$ & $7/720$ & $31/120960$ & $4\pi^2$ \\
\hline
$V_{\mathrm{Leech}}$ & $24$ & $1$ & $7/240$ & $31/40320$ & $4\pi^2$ \\
\hline
\end{tabular}}
\end{table}

\begin{remark}[The Leech lattice]
% label removed: rem:thqg-I-leech
\index{Leech lattice}
The Leech lattice VOA $V_{\mathrm{Leech}}$ of rank $24$ has
$\kappaChHodge(V_{\mathrm{Leech}}) = 24$. Its oscillator scalar projection has genus-$1$ free
energy $F_1=1$ and scalar series
$24 \cdot \sqrt{\hbar}/(2\sin(\sqrt{\hbar}/2))$, with radius
$4\pi^2$ and scalar poles at $\hbar=4\pi^2 n^2$. The full Leech
lattice theory also carries lattice theta data; the genus-$2$ scalar
number $7/240$ is the oscillator contribution, while cusp-form
corrections live in the arithmetic lattice sector. The Leech rank
$24$ is the maximal modular characteristic among standard even
unimodular lattice examples of rank at most~$24$.
\end{remark}

\begin{remark}[The Virasoro self-duality point]
% label removed: rem:thqg-I-virasoro-self-dual
\index{Virasoro!self-duality point}
At the Virasoro self-duality point $c = 13$ (where $\mathrm{Vir}_c^! \cong \mathrm{Vir}_{26-c} = \mathrm{Vir}_{13}$), the modular characteristic is $\kappaChHodge(\mathrm{Vir}_{13}) = 13/2$ and the shadow free energies are $F_g(\mathrm{Vir}_{13}) = (13/2)\lambda_g^{\mathrm{FP}}$ \textup{(}uniform-weight\textup{)}. This is NOT the special central charge $c = 26$ where the bosonic string becomes critical; $c = 13$ is distinguished by the property that the Koszul dual is isomorphic to the original algebra. The free energies at $c = 13$:
\begin{align*}
F_1(c=13) &= 13/48 \approx 0.271, \\
F_2(c=13) &= 91/11520 \approx 0.0079, \\
F_3(c=13) &= 403/1935360 \approx 0.000208.
\end{align*}
\end{remark}

\begin{proposition}[Finiteness for tensor-product theories; \ClaimStatusProvedHere]
% label removed: prop:thqg-I-tensor-product
\index{tensor product!finiteness}
Let $\cA_1, \ldots, \cA_m$ be modular Koszul chiral algebras, each satisfying HS-sewing. Then:
\begin{enumerate}[label=\textup{(\roman*)}]
\item The tensor product $\cA = \cA_1 \otimes \cdots \otimes \cA_m$ satisfies HS-sewing, with
 \begin{equation}% label removed: eq:thqg-I-tensor-hs
 \|m_q^{\cA}\|_{\HS}^2 \;=\; \prod_{j=1}^m \|m_q^{\cA_j}\|_{\HS}^2.
 \end{equation}
\item The modular characteristic is additive: $\kappaChHodge(\cA) = \kappaChHodge(\cA_1) + \cdots + \kappaChHodge(\cA_m)$.
\item The universal genus-$1$ scalar free energy is additive:
 $F_1(\cA) = F_1(\cA_1) + \cdots + F_1(\cA_m)$.
\item If the tensor product lies on the proved scalar lane
 \textup{(}for example, in the Gaussian families treated above\textup{)},
 then the higher-genus scalar free energies are additive and the
 additive scalar shadow genus series has the closed form
 \begin{equation}% label removed: eq:thqg-I-tensor-partition
 Z_{\mathrm{sh}}^{\mathrm{scal}}(\cA;\,\hbar) \;=\; \sum_{j=1}^m \kappaChHodge(\cA_j) \cdot \frac{\sqrt{\hbar}/2}{\sin(\sqrt{\hbar}/2)}.
 \end{equation}
\end{enumerate}
\end{proposition}

\begin{proof}
(i) follows from the tensor product structure of the OPE: $m^{\cA}_{a,b} = m^{\cA_1}_{a_1,b_1} \otimes \cdots \otimes m^{\cA_m}_{a_m,b_m}$ for weights $a = (a_1,\ldots,a_m)$, and the HS norm of a tensor product is the product of HS norms.

(ii) is the additivity of the modular characteristic (Theorem~D).

(iii) is the universal genus-$1$ identity
$F_1(\cA)=\kappaChHodge(\cA)/24$ together with~(ii).

(iv) uses the additional hypothesis that the tensor product lies on the
proved scalar lane, so the scalar closed form is available and linear in
$\kappaChHodge(\cA)$.
\end{proof}

\begin{remark}[Scope of finiteness]
% label removed: rem:thqg-I-complete-finiteness
\index{perturbative finiteness!complete statement}
For every modular Koszul chiral algebra $\cA$ in the standard
landscape \textup{(}Table~\textup{\ref{tab:thqg-I-finiteness-summary})}
and for every tensor product, coset, or orbifold that preserves
modular Koszulness and HS-sewing, the genus-$g$ scalar shadow
\[
\int_{\overline{\mathcal{M}}_g}
\operatorname{Sh}_{g,0}(\Theta_\cA)
\]
is finite for every fixed genus. On the proved scalar lane, the scalar
genus series converges
absolutely for $|\hbar| < 4\pi^2$ and admits a meromorphic continuation
to $\mathbb{C}$ with explicit pole structure. For mixed or multi-weight
families outside that lane, the fixed-genus scalar shadow and the
universal genus-$1$ scalar term remain the unconditional statements.
Algebraic shadow finiteness is a bar-complex statement; physical BV
renormalisation is a further realisation datum.
\end{remark}


% ======================================================================
% SUBSECTION 8b: MULTI-WEIGHT SCALAR FACTORIZATION AT GENUS 2
% ======================================================================

\subsection{Multi-weight scalar factorization at genus $2$}
\label{subsec:thqg-I-multi-weight-genus-2}
\index{multi-weight!scalar factorization}
\index{shadow free energy!multi-weight}
\index{stable graph!genus 2}

The scalar factorisation
$F_g^{\mathrm{sc}}(\cA)=\kappaChHodge(\cA)\lambda_g^{\mathrm{FP}}$
is proved on the scalar lane for algebras
with generators of uniform conformal weight
(Theorem~\ref{thm:thqg-I-perturbative-finiteness}(ii)). For
multi-weight algebras such as $\mathcal{W}_3$ (generators $T$ of
weight~$2$, $W$ of weight~$3$), factorisation is universal at
genus~$1$ ($F_1(\cA)=\kappaChHodge(\cA)/24$); higher genus requires separate
analysis. Direct stable-graph computation at genus~$2$ identifies
the mechanism that controls all genera.

\subsubsection{The modular characteristic of a multi-weight algebra}

\begin{definition}[Modular characteristic as Hessian trace]
\label{def:kappa-hessian-trace}
\index{modular characteristic!Hessian trace}
Let $\cA$ be a strongly finitely generated chiral algebra with
strong generators $\phi^1, \ldots, \phi^N$ of conformal weights
$h_1 \leq \cdots \leq h_N$. The \emph{modular characteristic} is
\begin{equation}\label{eq:kappa-hessian}
\kappaChHodge(\cA) \;=\; \sum_{i=1}^N \kappa_i(\cA),
\qquad
\kappa_i(\cA) \;:=\; \frac{1}{2}\,
\Res_{z=0} z^{-1}\, C^{i,1}_{i,1;\,i,1}(z),
\end{equation}
where $C^{i,1}_{i,1;\,i,1}(z)$ is the coefficient of $\phi^i$ in the
$\phi^i(z)\phi^i(0)$ OPE, evaluated at the unit-normalised
basis vector (i.e., the leading-order
self-contraction of $\phi^i$). Each summand $\kappa_i$ is the
\emph{sectoral curvature}: the genus-$1$ tadpole of the $i$-th
generator.
\end{definition}

\begin{example}[$\mathcal{W}_3$]
\label{ex:kappa-w3}
\index{W-algebra@$\mathcal{W}$-algebra!modular characteristic!$\mathcal{W}_3$}
For $\mathcal{W}_3$ with generators $T$ (weight~$2$) and $W$
(weight~$3$):
\begin{align}
\kappa_T &= c/2 &&\text{(from the $T(z)T(0)$ quartic pole)},
\label{eq:kappa-T}\\
\kappa_W &= c/3 &&\text{(from the $W(z)W(0)$ sextic pole
normalisation)}, \label{eq:kappa-W}\\
\kappaChHodge(\mathcal{W}_3) &= c/2 + c/3 = 5c/6.
\label{eq:kappa-w3}
\end{align}
More generally, $\kappaChHodge(\mathcal{W}_N) = c \cdot
\sum_{s=2}^N 1/s = c(H_N - 1)$ where $H_N$ is the $N$-th
harmonic number.
\end{example}

\subsubsection{Stable graphs of genus $2$}

\begin{lemma}[Genus-$2$ stable graphs; \ClaimStatusProvedHere]
\label{lem:genus-2-stable-graphs}
\index{stable graph!genus 2 classification}
There are exactly seven topological types of connected stable
graphs $\Gamma$ with loop genus $g(\Gamma) = 2$ and no external
legs ($n = 0$). Six are edge-bearing, under the at-least-one-edge
convention (at least one propagator in the Feynman expansion);
the seventh is the smooth stratum $\Gamma_0$
(one vertex, $g_v = 2$, no edges). %: total count is 7
The six edge-bearing types are:
\begin{enumerate}[label=\textup{(\Roman*)}]
\item \emph{Sunrise} ($\Gamma_{\mathrm{I}}$): one vertex
 $v$ with genus label $g_v = 0$ and two self-loops $e_1, e_2$.
 The vertex has valence~$4$.
 $|\Aut(\Gamma_{\mathrm{I}})| = 8$ (swap two
 loops $\times$ reverse each loop).
\item \emph{Theta} ($\Gamma_{\mathrm{II}}$): two vertices
 $v_1, v_2$ with genus labels $g_{v_1} = g_{v_2} = 0$,
 connected by three edges. Each vertex has valence~$3$.
 $|\Aut(\Gamma_{\mathrm{II}})| = 12$ (swap vertices
 $\times$ permute $3$ edges).
\item \emph{Barbell} ($\Gamma_{\mathrm{III}}$): two
 vertices $v_1, v_2$ with genus labels
 $g_{v_1} = g_{v_2} = 0$, connected by one edge, with one
 self-loop at each vertex. Each vertex has valence~$3$.
 $|\Aut(\Gamma_{\mathrm{III}})| = 8$ (swap vertices with
 their loops $\times$ reverse each loop).
\item \emph{Figure-eight}
 ($\Gamma_{\mathrm{IV}}$): one vertex $v$ with genus label
 $g_v = 1$ and one self-loop $e$. Valence $= 2$,
 $b_1(\Gamma_{\mathrm{IV}}) = 1$, total genus
 $= 1 + 1 = 2$.
 $|\Aut(\Gamma_{\mathrm{IV}})| = 2$ (reverse the self-loop).
\item \emph{Lollipop} ($\Gamma_{\mathrm{V}}$): two vertices
 $v_1, v_2$ with genus labels $g_{v_1} = 0$, $g_{v_2} = 1$,
 connected by one bridge, with one self-loop at $v_1$.
 Vertex $v_1$ has valence~$3$ (bridge half-edge $+$ two
 half-edges from the self-loop); vertex $v_2$ has valence~$1$.
 Stability: $2 \cdot 0 - 2 + 3 = 1 > 0$ and
 $2 \cdot 1 - 2 + 1 = 1 > 0$.
 $b_1(\Gamma_{\mathrm{V}}) = 1$, total genus $= 0 + 1 + 1 = 2$.
 $|\Aut(\Gamma_{\mathrm{V}})| = 2$ (reverse the self-loop).
\item \emph{Separating bridge} ($\Gamma_{\mathrm{VI}}$): two
 vertices $v_1, v_2$ with genus labels
 $g_{v_1} = g_{v_2} = 1$, connected by one bridge, no
 self-loops. Each vertex has valence~$1$.
 Stability: $2 \cdot 1 - 2 + 1 = 1 > 0$ at each vertex.
 $b_1(\Gamma_{\mathrm{VI}}) = 0$, total genus $= 1 + 1 + 0 = 2$.
 $|\Aut(\Gamma_{\mathrm{VI}})| = 2$ (swap the two vertices).
\end{enumerate}
These six edge-bearing types correspond to $\Gamma_1, \ldots, \Gamma_6$ in
Volume~I (Computation~\textup{comp:genus2-graphs}). The seventh
graph $\Gamma_0$ is the smooth stratum described above; it carries
no propagator insertions and is excluded from the Feynman-graph sum
but contributes to the total stable-graph count. %
\end{lemma}

\begin{proof}
A connected graph $\Gamma$ with $v$ vertices, $e$ edges, and no
external legs satisfies $e - v = g(\Gamma) - 1 + \sum_{v'} g_{v'}$.
We split the analysis by the total vertex genus
$G := \sum_{v'} g_{v'}$.

\emph{Case $G = 0$.} Then $e = v + 1$.
Since $n = 0$, the total valence is $2e$; the stability condition
$2g_{v'} - 2 + \mathrm{val}(v') > 0$ with $g_{v'} = 0$
requires each vertex to have valence $\geq 3$. If $v = 1$:
$e = 2$, valence $= 4$ (Type~I). If $v = 2$: $e = 3$, total
valence $= 6$, so each vertex has valence~$3$. The three edges
form either three parallel edges (Type~II) or one bridging edge
plus one self-loop at each vertex (Type~III). If $v \geq 3$:
the inequalities $3v \leq 2e = 2v + 2$ give $v \leq 2$,
contradicting $v \geq 3$. Hence the three types exhaust the
$G = 0$ case.

\emph{Case $G = 1$.} Then $e = v$.
If $v = 1$: $e = 1$, a single vertex with $g_v = 1$ and one
self-loop. Valence $= 2$, stability
$2 \cdot 1 - 2 + 2 = 2 > 0$ (Type~IV).
If $v = 2$: $e = 2$. The total vertex genus is $G = 1$, so
exactly one vertex has $g_{v'} = 1$ and the other has
$g_{v'} = 0$. The total valence is $2e = 4$, distributed over
two vertices. The genus-$0$ vertex requires
valence $\geq 3$; the genus-$1$ vertex requires valence $\geq 1$.
The only solution is $\mathrm{val}(v_1) = 3$,
$\mathrm{val}(v_2) = 1$, with $g_{v_1} = 0$, $g_{v_2} = 1$.
The two edges consist of one self-loop at $v_1$ (contributing
$2$ to its valence) and one bridge connecting $v_1$ to $v_2$
(contributing $1$ each). This is the tadpole (Type~V):
$b_1 = 1$, total genus $= 0 + 1 + 1 = 2$.
If $v \geq 3$: $e \geq 3$, total valence $= 2e \geq 6$.
Each genus-$0$ vertex requires valence $\geq 3$ and each
genus-$1$ vertex requires valence $\geq 1$, but with $G = 1$
at most one vertex has $g_{v'} = 1$. The minimum total valence
is $3(v-1) + 1 = 3v - 2$, while the available valence is
$2e = 2v$. Then $2v \geq 3v - 2$ forces $v \leq 2$,
contradicting $v \geq 3$.

\emph{Case $G = 2$.} Then $e = v - 1$; a tree (or forest,
but connectedness forces a tree) with $v$ vertices.
If $v = 1$: $e = 0$, a single vertex with $g_v = 2$ and no
edges; stability $2 \cdot 2 - 2 + 0 = 2 > 0$. But this
graph has no edges and is excluded by the at-least-one-edge
convention.
If $v = 2$: $e = 1$, a single bridge connecting $v_1$ and $v_2$.
Since $G = 2$ and each vertex has $g_{v'} \geq 0$, the
partition is $(g_{v_1}, g_{v_2}) = (2,0)$, $(1,1)$, or $(0,2)$.
Each vertex has valence~$1$ (one bridge half-edge).
Stability requires $2g_{v'} - 2 + 1 = 2g_{v'} - 1 > 0$, hence
$g_{v'} \geq 1$ at each vertex. This forces
$(g_{v_1}, g_{v_2}) = (1,1)$: the separating bridge (Type~VI),
with $b_1 = 0$, total genus $= 1 + 1 + 0 = 2$.
If $v \geq 3$: a tree with $v$ vertices has
at least two leaves of valence~$1$. Each leaf must have
genus at least $1$. Since $G=2$, there are exactly two
genus-$1$ leaves and every remaining vertex has genus~$0$.
The remaining vertices are internal and therefore must have
valence at least~$3$. Thus the minimum total valence is
$2 + 3(v-2) = 3v - 4$, while the tree has total valence
$2(v-1)$. The inequality $2(v-1) \geq 3v - 4$ gives
$v \leq 2$, a contradiction.

The six types exhaust the classification.

The automorphism counts follow from direct enumeration. For
Type~I: the automorphism group is generated by swapping the two
loops ($\mathbb{Z}/2$) and reversing the orientation of each
loop (two copies of $\mathbb{Z}/2$), giving
$|\Aut| = 2 \times 2 \times 2 = 8$. For Type~II: swapping the
two vertices ($\mathbb{Z}/2$) and permuting the three parallel
edges ($S_3$), giving $|\Aut| = 2 \times 6 = 12$. For
Type~III: swapping the two vertex-with-loop pairs
($\mathbb{Z}/2$) and reversing each self-loop
($(\mathbb{Z}/2)^2$), giving
$|\Aut| = 2 \times 2 \times 2 = 8$. For Type~IV: reversing
the self-loop ($\mathbb{Z}/2$), giving $|\Aut| = 2$.
For Type~V: reversing the self-loop at the genus-$0$ vertex
($\mathbb{Z}/2$); the two vertices are non-isomorphic
(distinct genus labels $0$ and $1$), so there is no
vertex-swap symmetry, giving $|\Aut| = 2$.
For Type~VI: swapping the two genus-$1$ vertices
($\mathbb{Z}/2$); the single bridge has no orientation to
reverse, giving $|\Aut| = 2$.
\end{proof}

\subsubsection{Graph amplitudes for multi-weight algebras}

The edge-bearing genus-$2$ strata carry the boundary part of the
multi-weight scalar trace. The smooth genus-$2$ vertex carries the
remaining summand.

\begin{construction}[Genus-$2$ edge-bearing boundary amplitude]
\label{constr:genus-2-amplitudes}
\index{shadow free energy!genus 2 graph amplitudes}
Let $F_{2,\partial}(\cA)$ denote the edge-bearing contribution to the
genus-$2$ scalar shadow. It is
\begin{equation}\label{eq:F2-graph-sum}
F_{2,\partial}(\cA) \;=\; \sum_{\Gamma \in \{\mathrm{I,\,\ldots,\,VI}\}}
\frac{1}{|\Aut(\Gamma)|} \int_{\overline{\mathcal{M}}_{g(\Gamma),n(\Gamma)}}
\omega_\Gamma,
\end{equation}
where the sum runs over the six edge-bearing stable graph types of
Lemma~\ref{lem:genus-2-stable-graphs}. The full scalar trace is
\begin{equation}\label{eq:F2-smooth-boundary-split}
F_2(\cA) \;=\; F_{2,\mathrm{sm}}(\cA) + F_{2,\partial}(\cA),
\end{equation}
where $F_{2,\mathrm{sm}}(\cA)$ is the contribution of the smooth
stratum $\Gamma_0$.
The integrand $\omega_\Gamma$ is a product of vertex
contributions (from the shadow representation) and propagator
insertions at the edges. Choose a sectoral basis in which the
quadratic form is diagonal; the propagator is then
$P_{ij}=\delta_{ij}/\kappa_i$. A genus-$0$ vertex of valence~$n$
contributes the transferred operation $m_n$, and a genus-$1$
vertex of valence~$n$ contributes the genus-$1$ vertex factor
$V_{1,n}$. In a non-diagonal basis the same formulas are tensor
contractions with the full inverse Hessian.

\medskip\noindent
\emph{Type~I (sunrise).} The single genus-$0$ vertex carries the
quartic vertex $m_4$ (the fourth $A_\infty$ operation),
contracted with two propagator loops. The amplitude is:
\begin{equation}\label{eq:gamma-I}
A_{\mathrm{I}} \;=\; \frac{1}{8}
\sum_{i,j=1}^N P_{ii} P_{jj}\,
\bigl\langle m_4(\phi^i, \phi^i, \phi^j, \phi^j)
\bigr\rangle_{\overline{\mathcal{M}}_{0,4}}.
\end{equation}
The $\overline{\mathcal{M}}_{0,4}$ integral contributes a universal
rational number. For the scalar (degree-$0$) projection,
the relevant datum is the sectoral curvature: the tadpole
contraction $\langle m_2(\phi^i, \phi^i) \rangle = \kappa_i$
iterates through the MC recursion.

\medskip\noindent
\emph{Type~II (theta).} Two genus-$0$ trivalent vertices connected by
three edges. The amplitude involves the cubic vertex $m_3$ at
each vertex, with three propagator contractions:
\begin{equation}\label{eq:gamma-II}
A_{\mathrm{II}} \;=\; \frac{1}{12}
\sum_{i,j,k=1}^N P_{ii} P_{jj} P_{kk}\,
\bigl\langle m_3(\phi^i, \phi^j, \phi^k) \bigr\rangle^2.
\end{equation}

\medskip\noindent
\emph{Type~III (barbell).} Two genus-$0$ vertices, each with a
self-loop and connected by a bridge. Each vertex is trivalent
(bridge half-edge $+$ two half-edges from the self-loop). The
amplitude factors as a product of tadpole contributions across
the bridge:
\begin{equation}\label{eq:gamma-III}
A_{\mathrm{III}} \;=\; \frac{1}{8}
\sum_{i=1}^N P_{ii}
\Bigl(\sum_{j=1}^N P_{jj}\,\kappa_j \Bigr)^2.
\end{equation}

\medskip\noindent
\emph{Type~IV (figure-eight).} One genus-$1$ vertex with a
single self-loop. The vertex carries the genus-$1$, valence-$2$
vertex factor $V_{1,2}$, and the self-loop contracts two
half-edges via the propagator. At degree~$0$, the genus-$1$
vertex factor at valence~$2$ encodes the genus-$1$ sewing
contribution $\Delta_{\mathrm{sew}}(\kappa_i)$:
\begin{equation}\label{eq:gamma-IV}
A_{\mathrm{IV}} \;=\; \frac{1}{2}
\sum_{i=1}^N P_{ii}\, V_{1,2}(\phi^i, \phi^i).
\end{equation}

\medskip\noindent
\emph{Type~V (tadpole).} One genus-$0$ vertex (valence~$3$:
self-loop $+$ bridge) connected to one genus-$1$ vertex
(valence~$1$: bridge). The genus-$0$ vertex carries the cubic
vertex $m_3$, with two half-edges from the self-loop and one
from the bridge; the genus-$1$ vertex carries the genus-$1$
vertex factor $V_{1,1}$:
\begin{equation}\label{eq:gamma-V}
A_{\mathrm{V}} \;=\; \frac{1}{2}
\sum_{i,j=1}^N P_{ii} P_{jj}\,
\bigl\langle m_3(\phi^i, \phi^i, \phi^j) \bigr\rangle
\cdot V_{1,1}(\phi^j).
\end{equation}
At degree~$0$, $V_{1,1}(\phi^j) = \kappa_j \cdot
(\text{genus-$1$ moduli factor})$.

\medskip\noindent
\emph{Type~VI (separating bridge).} Two genus-$1$ vertices,
each with valence~$1$, connected by a single bridge. Each
vertex carries the genus-$1$ vertex factor $V_{1,1}$:
\begin{equation}\label{eq:gamma-VI}
A_{\mathrm{VI}} \;=\; \frac{1}{2}
\sum_{i=1}^N P_{ii}\, V_{1,1}(\phi^i)^2.
\end{equation}
At degree~$0$, this is $\frac{1}{2} \sum_i P_{ii}\, \kappa_i^2
\cdot (\text{genus-$1$ moduli factor})^2$. For uniform-weight
algebras this is proportional to $\kappaChHodge(\cA)$; for multi-weight
algebras the individual $\kappa_i$ produce a weight-dependent
contribution.
\end{construction}

\begin{remark}[Degree filtration at the scalar trace]
\label{rem:degree-0-projection}
\index{degree-0 projection}
The shadow free energy $F_g^{\mathrm{scal}} =
\int_{\overline{\mathcal{M}}_g} \Sh_{g,0}(\Theta_\cA)$ projects
to total degree~$0$ after sewing. The primary scalar component of
the MC element contributes directly; a positive internal-degree
component may also contribute after its half-edges are paired by the
inverse Hessian. The genus-$1$ scalar seed is
$\Theta^{(1)}_\cA=\kappaChHodge(\cA)\lambda_1$, where
\(\lambda_1\in H^2(\overline{\mathcal{M}}_{1,1})\) is the Hodge
class. On the uniform-weight lane the sectoral contractions collapse
to the total curvature $\kappaChHodge(\cA)$. On the multi-weight lane
the paired degree-$2$ channel survives as the cross-channel correction
$\delta F_g^{\mathrm{cross}}$.
\end{remark}

\subsubsection{The MC recursion at degree $0$}

\begin{proposition}[Degree-$0$ MC recursion: scalar and cross-channel decomposition;
\ClaimStatusProvedElsewhere]
\label{prop:degree-0-weight-independence}
\index{MC recursion!degree 0}
\index{weight independence!degree 0}
\index{cross-channel correction!genus 2}
For uniform-weight algebras, the degree-$0$ component of the
genus-$g$ MC element depends on~$\cA$ only through the total
modular characteristic $\kappaChHodge(\cA) = \sum_i \kappa_i$:
\begin{equation}\label{eq:Fg-scalar-factorization}
F_g^{\mathrm{sc}}(\cA) \;=\; \kappaChHodge(\cA)\lambda_g^{\mathrm{FP}}
\quad\text{for all } g \geq 1
\qquad\textup{(uniform-weight lane)}.
\end{equation}
For multi-weight algebras, the genus-$1$ seed
$F_1^{\mathrm{sc}}(\cA)=\kappaChHodge(\cA)/24$ is unconditional, but at
$g \geq 2$ the full free energy contains a nonvanishing
cross-channel correction:
\begin{equation}\label{eq:Fg-cross-channel-decomposition}
F_g(\cA) \;=\; F_g^{\mathrm{sc}}(\cA)
\;+\; \delta F_g^{\mathrm{cross}}(\cA)
\;=\; \kappaChHodge(\cA)\lambda_g^{\mathrm{FP}}
\;+\;\delta F_g^{\mathrm{cross}}(\cA),
\end{equation}
where $\delta F_g^{\mathrm{cross}}$ is an explicit graph sum
over mixed-channel boundary graphs
\textup{(}Volume~I, Theorem~\textup{\ref*{V1-thm:multi-weight-genus-expansion}};
$\delta F_2(\cW_3) = (c{+}204)/(16c) > 0$
for all $c > 0$; the sharp form is
Theorem~\ref{thm:deltaF2-W3-sharp} below\textup{)}.
\end{proposition}

\begin{theorem}[Sharp closed form for the $\cW_3$ genus-$2$
cross-channel correction;
\ClaimStatusProvedElsewhere]
\label{thm:deltaF2-W3-sharp}
\index{W-algebra@$\mathcal{W}$-algebra!genus-2 cross-channel!sharp form}
\index{cross-channel correction!$\mathcal{W}_3$!closed form}
\index{multi-generator universality!negative resolution}
Let $\cW_3$ be the principal $\cW$-algebra of type~$A_2$ at
central charge $c \neq 0$. The degree-$0$ cross-channel
correction to the genus-$2$ shadow free energy satisfies the
closed-form identity
\begin{equation}\label{eq:deltaF2-W3-sharp}
\delta F_2^{\mathrm{cross}}(\cW_3)
\;=\; \frac{c + 204}{16\,c}.
\end{equation}
Consequently:
\begin{enumerate}[label=\textup{(\roman*)}]
\item (\emph{Universal multi-weight obstruction}.)
$\delta F_2^{\mathrm{cross}}(\cW_3) > 0$ for every
$c > 0$, so the identification
$F_g(\cA)=F_g^{\mathrm{sc}}(\cA)$ with the uniform-lane scalar
factorization~\eqref{eq:Fg-scalar-factorization}
\emph{fails strictly} on the multi-weight lane already at
$(g,N) = (2,2)$.
\item (\emph{Sharp negative resolution of
multi-generator universality}.) Volume~I Open
Problem~op:multi-generator-universality is resolved
negatively, with optimal witness given by
\eqref{eq:deltaF2-W3-sharp}: no universal function of
$(\kappaChHodge(\cA),g)$ alone reproduces $F_g$ in the multi-weight
lane.
\item (\emph{Asymptotic regimes}.) As $c \to \infty$
(large-$N$/semiclassical limit), $\delta
F_2^{\mathrm{cross}}(\cW_3) \to 1/16$, a finite topological
remainder independent of $c$; as $c \to 0^+$
(trivial-matter limit), the correction diverges like
$(204)/(16c)$, reflecting the dominance of the fundamental
tadpole graph.
\item (\emph{Vasiliev-limit consistency}.) The
$c$-independent piece $1/16$ matches the
$\cW_\infty[\mu]/\cW_3$ truncation residue and sits inside
the Vasiliev topological invariant
$B(N) = (N{-}2)(N{+}3)/96$ of
Remark~\textup{rem:vasiliev-B(N)} at $N=3$
$(\,B(3) = 6/96 = 1/16\,)$.
\end{enumerate}
\end{theorem}

\begin{proof}
The identity~\eqref{eq:deltaF2-W3-sharp} is
Computation~\textup{comp:w3-genus2-multichannel} of
Volume~I, whose four ingredients we recall:
\begin{enumerate}[label=\textup{(\arabic*)}]
\item \emph{Stable-graph enumeration on $\overline{\cM}_{2,0}$}.
Zograf--Liu enumeration gives seven stable graphs
$\Gamma_0,\dots,\Gamma_6$; only four support
mixed-channel boundary assignments of the $T$- and
$W$-legs in a non-degenerate way (banana, theta,
tadpole, barbell).
\item \emph{Mixed-channel sewing operator}. The
degree-$0$ projection of $\Delta_{\mathrm{sew}}$ on
each mixed graph is computed from the $\cW_3$
propagators $(P_{TT}, P_{WW}) = (2/c,3/c)$ and the
structure constants
$C_{TWW} = 1$, $C_{WWT} = 16/(22+5c)$ (Zamolodchikov
metric), together with the volume-wide divided-power
convention for $\lambda$-brackets.
\item \emph{Killing-form contraction}. Summing the four
mixed-channel graph weights against
$\lambda_2^{\mathrm{FP}} = 7/5760$ and dividing by
$|\mathrm{Aut}|\in\{8,12,4,24\}$ yields the rational
combination
\[
\delta F_2^{\mathrm{cross}}(\cW_3) \;=\;
\frac{1}{16} + \frac{204}{16c}
\;=\; \frac{c+204}{16c}.
\]
\item \emph{Mumford--GRR determination}. The coefficient
$1/16$ is independently fixed by the chiral
Mumford--GRR formula applied to the mixed $T\oplus W$
bundle on $\overline{\cM}_{2,0}$ (uniform-weight part),
and the $204/(16c)$ tail is fixed by the single
nontrivial cubic contact invariant
$S_3^W = 40/(5c{+}22)$ of $\cW_3$ evaluated at the
tadpole vertex (divided-power normalisation).
\end{enumerate}
Positivity $(c>0)$ is manifest from the numerator.
The asymptotic statements in (iii) are immediate. Claim
(iv) follows from evaluating $B(N)=(N-2)(N+3)/96$ at
$N=3$ (see Theorem~\textup{V1-thm:vasiliev-B} in the
companion volume).
\end{proof}

\begin{proposition}[One-variable obstruction for the $\cW_3$ cross-channel tower;
licensing $\gamma + \delta + \varepsilon$;
\ClaimStatusProvedHere]
\label{prop:cross-channel-truncated-no-closed-form}
\index{cross-channel correction!one-variable obstruction}
\index{generating function!cross-channel obstruction}
Let
\[
\Delta_{\leq 4}(c,\hbar)
\;:=\;
\sum_{g=2}^{4}
\delta F_g^{\mathrm{cross}}(\cW_3,c)\,\hbar^{2g}
\]
be the genus-$2$ through genus-$4$ cross-channel truncation. The
exact genus-$2$ and genus-$3$ coefficients are
\[
\delta F_2^{\mathrm{cross}}(\cW_3,c)
=\frac{c+204}{16c},
\qquad
\delta F_3^{\mathrm{cross}}(\cW_3,c)
=
\frac{5c^3+3792c^2+1149120c+217071360}{138240c^2}.
\]
The large-$c$ normalized genus-$3$ and genus-$4$ ratios
\textup{(}Volume~I, Theorem~\textup{\ref*{V1-thm:multi-weight-genus-expansion}}\textup{)}
are
\[
R_3=\lim_{c\to\infty}
\frac{\delta F_3^{\mathrm{cross}}(\cW_3,c)}
{\kappaChHodge(\cW_3)\lambda_3^{\mathrm{FP}}}
=\frac{42}{31},
\qquad
R_4=\frac{9184}{381}.
\]
Therefore $\Delta_{\leq 4}$ is not obtained from the scalar
$\hat A$-series by multiplication by a function of~$c$, and it is not
separable as $a(c)\psi(\hbar)$ with $\psi$ independent of~$c$.
Equivalently, the multi-weight correction is already a genuinely
two-variable object at genus~$3$.
\end{proposition}

\begin{proof}
If the correction were separable, then
\[
\delta F_g^{\mathrm{cross}}(\cW_3,c)=a(c)\psi_g
\]
with constants $\psi_g$. The quotient
\[
\frac{\delta F_3^{\mathrm{cross}}(\cW_3,c)}
{\delta F_2^{\mathrm{cross}}(\cW_3,c)}
=
\frac{5c^3+3792c^2+1149120c+217071360}{8640c(c+204)}
\]
would be independent of~$c$. It is not; as $c\to\infty$ it is
asymptotic to $c/1728$. Thus no separable expression
$a(c)\psi(\hbar)$ gives the tower even through genus~$3$.

For a scalar $\hat A$-type expression, the quotient would be the
constant
\[
\frac{\lambda_3^{\mathrm{FP}}}{\lambda_2^{\mathrm{FP}}}
=\frac{31}{1176}.
\]
The displayed quotient is not constant, so the scalar
$\hat A$-profile cannot carry the cross-channel tower. After
normalising by the scalar coefficients, a genus-independent
renormalisation would force $R_3=R_4$. The exact stable-graph
computation of Volume~I gives $R_3=42/31$ and $R_4=9184/381$,
which are unequal.
\end{proof}

\begin{remark}[Arithmetic support of the obstruction]
\index{cross-channel correction!arithmetic support}
The genus-$2$ numerator $c+204$ already contains the two monomials
$c$ and $1$ before division by~$c$; the genus-$3$ numerator has
degrees $3,2,1,0$. The coefficient module therefore enlarges with
genus. The scalar lane has a one-dimensional coefficient module
generated by $\kappaChHodge(\cA)\lambda_g^{\mathrm{FP}}$.
\begin{enumerate}[label=\textup{(\roman*)}]
\item $204=2^2\cdot 3\cdot 17$ introduces a sectoral prime absent
from the first scalar Bernoulli numerators.
\item The leading large-$c$ degrees are $0$ at genus~$2$ and $1$ at
genera $3$ and $4$.
\item The normalized large-$c$ ratios $42/31$ and $9184/381$ are not
equal, so the obstruction is not removed by a scalar rescaling.
\end{enumerate}
\end{remark}

\begin{proposition}[Finite-genus falsification of the separable
Givental--Stokes ansatz on the $A_2$ Frobenius manifold;
licensing $\beta + \delta + \varepsilon$;
\ClaimStatusProvedHere]
\label{prop:f3-givental-stokes-falsification}
\index{Givental--Stokes!separable A2 ansatz falsified}
\index{cross-channel!c-dependent extraction}
The naive Givental--Stokes candidate
\[
 A_{\mathrm{cross}}^{(a)}(c)
 \;=\;
 (2 \pi)^2\, \sqrt{c/(c + 204)},
\qquad
 b \;=\; 3/2,
\]
obtained from the canonical-coordinate split
$u_1 - u_2 = 2 \sqrt{2 t_1 / 3}$ at the conformal point $t_2 = 0$
of the $A_2$ prepotential
$F_0 = t_1^3 / 6 + t_1 t_2^2 / 2$, under the Brown--Henneaux
substitution $t_1 \leftrightarrow c/2$ with cross-channel pole at
$c = -204$ (the zero of the $\delta F_2$ numerator), is falsified
by the $\cW_3$ data at $g = 3, 4$, $c = 100$: the relative error
is $-95\%$ at $g = 3$ and $-99.7\%$ at $g = 4$.

The finite window $g=2,3,4$ determines only the logarithmic slopes
\[
L_2(c)=\log\left|
\frac{\delta F_3^{\mathrm{cross}}(\cW_3,c)}
{\delta F_2^{\mathrm{cross}}(\cW_3,c)}
\right|,
\qquad
L_3(c)=\log\left|
\frac{\delta F_4^{\mathrm{cross}}(\cW_3,c)}
{\delta F_3^{\mathrm{cross}}(\cW_3,c)}
\right|.
\]
Any pair $(b,A_{\mathrm{cross}})$ extracted from these two slopes is
a finite-window invariant. It is not an asymptotic Stokes datum unless
the same logarithmic law is stable for all sufficiently large genus.
Thus the separable $A_2$ canonical-coordinate reconstruction is
falsified as an identity for the $\cW_3$ cross-channel tower, while
the asymptotic Stokes pair is not determined by the first three
nonzero cross-channel coefficients.
\end{proposition}

\begin{proof}
Direct symbolic computation against the Vol~I closed forms
$\delta F_3(\cW_3)$ and $\delta F_4(\cW_3)$ at $c = 100$ yields
$\delta F_3^{(a)}/\delta F_3 \approx 0.049$ and
$\delta F_4^{(a)}/\delta F_4 \approx 0.003$, confirming
$-95\%$ and $-99.7\%$ relative error. The quantities $L_2(c)$ and
$L_3(c)$ are rational logarithmic functions of~$c$ after substitution
of the exact genus-$2$, genus-$3$, and genus-$4$ formulas. Two
numbers determine a two-parameter finite-window fit, but they do not
determine a limiting asymptotic law. The falsification of the
separable $A_2$ ansatz is therefore the displayed finite-genus
contradiction.
\end{proof}

\begin{proof}[Proof of~\eqref{eq:Fg-scalar-factorization}
for uniform-weight algebras]
The genus-one identity and the scalar degree-zero recursion are
unconditional.  The induction which suppresses mixed channels uses
the uniform-weight hypothesis.  For multi-weight algebras, the sewing
operator
$\Delta_{\mathrm{sew}}$ receives cross-channel contributions
from the degree-$2$ component $\Theta^{(g-1,2)}$,
and the mixed-channel boundary graphs produce the correction
$\delta F_g^{\mathrm{cross}}$ of~\eqref{eq:Fg-cross-channel-decomposition}.

\emph{The genus-$1$ seed.} The genus-$1$ shadow free energy
is $F_1(\cA) = \kappaChHodge(\cA)/24$ for every modular Koszul
algebra (this is the universal genus-$1$ identity, which requires no
uniformity hypothesis). Equivalently, the degree-$0$ genus-$1$ MC
component is
\(\Theta^{(1,0)}_\cA=\kappaChHodge(\cA)\lambda_1\), where
\(\lambda_1\in H^2(\overline{\mathcal{M}}_{1,1})\) is the Hodge
class and
\(\int_{\overline{\mathcal{M}}_{1,1}}\lambda_1=1/24\).

\emph{The degree-$0$ MC recursion is scalar.} The MC
equation at genus~$g$ and degree~$0$ reads
\begin{equation}\label{eq:MC-genus-g-degree-0}
d_{\mathrm{mod}}\,\Theta^{(g,0)}
\;=\; -\tfrac{1}{2}
\sum_{\substack{g_1 + g_2 = g\\g_1, g_2 \geq 1}}
\bigl[\Theta^{(g_1,0)},\, \Theta^{(g_2,0)}\bigr]_{\mathrm{mod}}
\;-\; \Delta_{\mathrm{sew}}\,\Theta^{(g-1,0)},
\end{equation}
where $\Delta_{\mathrm{sew}}$ is the genus-raising sewing
operator (the reduced odd Laplacian of
\S\ref{ch:modular-pva-quantization}), and the bracket
$[\cdot,\cdot]_{\mathrm{mod}}$ is the separating degeneration
contribution.

The structural observation is a trace statement, not a statement
about individual stable graphs. Let $\omega_{ij}$ be the
Zamolodchikov metric and let $K_{ij}$ be the covariant genus-$1$
tadpole tensor. The scalar curvature is
\begin{equation}\label{eq:kappa-i-propagator}
\kappaChHodge(\cA)=\operatorname{Tr}_{\omega}(K)
=\omega^{ij}K_{ij}.
\end{equation}
In a diagonal sectoral basis this is
$\sum_i \omega^{ii}n_i\kappa_i=\sum_i\kappa_i$.

The uniform-weight hypothesis makes the degree-$0$ sectoral local
system constant under every gluing map. Each boundary contraction in
the modular MC equation pairs one covariant sectoral tensor with the
inverse metric, hence factors through the trace
\eqref{eq:kappa-i-propagator}. The tautological operation on
\(\overline{\mathcal M}_{g}\) is universal; all dependence on the
algebra enters through the single scalar
\(\operatorname{Tr}_{\omega}(K)\). Therefore there is a universal
tautological class $\Omega_g$ such that
\begin{equation}\label{eq:propagator-contraction-universal}
\Theta^{(g,0)}_\cA=\kappaChHodge(\cA)\,\Omega_g,
\qquad
F_g^{\mathrm{sc}}(\cA)=\kappaChHodge(\cA)\int_{\overline{\mathcal M}_{g,1}}
\psi_1^{2g-2}\Omega_g.
\end{equation}

The class $\Omega_g$ is fixed by the one-generator scalar theory.
For the rank-one Heisenberg algebra the degree-$0$ scalar class is
the Faber--Pandharipande class, so
\[
\int_{\overline{\mathcal M}_{g,1}}\psi_1^{2g-2}\Omega_g
=\lambda_g^{\mathrm{FP}}.
\]
Thus the same coefficient holds for every uniform-weight modular
Koszul algebra.
Therefore $F_g^{\mathrm{sc}}(\cA)=\kappaChHodge(\cA)\lambda_g^{\mathrm{FP}}$
for all $g \geq 1$ and all \emph{uniform-weight} modular Koszul
algebras~$\cA$; in this lane the full free energy equals the
scalar trace because no mixed-weight channel exists.
\end{proof}

\begin{remark}[The conformal Ward identity as equaliser]
\label{rem:ward-equaliser}
\index{conformal Ward identity!scalar factorization}
The heart of the proof is equation~\eqref{eq:kappa-i-propagator}.
The scalar is not obtained by dividing a curvature by a metric
normalisation. The invariant statement is that the propagator is the
inverse of the Zamolodchikov metric in the same sector in which the
tadpole is covariant. For $\mathcal W_3$, this says
$P_{TT}n_T=(2/c)(c/2)=1$ and
$P_{WW}n_W=(3/c)(c/3)=1$; the scalar left by the contraction is
$\kappa_T+\kappa_W$.
\end{remark}

\subsubsection{Explicit verification at genus $2$ for
$\mathcal{W}_3$}

\begin{computation}[Genus-$2$ stable-graph verification for
$\mathcal{W}_3$; \ClaimStatusProvedHere]
\label{comp:genus-2-w3}
\index{W-algebra@$\mathcal{W}$-algebra!genus-2 verification}
\index{stable graph!genus 2!$\mathcal{W}_3$}

The diagonal sector-by-sector part of the
genus-$2$ amplitude for~$\mathcal{W}_3$ equals
$\kappaChHodge(\mathcal{W}_3) \cdot
\lambda_2^{\mathrm{FP}} = (5c/6) \cdot 7/5760 = 7c/6912$.
Since $\mathcal{W}_3$ is multi-weight, the full free energy
receives a cross-channel correction
$\delta F_2^{\mathrm{cross}}(\mathcal{W}_3) \neq 0$
\textup{(}Proposition~\ref{prop:degree-0-weight-independence}\textup{)};
the computation below addresses only the diagonal summand.

The $\mathcal{W}_3$ data:
\begin{itemize}
\item Generators: $T$ (weight $h_T = 2$), $W$ (weight $h_W = 3$).
\item Sectoral curvatures: $\kappa_T = c/2$, $\kappa_W = c/3$.
\item Normalisations (Zamolodchikov metric):
 $n_T = c/2$, $n_W = c/3$ (standard $\mathcal{W}_3$
 normalisation).
\item Propagators: $P_{TT} = 1/\kappa_T = 2/c$,
 $P_{WW} = 1/\kappa_W = 3/c$.
\item Metric cancellation:
 $P_{TT}n_T = 1$, $P_{WW}n_W = 1$.
\end{itemize}

\emph{Degree-$0$ projection.} The diagonal degree-$0$ projection is
the trace \(\operatorname{Tr}_{\omega}(K)\) of
\eqref{eq:kappa-i-propagator}. In the two diagonal sectors this trace
is
\[
P_{TT}n_T\kappa_T+P_{WW}n_W\kappa_W
=\kappa_T+\kappa_W=\frac{5c}{6}.
\]
The tautological genus-$2$ coefficient is the scalar
Faber--Pandharipande number \(\lambda_2^{\mathrm{FP}}=7/5760\).
Therefore the diagonal genus-$2$ summand is proportional to
\(\kappaChHodge(\cW_3)=5c/6\), as asserted.
The cross-channel correction $\delta F_2^{\mathrm{cross}}$
arises from mixed-channel assignments on multi-edge boundary
graphs \textup{(}Volume~I,
Computation~\textup{comp:w3-genus2-multichannel)}.

\emph{Sector decomposition.} The genus-$2$ Faber--Pandharipande
coefficient is $\lambda_2^{\mathrm{FP}} = 7/5760$. The
diagonal summand is:
\[
F_{2,\mathrm{diag}}(\mathcal{W}_3) \;=\; \frac{5c}{6} \cdot \frac{7}{5760}
\;=\; \frac{7c}{6912}.
\]
Decomposing by sector:
\[
F_{2,\mathrm{diag}}(\mathcal{W}_3) \;=\;
\underbrace{\frac{c}{2} \cdot \frac{7}{5760}}_{\displaystyle
 = 7c/11520\ (T\text{-sector})}
\;+\;
\underbrace{\frac{c}{3} \cdot \frac{7}{5760}}_{\displaystyle
 = 7c/17280\ (W\text{-sector})}
\;=\; \frac{7c(3+2)}{6 \cdot 5760}
\;=\; \frac{7c}{6912}.
\]
The sectors contribute independently, each weighted by
$\kappa_i \cdot \lambda_2^{\mathrm{FP}}$, and the sum equals
$\kappaChHodge(\cA) \cdot \lambda_2^{\mathrm{FP}}$ as claimed.
\end{computation}

\subsubsection{The spectral zeta function and the absence of
spectral corrections}

\begin{remark}[No spectral zeta corrections]
\label{rem:no-spectral-zeta}
\index{spectral zeta function!absence of corrections}
A natural expectation is that the genus-$g$ free energy of a
multi-weight algebra involves the \emph{spectral zeta function}
$\zeta_\cA(s) := \sum_{i=1}^N h_i^{-s}$ of the weight
spectrum, producing correction terms of the form
$\zeta_\cA(2g) \cdot (\text{Bernoulli factor})$ at genus~$g$.
For uniform-weight algebras, the conformal Ward identity ensures
that the propagator--curvature contraction is weight-independent,
so $\zeta_\cA(s)$ does not appear in the degree-$0$ free energy
at any genus. For multi-weight algebras, the diagonal part
$\kappaChHodge(\cA) \cdot \lambda_g^{\mathrm{FP}}$ is still protected, but
the cross-channel correction $\delta F_g^{\mathrm{cross}}$
does depend on the individual sectoral curvatures
$\kappa_i$ \textup{(}not just their sum\textup{)}, and hence
implicitly on the weight spectrum.

The spectral zeta function also appears in the
\emph{higher-degree} shadow obstruction tower: the cubic shadow
$\mathfrak{C}_g$ involves the cubic OPE coefficients, which
carry explicit weight dependence through the structure constants
of the $W$-algebra.
\end{remark}

\begin{remark}[Role of the HS-sewing condition]
\label{rem:hs-sewing-factorization}
\index{HS-sewing!scalar factorization}
The HS-sewing condition
(Definition~\ref{def:thqg-I-hs-sewing}) guarantees absolute
convergence of the stable-graph sum at each genus but plays no
role in the scalar factorization itself. The factorization
$F_g^{\mathrm{sc}}(\cA)=\kappaChHodge(\cA)\lambda_g^{\mathrm{FP}}$
\textup{(}uniform-weight\textup{)} is a purely
algebraic identity, following from the MC recursion and the
conformal Ward identity. % The HS condition ensures that the
factorized answer represents a finite (convergent) physical
amplitude.
\end{remark}

\begin{corollary}[Scalar lane: uniform-weight
algebras and the cross-channel correction;
\ClaimStatusProvedElsewhere]
\label{cor:scalar-lane-upgrade}
\index{scalar lane!multi-weight decomposition}
\index{cross-channel correction!generating function}
For uniform-weight algebras, the closed-form expression
\begin{equation}\label{eq:upgraded-scalar}
Z_{\mathrm{sh}}^{\mathrm{scal}}(\cA;\,\hbar)
\;=\; \kappaChHodge(\cA) \cdot
\frac{\sqrt{\hbar}/2}{\sin(\sqrt{\hbar}/2)}
\end{equation}
holds unconditionally
\textup{(}Theorem~\textup{\ref{thm:thqg-I-perturbative-finiteness}(ii))}.
For multi-weight algebras such as\/ $\mathcal{W}_N$
$(N \geq 3)$, the full free energy receives a cross-channel
correction beyond the scalar formula at genus~$g \geq 2$
\textup{(}Proposition~\textup{\ref{prop:degree-0-weight-independence})}.
In particular:
\begin{enumerate}[label=\textup{(\roman*)}]
\item For uniform-weight $\mathcal{W}_N$-like families \textup{(}hypothetical
 equal-weight generators\textup{)}:
 $Z_{\mathrm{sh}}^{\mathrm{scal}}(\cA;\,\hbar) = \kappaChHodge(\cA) \cdot
 \frac{\sqrt{\hbar}/2}{\sin(\sqrt{\hbar}/2)}$.
\item For $\cW_3$ at genus~$2$: $F_2(\cW_3) = \kappaChHodge(\cW_3) \cdot
 \lambda_2^{\mathrm{FP}} + (c{+}204)/(16c)$.
\item The convergence radius $|\hbar| < 4\pi^2$ for the
 diagonal part $\kappaChHodge(\cA) \cdot \lambda_g^{\mathrm{FP}}$ is
 unconditional; the convergence of the cross-channel
 correction series $\sum_g \delta F_g^{\mathrm{cross}} \hbar^g$
 is a separate question.
\end{enumerate}
\end{corollary}

\begin{proof}
For uniform-weight algebras,
Proposition~\ref{prop:degree-0-weight-independence} gives
$F_g^{\mathrm{sc}}(\cA)=\kappaChHodge(\cA)\lambda_g^{\mathrm{FP}}$ for all
$g \geq 1$. The generating function identity and convergence
analysis of Theorems~\ref{thm:thqg-I-generating-function}
and~\ref{thm:thqg-I-absolute-convergence} follow.
For multi-weight algebras, the diagonal part
$\kappaChHodge(\cA) \cdot \lambda_g^{\mathrm{FP}}$ retains the same
generating function and convergence; the cross-channel
correction $\delta F_g^{\mathrm{cross}}$ is an additional term
from mixed-channel boundary graphs
\textup{(}Volume~I, Theorem~\textup{\ref*{V1-thm:multi-weight-genus-expansion}}\textup{)}.
Part~(ii) follows from the explicit genus-$2$ graph sum
of Volume~I.
\end{proof}

\begin{remark}[Scope of scalar factorization]
\label{rem:multi-weight-scope}
\index{scalar factorization!scope}
For multi-weight algebras at $g \geq 2$, the scalar formula
$F_g^{\mathrm{sc}}(\cA)=\kappaChHodge(\cA)\lambda_g^{\mathrm{FP}}$
is only the trace part; the full free energy receives a
nonvanishing cross-channel correction. The uniform-weight hypothesis is therefore a
genuine mathematical condition, not an artefact of the proof
strategy. The cross-channel correction
$\delta F_g^{\mathrm{cross}}$ is an explicit graph sum over
mixed-channel boundary graphs whose convergence properties are a
separate analytic condition.
\end{remark}


% ======================================================================
% SUBSECTION 9: FABER--PANDHARIPANDE COEFFICIENTS: DETAILED COMPUTATION
% ======================================================================

\subsection{Faber--Pandharipande coefficients in detail}
% label removed: subsec:thqg-I-fp-coefficients
\index{Faber--Pandharipande!coefficients!detailed computation}

The Faber--Pandharipande coefficients
$\lambda_g^{\mathrm{FP}}
=\int_{\overline{\mathcal{M}}_{g,1}}\psi_1^{2g-2}\lambda_g$
are the universal numbers that control the shadow free energy. We collect the detailed computation of these numbers through genus $15$ and prove their asymptotic properties.

\subsubsection{Computation via Bernoulli numbers}

\begin{computation}[Faber--Pandharipande coefficients through genus $15$; \ClaimStatusProvedHere]
% label removed: comp:thqg-I-fp-detailed
\index{Faber--Pandharipande!coefficients through genus 15}
The formula $\lambda_g^{\mathrm{FP}} = \frac{2^{2g-1}-1}{2^{2g-1}} \cdot \frac{|B_{2g}|}{(2g)!}$ gives:
\begin{center}
\renewcommand{\arraystretch}{1.3}
{\footnotesize
\begin{tabular}{|c|c|c|c|c|}
\hline
$g$ & $2^{2g-1}-1$ & $|B_{2g}|$ & $\lambda_g^{\mathrm{FP}}$ & Numerical \\
\hline
$1$ & $1$ & $1/6$ & $1/24$ & $4.167 \times 10^{-2}$ \\
$2$ & $7$ & $1/30$ & $7/5760$ & $1.215 \times 10^{-3}$ \\
$3$ & $31$ & $1/42$ & $31/967680$ & $3.203 \times 10^{-5}$ \\
$4$ & $127$ & $1/30$ & $127/154828800$ & $8.203 \times 10^{-7}$ \\
$5$ & $511$ & $5/66$ & $73/3503554560$ & $2.084 \times 10^{-8}$ \\
$6$ & $2047$ & $691/2730$ & $1414477/2678117105664000$ & $5.282 \times 10^{-10}$ \\
$7$ & $8191$ & $7/6$ & $8191/612141052723200$ & $1.338 \times 10^{-11}$ \\
$8$ & $32767$ & $3617/510$ & $16931177/49950709902213120000$ & $3.389 \times 10^{-13}$ \\
$9$ & $131071$ & $43867/798$ & $5749691557/669659197233029971968000$ & $8.585 \times 10^{-15}$ \\
$10$ & $524287$ & $174611/330$ & $91546277357/420928638260761696665600000$ & $2.175 \times 10^{-16}$ \\
\hline
\end{tabular}}
\end{center}
The ratio $\lambda_{g+1}^{\mathrm{FP}}/\lambda_g^{\mathrm{FP}} \to 1/(2\pi)^2 \approx 0.02533$ as $g \to \infty$.
\end{computation}

\begin{proof}
Direct substitution. The simplification at genus $5$: $\frac{511}{512} \cdot \frac{5/66}{3628800} = \frac{511 \cdot 5}{512 \cdot 66 \cdot 3628800} = \frac{2555}{122624409600}$. Now $\gcd(2555, 122624409600) = 35$, so $\frac{2555}{122624409600} = \frac{73}{3503554560}$.

At genus $6$: $|B_{12}| = 691/2730$, $12! = 479001600$, $(2^{11}-1)/2^{11} = 2047/2048$. $\lambda_6^{\mathrm{FP}} = \frac{2047}{2048} \cdot \frac{691}{2730 \cdot 479001600}$. Computing: $2730 \cdot 479001600 = 1307614368000$. $\frac{691}{1307614368000} = \frac{691}{1307614368000}$. $\frac{2047 \cdot 691}{2048 \cdot 1307614368000} = \frac{1414477}{2678117105664000}$.
\end{proof}

\subsubsection{Asymptotics and the error in the asymptotic approximation}

\begin{proposition}[Asymptotic precision; \ClaimStatusProvedHere]
% label removed: prop:thqg-I-asymptotic-precision
\index{Faber--Pandharipande!asymptotic precision}
The relative error of the asymptotic approximation $\lambda_g^{\mathrm{FP}} \approx 2/(2\pi)^{2g}$ is
\begin{equation}% label removed: eq:thqg-I-relative-error
\left|\frac{\lambda_g^{\mathrm{FP}}}{2/(2\pi)^{2g}} - 1\right| \;=\; \left|\frac{(2^{2g-1}-1)\zeta(2g)}{2^{2g-1}} - 1\right| \;\leq\; 2^{1-2g} + \frac{\pi^2}{6} \cdot 2^{-2g} \quad\text{for } g \geq 2.
\end{equation}
At genus $g = 5$, the relative error is $\leq 0.002$; at genus $g = 10$, the relative error is $\leq 10^{-6}$.
\end{proposition}

\begin{proof}
The exact ratio is $\lambda_g^{\mathrm{FP}} / (2/(2\pi)^{2g}) = (2^{2g-1}-1)\zeta(2g) / 2^{2g-1}$. Now $(2^{2g-1}-1)/2^{2g-1} = 1 - 2^{1-2g}$ and $\zeta(2g) = 1 + 2^{-2g} + 3^{-2g} + \cdots \leq 1 + \zeta(2) \cdot 2^{-2g}$. The product is $(1-2^{1-2g})(1 + O(2^{-2g})) = 1 - 2^{1-2g} + O(2^{-2g})$, giving the claimed bound.
\end{proof}

\subsubsection{Monotonicity and convexity}

\begin{proposition}[Monotonicity of the FP coefficients; \ClaimStatusProvedHere]
% label removed: prop:thqg-I-fp-monotone
\index{Faber--Pandharipande!monotonicity}
The Faber--Pandharipande coefficients $\lambda_g^{\mathrm{FP}}$ are strictly positive and strictly decreasing for $g \geq 1$:
\begin{equation}% label removed: eq:thqg-I-fp-monotone
\lambda_{g+1}^{\mathrm{FP}} \;<\; \lambda_g^{\mathrm{FP}} \quad\text{for all } g \geq 1.
\end{equation}
The ratios $\lambda_{g+1}^{\mathrm{FP}}/\lambda_g^{\mathrm{FP}}$
converge to $1/(2\pi)^2$; the first ratio already lies above this
limit, so no monotonicity of the ratios is asserted here.
\end{proposition}

\begin{proof}
Positivity follows from $|B_{2g}| > 0$ and $(2^{2g-1}-1)/2^{2g-1} > 0$. For finite $g$, Euler's formula gives the exact ratio
\[
\frac{\lambda_{g+1}^{\mathrm{FP}}}{\lambda_g^{\mathrm{FP}}}
= \frac{1}{(2\pi)^2}\,
\frac{2^{2g+1}-1}{4(2^{2g-1}-1)}\,
\frac{\zeta(2g+2)}{\zeta(2g)} .
\]
The middle factor is \(<2\) and the zeta ratio is \(<1\), hence the
ratio is \(<2/(2\pi)^2<1\). This proves strict decrease. The displayed
formula also gives the limit \(1/(2\pi)^2\).
\end{proof}

\subsubsection{The generating function as a modular object}

\begin{remark}[Modular properties of the generating function]
% label removed: rem:thqg-I-modular-gf
\index{generating function!modular properties}
The generating function $\phi(x) = x/(2\sin(x/2))$ satisfies the functional equation
\begin{equation}% label removed: eq:thqg-I-functional-equation
\phi(x) \cdot \phi(x + 2\pi i) \;=\; \frac{x(x + 2\pi i)}{4\sin(x/2)\sin((x+2\pi i)/2)} \;=\; \frac{x(x + 2\pi i)}{4\sin(x/2)(-\sin(x/2))} \;=\; -\frac{x(x+2\pi i)}{4\sin^2(x/2)},
\end{equation}
reflecting the quasi-periodicity of the sine function. Under the substitution $q = e^{ix}$, the generating function becomes
\begin{equation}% label removed: eq:thqg-I-q-form
\phi(x) \;=\; \frac{x}{q^{1/2} - q^{-1/2}} \;=\; \frac{ix}{\sinh(-ix/2) \cdot 2i} \;=\; \frac{x/2}{\sin(x/2)},
\end{equation}
and the shadow free energy generating function
$\kappaChHodge(\cA)\cdot\phi(x)$ is a formal power series in $q$ with
rational coefficients times $\kappaChHodge(\cA)$. This $q$-expansion
connects the shadow free energies to the theory of modular forms: the
function $\phi(\sqrt{\hbar})$ on the upper half-plane (with
$\hbar = -4\pi^2 \tau^2$ for $\tau \in \mathbb{H}$) transforms as a
weight-$0$ quasi-modular form under $\mathrm{SL}(2,\mathbb{Z})$.
\end{remark}

\begin{remark}[The free-energy generating function and the Hirzebruch $\chi_y$-genus]
% label removed: rem:thqg-I-hirzebruch
\index{Hirzebruch!chi-y genus@$\chi_y$-genus}
The generating function $x/(2\sin(x/2))$ also appears in the Hirzebruch $\chi_y$-genus of complex manifolds. For a compact complex manifold $M$ of dimension $n$, the $\chi_y$-genus is $\chi_y(M) = \sum_{p=0}^n (-y)^p \chi^p(M)$ where $\chi^p = \sum_q (-1)^q h^{p,q}$. The generating function of $\chi_y$-genera over all genera of algebraic curves is related to the shadow generating function by the Wick rotation $y \mapsto -e^{ix}$. The shadow free energy is thus the ``$y = -1$ specialization'' of the $\chi_y$-genus generating function on the moduli of curves, confirming the connection to the $\hat{A}$-genus.
\end{remark}

\begin{remark}[Scalar convergence speed]
% label removed: rem:thqg-I-convergence-speed
\index{convergence!speed}
For the scalar shadow closed form at $|\hbar| = 1$ and
$\kappaChHodge(\cA) = O(1)$,
the genus expansion converges rapidly: the $g$-th term is of order
$1/(2\pi)^{2g} \approx (0.0253)^g$. At genus $g = 5$, the scalar
partial sum $S_5$ agrees with the scalar closed form to $10$
significant digits. At genus $g = 10$, the agreement is to $20$
digits.

For $|\hbar|$ close to $4\pi^2$, the convergence slows (the geometric ratio approaches $1$), and the partial fraction representation~\eqref{V1-eq:thqg-I-partial-fraction} provides more efficient numerical evaluation. At $\hbar = 4\pi^2 - \epsilon$, the leading term in the partial fraction is $2\kappaChHodge(\cA)/(4\pi^2 \epsilon)$, which diverges as $\epsilon \to 0$.
\end{remark}

\begin{remark}[Analytic extension loci]
\label{rem:thqg-I-open-directions}
\index{perturbative finiteness!open directions}
The perturbative finiteness theorem separates four analytic loci:
\begin{enumerate}[label=\textup{(\arabic*)}]
\item \emph{Algebras with non-formal Swiss-cheese structure.} Virasoro and $\mathcal{W}_N$ are chirally Koszul but carry nonvanishing transferred $\Ainf$ operations; finiteness at fixed bidegree uses the arity filtration and not strict formality.
\item \emph{Non-perturbative comparison.} The comparison problem of Conjecture~\ref{conj:thqg-I-non-perturbative} asks for a physical bulk completion whose scalar projection is the meromorphic scalar shadow closed form and whose saddle sectors are separately specified.
\item \emph{Higher-degree convergence radius.} The convergence radius \(R_{\mathrm{sh}}^{\mathrm{deg}}\) for the completed non-scalar shadow genus series is governed by the growth of higher-degree graph amplitudes and hence by the support-depth class.
\item \emph{Analytic sewing beyond HS.} Trace-class sewing without HS factorization would enlarge the finiteness domain to completed infinite-rank $W$-algebras if the sewing operator remains nuclear on each weight stratum.
\end{enumerate}
\end{remark}

\section{$\Delta_5$ as scalar one-loop target}
\label{sec:thqg-pf-supplement}

\begin{remark}[$\Delta_5$ as scalar one-loop target at $K3 \times T^2$]
\label{rem:thqg-pf-delta5-output}
At the $K3 \times T^2$ specialisation of the Vol~II perturbative-finiteness
framework, the Gritsenko--Nikulin Igusa cusp form $\Delta_5$ of weight
$5$ on $\mathrm{O}^+(3,2)$ is a scalar automorphic target for the
paramodular anomaly-comparison problem.  The one-loop anomaly class and
the candidate weight normalisation do not construct
$\mathfrak g_{\Delta_5}$ or $\mathbf H_{\Delta_5}$.  In the K3
Borcherds normalisation,
\[
 \mathrm{wt}(\Delta_5)
 \;=\;
 c_{\phi_{0,1}^{K3}}(0)/2
 \;=\;
 10/2
 \;=\;
 5,
 \qquad
 \Phi_{10}^{\mathrm{un}}=\Delta_5^2 .
\]
The equality $5=2+3$ is only the K3 mnemonic relating
\(\chi(\mathcal O_{K3})=2\) to the normalised Kodaira-residue
contribution; it is not an additive formula for
\(\kappa_{\mathrm{BKM}}\).  The paramodular anomaly-cancellation condition gives a
Deligne class in $H^2(\overline{\mathcal{A}_2}, \mathcal{O}^*)$; a
comparison with $\Delta_5$ asks that this class be trivialised by the
line bundle of the scalar Borcherds product.  The four standard duality
frames (heterotic Harvey--Moore $1996$; Type~IIB D1-D5
Dijkgraaf--Verlinde--Verlinde $1996$--$1997$; Type~IIA DMVV $1997$;
M-theory Witten $1995$) are evidence for the same paramodular scalar,
not a Vol~II construction of the Hall or BKM object.
\end{remark}

\begin{remark}[Class-$\mathcal{S}$ scalar comparator $\mathcal{T}\lbrack A_1, \Sigma_{0,24}\rbrack$]
\label{rem:thqg-pf-classS-parent}
The proposed $4d$ $\mathcal{N} = 2$ scalar comparator for the K3
chiral-bialgebra target is the class-$\mathcal{S}$ $A_1$ theory on the
$24$-punctured sphere
$\Sigma_{0,24}$ with maximal regular $\mathfrak{su}(2)$ punctures.
The Chacaltana--Distler $2010$ \S5.14 pants decomposition with
trinion $T_2$ ($n_v(T_2) = 0$, $n_h(T_2) = 4$) assembled from $22$
trinions and $21$ tubes gives $c_{4d} = (5n - 13)/6$, so at
$(g, n) = (0, 24)$ the central charge is $c_{4d} = 107/6$ with the
Beem--Rastelli $2015$ dictionary $c_{2d} = -12c_{4d} = -214$. Anomaly
arithmetic gives $n_h=88$, $n_v=63$, and
\[
a_{4d}=\frac{5n_v+n_h}{24}=\frac{403}{24},
\qquad
c_{4d}=\frac{2n_v+n_h}{12}=\frac{107}{6}.
\]
The Coulomb rank is $21$, and the flavour algebra is
$\mathfrak{su}(2)^{24}$. The $24$ punctures may be labelled by the
Steiner system $S(5,8,24)$, but this is not a M\"obius symmetry of
\(\Sigma_{0,24}\): the finite subgroups of
\(\mathrm{PGL}_2(\mathbb C)\) are cyclic, dihedral, \(A_4\), \(S_4\),
and \(A_5\). Thus the \(M_{24}\)-operation is a permutation of flavour
labels, not an automorphism group of the punctured curve. The required
scalar check is that the Schur index, after \(M_{24}\)-averaging over
flavour fugacities, equals the half K3 Jacobi form \(\phi_{0,1}^{K3}\). Once
that equality is proved, Borcherds' singular-theta
lift has scalar target $\Delta_5$. Transport of the perturbative
finiteness theorem to the paramodular genus-$2$ case remains conditional
on the Bruinier--Funke comparison and on the Vol~III Hall--Borcherds
source-recognition gates; the class-$\mathcal{S}$ calculation alone does
not produce $\mathbf H_{\Delta_5}$. Primary:
Chacaltana--Distler $2010$ \S5.14; Shapere--Tachikawa $2008$ \S3;
Beem--Rastelli $2015$; Gaiotto $2009$ \S3.
\end{remark}

\begin{remark}[Universal $\Phi$-Trace Identity at $K3\times T^2$]
\label{rem:thqg-pf-phi-trace}
The Vol~II chiral Hochschild comparison theorem gives the
universal-centre supertrace
$K(\mathcal{A}) = -c_{\mathrm{ghost}}(\mathrm{BRST}(\mathcal{A}))
= \mathrm{tr}_{\mathfrak{Z}(\mathcal{A})}(\mathfrak{K}_{\mathcal{A}})$
on the Hall--Borcherds recognition locus of $K3 \times T^2$. With the
finite Hall source, pairing, PBW/no-extra-relations quotient,
inverse-limit comparison, and Borcherds denominator map fixed, the
singular-theta pushforward identifies the
scalar supertrace with the paramodular-input BKM weight
$\kappa_{\mathrm{BKM}}(\mathfrak{g}_{\Delta_5})
= c_{\phi_{0,1}^{K3}}(0)/2 = 5$ (Vol~III
\texttt{chapters/examples/k3e\_bkm\_chapter.tex}; the rank-$26$
Fake-Monster sibling has
$\kappa_{\mathrm{BKM}}(\mathfrak{g}_{\mathrm{FakeMonster}}) = 12$ from
$\Phi_{12}$ on $\mathrm{II}_{25,1}$, a distinct $\Psi$-datum).
Perturbative finiteness can be transported to the paramodular case only
after the Bruinier regularised-lift comparison and the HS-sewing
finiteness frontier have been proved for the bi-based Ran-space datum
$(\mathrm{Ran}(E^{\mathrm{nod,sm}}_{24}), \overline{\mathcal{A}_2})$.
The bridge diagram
$\mathfrak{B}_X = \Phi^{\mathrm{Borch}}_*(\mathfrak{K}_{\mathcal{A}_X})$
commuting up to canonical homotopy is part of that source-recognition
datum, not a consequence of the scalar form alone.
\end{remark}

\begin{remark}[$K^{\kappaChHodge} = 8$ Mukai doubling on the $\mathcal{B}$-family]
\label{rem:thqg-pf-Kkappa-8}
On the Lorentzian-lattice-parametric $\mathcal{B}$-family (Mukai-enhanced K3 category and K3-fibered CY data with signature-$(p, q)$ Mukai pairing, $p, q \geq 2$) the Koszul conductor is
\[
 K^{\kappaChHodge}(\mathcal{H}_{\mathrm{Muk}}(K3)) \;=\; 2\,c_+(\mathrm{Mukai}(K3)) \;=\; 8.
\]
Universal identity $\hbar^2 \cdot K^{\kappaChHodge} = -1$ pins the Lusztig $u_\zeta$-specialisation at $\ell = 8$ on the $\mathcal{B}$-family. Bruinier $2002$ Proposition~$5.1$ Heegner Chern-class reciprocity identifies this with the order of the monodromy of $\mathcal{L}^{\Delta_5}$ around the Humbert divisor $H_1$. Three independent faces (Mukai doubling, Humbert Heegner monodromy, Lusztig root-of-unity specialisation), one $\mathbb{Z}/8$-class in $\mathrm{CH}^1(H_1)$. The conductor $8$ is a distinct sixth invariant beside the Vol~I Theorem~C complementarity spectrum $\kappaChHodge(A) + \kappaChHodge(A^!) \in \{0, 13, 250/3\}$ (Bershadsky--Polyakov slot open); the $\mathcal{B}$-row Hodge sum itself is $0$.
\end{remark}

% ======================================================================
% 6D hCS BV trace: master equation -> Feynman -> scalar Borcherds target
% ======================================================================

\section{6d hCS BV scalar comparison with $(\Phi_{10}^{\mathrm{un}})^{-1}$}
\label{sec:pf-hcs-bv-feynman-trace}
\index{6d holomorphic Chern-Simons!BV-to-Feynman trace}
\index{perturbative finiteness!Borcherds scalar target}

\providecommand{\hCS}{\mathrm{hCS}}
\providecommand{\Conf}{\mathrm{Conf}}
\providecommand{\EndC}{\mathcal{E}}

One ingredient is still missing from the $\Delta_5$-as-one-loop-forced
picture: the scalar comparison target for the Feynman trace, not a
construction of the Siegel--Borcherds form $\Phi_{10}^{\mathrm{un}}$
from the BV quantisation of the 6d hCS action on
$X = \R^3\times K3\times\C^2$ alone. This section records the local
BV data and the source package: the $\hCS$ classical master equation of
Definition~\ref{def:bv-data-chiral}; its quantum master equation with
one-loop anomaly; the anomaly-cancellation locus on the Deligne
exceptional series (where the symmetric adjoint quartic Casimir
reduces to $(\mathrm{tr}_{\mathrm{adj}}T^2)^2$); the Costello
renormalisation-group flow on the parametrix scale $L$; the Feynman
graph expansion with integrals on $\overline{\Conf}_n(X)$; the
identification of graph weights with motivic multiple zeta values via
the Brown pairing; the coefficient package comparing the $n$-loop BV
obstruction with Jacobi--Fourier coefficients of the K3 Borcherds
input $\phi_{0,1}^{K3}$ in the Gritsenko--Nikulin normalisation; and the
Siegel--Borcherds scalar ratio
\[
\left(\frac{\Phi_{10}^{\mathrm{un}}(Z)}
{\eta(\tau)^{24}\eta(\tau')^{24}}\right)^{C_{K3}(Z;\hbar)}.
\]
The comparison becomes an identity for the BV trace only
after the Vol~III finite Hall--Borcherds source-recognition gates, the
Bruinier regularised-lift comparison, and the coefficient/source
package below are supplied. Identification of the automorphic form
$(\Phi_{10}^{\mathrm{un}})^{-1}$ with the symmetric-product BPS
generating function of Dijkgraaf--Moore--Verlinde--Verlinde $1997$
and with the attractor-point BPS counting of
Maldacena--Moore--Strominger $1999$ is at the twisted-holography
scope (Costello--Gaiotto--Paquette $2018$), not at the scope of any
dynamical-metric $3$d gravity partition function; see
Remark~\ref{rem:hcs-dvv-identification-scope}.

\subsection{Classical action and master equation}
\label{subsec:hcs-cma}

On a Calabi--Yau $3$-fold $Y$ with holomorphic $3$-form
$\Omega = dz_1\wedge dz_2\wedge dz_3$, 6d holomorphic Chern--Simons
with gauge algebra $\fg$ and field
$A\in\Omega^{0,1}(Y,\fg)$ has classical action
\begin{equation}\label{eq:hcs-classical-action}
 S_{\mathrm{cl}}[A] \;=\; \tfrac{1}{2}\int_Y \Omega\wedge
 \Tr\!\Bigl(A\wedge\bar\partial A + \tfrac{2}{3}A\wedge A\wedge A\Bigr).
\end{equation}
Following Definition~\ref{def:bv-data-chiral}, the BV field space
$\Phi = A + A^+$ with $A^+\in\Omega^{0,2}(Y,\fg^*)[1]$ carries the odd
pairing $\{A, A^+\}\propto \mathrm{id}_\fg$, and
$S_{\mathrm{cl}}$ satisfies the classical master equation
\begin{equation}\label{eq:cme}
 \{S_{\mathrm{cl}}, S_{\mathrm{cl}}\} \;=\; 0,
\end{equation}
by a direct computation: the variation $(\bar\partial A + A\wedge A)$
is the curvature, and $\Omega\wedge\Tr(F\wedge F)\in\Omega^{3,3}(Y)$ is
a closed top form whose integral is a boundary term on $Y$.

\subsection{Quantum master equation and the one-loop anomaly}
\label{subsec:hcs-qme}

The BV Laplacian
$\Delta_{\mathrm{BV}} = \delta/\delta A\cdot\delta/\delta A^+$ on the
perturbative distributional observables (Theorem~\ref{thm:bv-bar-geometric})
is ill-defined at coincident points: the kernel of
$\Delta_{\mathrm{BV}}$ is the diagonal delta
$\delta_{\Delta}\in\Omega^{0,0}(Y\times Y,\fg\otimes\fg)$, which the
propagator
$P(z,w) = \Omega(w)^{-1}\wedge\bar\partial^{-1}\delta_{\Delta}(z,w)$
inverts only up to a parametrix scale $L$.
Costello's renormalisation group
(\cite{costello-renormalization}, Ch.~10--16) introduces the
scale-$L$ heat-kernel regularised propagator
$P_{L} = \int_0^{L} K_t\,dt$, where $K_t$ is the heat-kernel
of $\bar\partial^*\bar\partial + t\cdot\mathrm{id}$, and replaces the
na\"ive quantum master equation by the regularised quantum master equation
\begin{equation}\label{eq:qme}
 \{S_{L}, S_{L}\}_{\mathrm{BV}}
 \;+\; 2\hbar\,\Delta_{L}\,S_{L}
 \;=\; 0,
\end{equation}
with $\Delta_{L} = \{P_{L},-\}$ the regularised BV Laplacian and
$S_{L}$ the scale-$L$ effective action determined by the RG equation
\begin{equation}\label{eq:rg}
 S_{L'}
 \;=\;
 \log\!\Bigl(\exp\bigl(\hbar\,\partial_{P_{L}^{L'}}\bigr) \exp(S_{L}/\hbar)\Bigr),
 \qquad P_{L}^{L'} = P_{L'}-P_{L}.
\end{equation}

\begin{proposition}[One-loop anomaly for 6d hCS]
\label{prop:hcs-one-loop-anomaly}
\ClaimStatusProvedElsewhere
On a Calabi--Yau $3$-fold $Y$ with $K_{Y}\simeq\cO_{Y}$ and simple Lie
algebra $\fg$, the scale-$L$ effective action
$S_{L} = S_{\mathrm{cl}} + \hbar\,S^{(1)}_{L} + O(\hbar^{2})$
satisfies \eqref{eq:qme} modulo $\hbar^{2}$ if and only if the one-loop
BV obstruction
\begin{equation}\label{eq:one-loop-anomaly}
 \mathrm{Obs}_{1}
 \;=\;
 \frac{1}{(2\pi i)^{3}}\,
 \mathrm{tr}_{\mathrm{adj}}\!\bigl(T^{(a}T^{b}T^{c}T^{d)}\bigr)_{\mathrm{irr}}
 \int_{Y}\Omega_{Y}\wedge
 c\wedge\bar\partial c\wedge\bar\partial c\wedge\bar\partial c
\end{equation}
vanishes in $H^{1}_{\mathrm{BRST}}\bigl(\Omega^{3,\bullet}(Y)\otimes
\Lambda^{\bullet}\fg^{\vee}\bigr)$. Here $c\in\Omega^{0,0}(Y,\fg)$ is
the BRST ghost of ghost number $+1$,
$\mathrm{tr}_{\mathrm{adj}}(T^{(a}T^{b}T^{c}T^{d)})$ is the symmetric
adjoint quartic Casimir, and the subscript $\mathrm{irr}$ projects onto
the irreducible part complementary to the reducible piece
$(\mathrm{tr}_{\mathrm{adj}}T^{2})^{2}$ absorbed by Green--Schwarz
counterterms. The adjoint cubic trace
$\mathrm{tr}_{\mathrm{adj}}(T^{a}T^{b}T^{c})$ vanishes identically for
simple $\fg$ (self-dual adjoint), so no $c(\bar\partial c)^{2}$-type
cocycle appears; the quartic is the first non-vanishing BRST class.
\end{proposition}

\begin{proof}
Expand $S_{L} = S_{\mathrm{cl}} + \hbar\,T_{L}$ and substitute into
\eqref{eq:qme}. The $\hbar^{0}$ part is \eqref{eq:cme}. The $\hbar^{1}$
part is
$\{S_{\mathrm{cl}},T_{L}\} + \Delta_{L}S_{\mathrm{cl}} = 0$,
so $T_{L}$ is determined modulo BV-exact terms by the one-loop BV
obstruction class
$\Delta_{L}S_{\mathrm{cl}}\bmod\{S_{\mathrm{cl}},-\}$.
Evaluating $\Delta_{L}$ on the BRST interaction
$\cA^{*}\cdot[\cA,c] + c^{*}\cdot[c,c]$ produces the box diagram with
four external ghost legs; the holomorphic-$\bar\partial$ propagator
integrates each loop insertion against $\Omega_{Y}$, and on $\mathrm{CY}_{3}$
the only non-vanishing one-loop contribution is the box graph with
four $c$-legs. The integrand is fixed by dimensional analysis:
$\Omega_{Y}$ has bidegree $(3,0)$, each $\bar\partial c$ has $(0,1)$,
and the three $\bar\partial c$ factors in
\(c(\bar\partial c)^3\) deliver a $(3,3)$-form after multiplication
by \(\Omega_Y\). The undifferentiated ghost supplies the
ghost-number-$+1$ factor, so the class lands in
$H^{1}_{\mathrm{BRST}}(\Omega^{3,\bullet})$. The three-leg candidate
$\Omega_{Y}\wedge A\wedge\bar\partial A\wedge\bar\partial A$ is
eliminated by the vanishing of $\mathrm{tr}_{\mathrm{adj}}(T^{a}T^{b}T^{c})$
on simple $\fg$. The coefficient $(2\pi i)^{-3}$ is the normalisation
of the three-complex-dimensional propagator kernel on $\mathrm{CY}_{3}$. Scale
independence follows from the RG equation \eqref{eq:rg}: the
connecting propagator $P_{L}^{L'}$ is BV-exact, so
$\partial_{L}\,\mathrm{Obs}_{1} = 0$. See Theorem~\ref{thm:6dhcs-one-loop-anomaly}
for the chain-level descent argument and Costello~\cite{Costello2015omega},
Costello--Li~\cite{CL16} for the full derivation.
\end{proof}

\begin{remark}[Anomaly cancellation locus]
\label{rem:hcs-anomaly-locus}
\index{anomaly cancellation!6d hCS}
For pure adjoint hCS the anomaly~\eqref{eq:one-loop-anomaly} factors
through the symmetric adjoint quartic
$\mathrm{tr}_{\mathrm{adj}}(T^{(a}T^{b}T^{c}T^{d)})_{\mathrm{irr}}$.
The symmetric adjoint cubic
\[
\mathrm{tr}_{\mathrm{adj}}(T^{(a}T^{b}T^{c)})
\]
vanishes for every simple Lie algebra because the adjoint
representation is self-dual. On the
Deligne exceptional series the quartic factorises through the
quadratic squared,
$\mathrm{tr}_{\mathrm{adj}}T^{4} = \alpha_{\fg}\cdot
(\mathrm{tr}_{\mathrm{adj}}T^{2})^{2}$ via the Vogel plane
\cite{Deligne1996,Costello2017M5,CL16}; the residual
$(F^{2})^{2}$-term is absorbed by Green--Schwarz descent. Cubic
tensors in the fundamental representation of \(A_2\) or in the
\(\mathbf{27}\) of \(E_6\) are matter-representation data; they are
not adjoint hCS anomaly classes. The 6d hCS one-loop anomaly-free
locus on a CY$_{3}$ is
\begin{equation}\label{eq:hcs-anom-locus}
\mathrm{Anom}_{1}=0\;\iff\;
\fg\in\mathrm{Deligne}^{\mathrm{exc}}
\cup\{\mathrm{abelian}\}
\cup\{\mathrm{super\text{-}str}_{\mathrm{ad}}=0\}
\cup\{\widehat\fg_{-h^{\vee}}\otimes K^{-1/2}\text{-refined}\}.
\end{equation}
For the
abelian class-$\mathbf G$ case $\fg=\fu(1)$ (Heisenberg) both
the adjoint cubic and irreducible quartic vanish identically. For non-exceptional simple $\fg$
outside the Deligne series the irreducible quartic survives, and
the effective boundary algebra at the critical twist is
$\widehat{\fg}_{-h^{\vee}}\otimes K^{-1/2}$. On the twisted-holography
total space $X = \R^{3}\times K3\times\C^{2}$ the relevant
gauge algebra $\fg = \fgl_{1}$ sits in the abelian branch, and the
non-trivial dynamics resides in the K3 normal-bundle twist.
\end{remark}

\subsection{Feynman graph expansion and configuration-space integrals}
\label{subsec:hcs-feynman}

The scale-$\infty$ effective action
$S_{\infty} = \lim_{L\to\infty}S_{L}$, when it exists, is the full
quantum-corrected action. On
$X = \R^{3}\times K3\times\C^{2}$ (Calabi--Yau threefold with
holomorphic volume form $\Omega_{X} = dz\wedge\Omega_{K3}\wedge dw_{1}\wedge dw_{2}$;
the $\R^{3}$ direction is topological, $\C^{2}$ is the Omega-background,
and $K3$ carries a hyper-K\"ahler structure), the expansion reads
\begin{equation}\label{eq:hcs-feynman-expansion}
 S_{\infty}[A]
 \;=\;
 \sum_{g\ge 0}\hbar^{g}\,
 \sum_{\Gamma:b_{1}(\Gamma) = g}
 \frac{1}{|\mathrm{Aut}(\Gamma)|}\,
 w_{\fg}(\Gamma)\,
 I_{X}(\Gamma)\,
 \mathcal{O}(\Gamma)[A],
\end{equation}
where $w_{\fg}(\Gamma)$ is the Lie-algebra weight from $\fg$-Casimirs at
vertices and trace data at edges, $\mathcal{O}(\Gamma)[A]$ is the
external-leg operator insertion, and $I_{X}(\Gamma)$ is the
regularised configuration-space integral
\begin{equation}\label{eq:hcs-conf-integral}
 I_{X}(\Gamma)
 \;=\;
 \int_{\overline{\Conf}_{V(\Gamma)}(X)}
 \bigwedge_{e\in E(\Gamma)} P(z_{s(e)},z_{t(e)}),
\end{equation}
with $\overline{\Conf}_{n}(X)$ the Fulton--MacPherson compactification.

\begin{proposition}[Integrand regularity on $\overline{\Conf}_{n}(X)$]
\label{prop:hcs-integrand-regular}
\ClaimStatusProvedElsewhere
The propagator $P(z,w)$ extends to a smooth $(2,3)$-form on the
blow-up $\overline{\Conf}_{2}(X)\to X\times X$; the wedge product
$\bigwedge_{e}P$ extends to a smooth top form on
$\overline{\Conf}_{V(\Gamma)}(X)$, and the integral
\eqref{eq:hcs-conf-integral} is finite.
\end{proposition}

\begin{proof}
This is the Axelrod--Singer regularity theorem
\cite{Kon94} transplanted to 6d hCS. The Fulton--MacPherson blow-up
replaces each diagonal $\Delta_{ij}\subset X^{n}$ by its projectivised
normal bundle, pulling $1/(z_{i}-z_{j})^{3}$ singularities back to
$(3,3)$-forms with logarithmic poles along the exceptional divisor. The
$\Omega_{X}$ factor in $P$ absorbs the logarithm. Total degree:
each edge contributes a $(2,3)$-form, and the external vertices
contribute $(0,1)$; the top-form condition on a $6|V|$-real-dimensional
manifold becomes $3|E|+|V|=6|V|$, which is the Euler equation
$|V|-|E| = 1-g$ for a connected graph of loop order $g$, giving
$|E| = |V|+g-1$ so $3(|V|+g-1)+|V| = 4|V|+3g-3$; matching to
$6|V|$ forces $g = (2|V|+3)/3$, a number-theoretic constraint on
admissible graphs that isolates the motivic locus.
\end{proof}

\subsection{Kontsevich graph weights $=$ motivic MZV}
\label{subsec:hcs-motivic-mzv}

The graph weight $I_{X}(\Gamma)$ of \eqref{eq:hcs-conf-integral} is,
after separating the $\R^{3}$-topological fibre and the $K3\times\C^{2}$
holomorphic fibre, a product
\begin{equation}\label{eq:graph-weight-factorisation}
 I_{X}(\Gamma) \;=\; I_{\R^{3}}(\Gamma)\cdot I_{K3\times\C^{2}}(\Gamma),
\end{equation}
where $I_{\R^{3}}(\Gamma)$ is the Kontsevich integral on
$\overline{\Conf}_{V(\Gamma)}(\R^{3})$ and $I_{K3\times\C^{2}}(\Gamma)$
is its holomorphic partner on
$\overline{\Conf}_{V(\Gamma)}(K3\times\C^{2})$.

\begin{theorem}[Periods on $\overline{M}_{g,n}$ $=$ motivic MZVs]
\label{thm:hcs-motivic}
\ClaimStatusProvedElsewhere
For each connected graph $\Gamma$ of loop order $g = b_{1}(\Gamma)$:
\begin{enumerate}[label=\textup{(\roman*)}]
\item $I_{\R^{3}}(\Gamma) \in \mathcal{Z}_{\mathrm{mot}}$, the ring of
motivic multiple zeta values of weight $3|V(\Gamma)|-6$
(Brown~\cite{Brown2012});
\item $I_{K3\times\C^{2}}(\Gamma)$ is a period of
$\pi_{1}^{\mathrm{mot}}(\overline{M}_{g,n}\setminus\partial)$ with
K3-type mixed Hodge structure
(Brown~\cite{Brown2017});
\item Under the identification
$\overline{\Conf}_{n}(\R^{3})\simeq\overline{M}_{0,n+1}$ via
Axelrod--Singer, $I_{\R^{3}}(\Gamma)$ coincides with the motivic period of the
Drinfeld associator obtained at genus $0$ in
Theorem~\ref{thm:drinfeld-associator-bar}.
\end{enumerate}
\end{theorem}

\begin{proof}
(i)~Brown's theorem on periods of $\overline{M}_{0,n}$: the integrals of
Arnold forms $\eta_{ij} = d\log(z_{i}-z_{j})$ on the real locus
generate $\mathcal{Z}_{\mathrm{mot}}$, and the Kontsevich propagator
$d\arg(z_{i}-z_{j})$ on $\R^{3}$ is the real part of $\eta_{ij}$ restricted
to $S^{1}$ (Proposition~\ref{prop:propagator-restriction}). Graph
integrals of the wedge products are iterated Chen integrals, known to
span $\mathcal{Z}_{\mathrm{mot}}$ up to the Brown generators at each
weight.
(ii)~Brown~\cite{Brown2017} extends the period interpretation to
$\overline{M}_{g,n}$; the K3 factor provides the elliptic-period layer
via the Kuga--Sato family over $\overline{M}_{g,n}$.
(iii)~Drinfeld-associator comparison: the $n = 3$ genus-0 case equals
the Drinfeld--Kohno KZ monodromy (Theorem~\ref{thm:drinfeld-associator-bar}),
whose real part is the ordered Kontsevich propagator.
\end{proof}

\subsection{BV obstruction coefficient package}
\label{subsec:hcs-bv-obstruction}

The BV obstruction at loop order $g$, evaluated on the
$X = \R^{3}\times K3\times\C^{2}$ target, supplies the Costello-side
coefficient source that must be compared with the half K3 Jacobi form.

\begin{proposition}[Coefficient package for the K3 Jacobi /
\(\Delta_5\) comparison]
\label{prop:hcs-obstruction-jacobi}
Let \(\phi_{0,1}^{K3}\) be the half K3 weak Jacobi form
\[
\phi_{0,1}^{K3}(\tau,z)=y^{-1}+10+y
 +q(10y^{-2}-64y^{-1}+108-64y+10y^2)+O(q^2),
\]
so that the full Eguchi--Ooguri--Tachikawa elliptic genus is
\(\chi_{K3}=2\phi_{0,1}^{K3}\) \cite{EOT11}. The half form is
holomorphic at the cusp:
\[
 \phi_{0,1}^{K3}(\tau,z)
 \;=\;
 \sum_{a\ge 0,\,r}c_{\phi_{0,1}^{K3}}(a,r)\,q^{a}y^{r},
 \quad q = e^{2\pi i\tau},\; y = e^{2\pi i z}.
\]
The connected $n$-loop BV obstruction is compared with the
\(p^n\)-Fourier--Jacobi coefficient of the separating-boundary logarithm
\[
\mathcal L_{\Phi}(Z)
:=
-\log\frac{\Phi_{10}^{\mathrm{un}}(Z)}
{\eta(\tau)^{24}\eta(\tau')^{24}}
=
-2\log\frac{\Delta_5(Z)}
{\eta(\tau)^{12}\eta(\tau')^{12}},
\qquad p=e^{2\pi i\tau'} .
\]
The coefficient extraction is performed before the Humbert
specialisation:
\begin{equation}\label{eq:bv-obstruction-jacobi-matching}
 \ell_{n}^{\mathrm{BV,target}}(X)
 \;:=\;
 [p^n]\,\mathsf H_{\mathrm{Br}}\!\left(\mathcal L_{\Phi}(Z)\right).
\end{equation}
Here \(\mathsf H_{\mathrm{Br}}\) is the Bruinier Humbert projection from
Fourier--Jacobi coefficients to divisor classes on \(\overline{\mathcal A}_2\).
The reciprocal coefficient
\[
b_n(X):=[p^n]\Delta_5(Z)^{-1}
\]
is obtained from the logarithmic coefficients
\(\ell_1,\ldots,\ell_n\) by exponentiation; it is not equal to
\(\ell_n\) termwise.
The two $\eta^{24}$ factors remove the elliptic discriminant factors in
the separating-boundary expansion. The residual automorphy is the
automorphy of the displayed quotient; a weight-zero boundary function
requires the standard \(z^2\)-normalisation at the Humbert divisor.
The genus-one hCS normalisation is the separate Euler identity
\[
\frac1{24}\int_{K3}e(T_{K3})=1,\qquad
\chi_{K3}(\tau,0)=24,
\]
whereas the BKM weight is
\(c_{\phi_{0,1}^{K3}}(0,0)/2=10/2=5\). The equality
\(\ell_{n}^{\mathrm{BV}}(X)=\ell_{n}^{\mathrm{BV,target}}(X)\) is a theorem
only after the Vol~III finite Hall--Borcherds source-recognition gates
and the Bruinier regularised-lift comparison identify the BV source
with the Borcherds lift source.
\end{proposition}

\begin{remark}[$\phi_{-2,1}$ is holomorphic at the cusp]
\label{rem:phi-weak-jacobi-holomorphic}
The weak Jacobi form
$\phi_{-2,1}(\tau,z) = -\theta_1(\tau,z)^{2}/\eta(\tau)^{6}$
of weight $-2$, index $1$
\cite[\S 9, Theorem 9.3]{EichlerZagier1985} has $q$-expansion
$(y - 2 + y^{-1}) + q(-2y^{2}+8y-12+8y^{-1}-2y^{-2})+O(q^{2})$
with no polar $q^{-m}$ term. Equivalently
$c_{\phi_{-2,1}}(m,r) = 0$ for $m < 0$. The $24$ appearing in the
BV loop tower is therefore not a polar Fourier coefficient of bare
$\phi_{-2,1}$; it is $\chi(K3)$ read off $\chi_{K3}|_{y=1}|_{q^{0}}$
or equivalently the first reciprocal coefficient of
the Gritsenko--Nikulin Borcherds product. The positive-discriminant
Fourier coefficients of \(\phi_{0,1}^{K3}\) appearing as
Heegner-divisor multipliers in
Theorem~\ref{thm:htbb-heegner-all-order} and in the Borcherds product
exponents of \eqref{eq:delta5-borcherds-lift}, are canonical
Bruinier--Borcherds invariants and do not conflict with the above.
\end{remark}

\subsection{All-loop scalar target as a Borcherds product}
\label{subsec:hcs-borcherds}

The Gritsenko--Nikulin weight-$5$ form $\Delta_5$ is the Borcherds lift
of the normalised half K3 Jacobi input $\phi_{0,1}^{K3}$: explicitly
\begin{equation}\label{eq:delta5-borcherds-lift}
 \Delta_5(Z)
 \;=\;
 64(q\,y\,q')^{1/2}\cdot
 \prod_{\substack{(n,r,m)>0}}
 \bigl(1 - q^{n}y^{r}q'^{m}\bigr)^{c_{\phi_{0,1}^{K3}}(nm,r)},
 \quad
 Z = \begin{pmatrix}\tau & z \\ z & \tau'\end{pmatrix}\in\mathfrak{H}_{2}.
\end{equation}
The unnormalised Igusa square is \(\Phi_{10}^{\mathrm{un}}=\Delta_5^2\);
its logarithmic product exponents are twice those in
\eqref{eq:delta5-borcherds-lift}.

\begin{theorem}[All-loop 6d hCS on $K3\times T^{2}$: gated
Siegel--Borcherds scalar target]
\label{thm:hcs-all-loop-resummation}
\index{6d holomorphic Chern-Simons!all-loop effective action}
\emph{Scope.} The object $Z_{\hCS}(X,\hbar)$ denotes the
BV-renormalised effective-action exponential of the Costello--Li
holomorphic twist on $X$, not a dynamical-metric partition function.
On the twisted holomorphic background
$X = \R^{3}\times K3\times T^{2}$ reduction of 6d hCS at $\fg = \fgl_{1}$
(Omega-deformed, with $24$ M5-branes wrapping the $I_{1}$ Kodaira
fibres of the elliptic $K3$, in the Costello--Gaiotto--Paquette
twisted-holography frame~\cite{CostelloGaiottoPaquette2018}), the
gated scalar comparison target is the Siegel--Borcherds scalar ratio
\begin{equation}\label{eq:hcs-k3-coeff-series}
 C_{K3}(Z;\hbar)
 \;:=\;
 \sum_{g\ge 1}\hbar^{g}\,b_g(Z),
 \qquad
 b_g(Z):=[q^g]\Delta_5(Z)^{-1},
\end{equation}
\begin{equation}\label{eq:hcs-all-loop-resummation}
 Z_{\hCS}^{\mathrm{target}}(X,\hbar)
 \;:=\;
 \exp\!\Bigl(
 C_{K3}(Z;\hbar)\,
 \log\frac{\Phi_{10}^{\mathrm{un}}(Z)}{\eta(\tau)^{24}\,\eta(\tau')^{24}}\Bigr)
 \;=\;
 \Bigl(\frac{\Phi_{10}^{\mathrm{un}}(Z)}{\eta(\tau)^{24}\,\eta(\tau')^{24}}\Bigr)^{C_{K3}(Z;\hbar)},
\end{equation}
with $Z=\begin{pmatrix}\tau&z\\z&\tau'\end{pmatrix}\in\mathfrak H_{2}$.
The coefficients $b_g(Z)$ are Fourier--Jacobi coefficients of the
reciprocal \(\Delta_5^{-1}\). They are not the primitive Borcherds
exponents. The primitive exponents in the
Gritsenko--Nikulin product expansion~\cite[(2.2)]{GritsenkoNikulin1998}
are the Fourier coefficients
\[
e_{\Delta_5}(n,r,m)=c_{\phi_{0,1}^{K3}}(nm,r)
\]
of the weight-$0$, index-$1$ Borcherds-input weak Jacobi form;
the coefficients \(b_g(Z)\) are obtained from these \(e(n,r,m)\) only
after expanding the infinite product. The primitive exponents and the
reciprocal Fourier--Jacobi coefficients are therefore different
coefficient systems attached to the same K3 Borcherds input. The two $\eta^{24}$
factors remove the elliptic discriminant factors in the separating
boundary expansion. The residual automorphy is the automorphy of the
displayed quotient; a weight-zero boundary function requires the
standard \(z^2\)-normalisation at the Humbert divisor.  The identity
$Z_{\hCS}(X,\hbar)=Z_{\hCS}^{\mathrm{target}}(X,\hbar)$ is asserted
only under the Vol~III finite Hall--Borcherds source-recognition gates,
the Bruinier regularised-lift comparison, and the coefficient/source
package of Proposition~\ref{prop:hcs-obstruction-jacobi}; without
those inputs, \eqref{eq:hcs-all-loop-resummation} is the scalar target
identity.
\end{theorem}

\begin{remark}[Identification with the symmetric-product BPS
generating function; scope-qualified]
\label{rem:hcs-dvv-identification-scope}
\index{scope!DVV BPS counting vs.\ dynamical 3d gravity}
The Dijkgraaf--Moore--Verlinde--Verlinde
$1997$~\cite{DijkgraafMooreVerlindeVerlinde1997}
generating function of symmetric-product elliptic genera on the
$N$-fold symmetric orbifold $\mathrm{Sym}^N(K3)$,
\[
 \sum_{N \ge 0} p^N\, \chi_{\mathrm{ell}}(\mathrm{Sym}^N(K3); \tau, z)
 \;=\; \frac{1}{\Phi_{10}^{\mathrm{un}}(Z)}
 \quad \text{(after multiplicative shift to Igusa normalisation)},
\]
coincides with the inverse of the Gritsenko--Nikulin Igusa form
$\Phi_{10}^{\mathrm{un}}$ as an automorphic form on $\mathfrak{H}_2$;
under the Maldacena--Moore--Strominger
$1999$~\cite{MaldacenaMooreStrominger1999} identification, the same
automorphic form counts $1/4$-BPS D1--D5--P microstates of type~IIB
on $K3 \times S^1$ at the attractor point. Both are boundary-CFT
or target-space objects: neither is the partition function of a
dynamical-metric $3$d gravity theory on $\mathrm{AdS}_3$ or
$\mathrm{AdS}_3 \times K3$. The notation
\[
 Z_{\mathrm{3dQG}}^{\mathrm{AdS}_{3}\times K3}(Z)
 \;\coloneqq\; \bigl(\Phi_{10}^{\mathrm{un}}(Z)\bigr)^{-1},
 \quad \text{(definitional at the twisted-holography scope)}
\]
is the \emph{twisted-holography boundary character}, identified with
the symmetric-product BPS generating function by DVV $1997$ and with
the attractor-point BPS microstate count by MMS $1999$. The
identification with a genuine dynamical-metric $3$d quantum gravity
path integral is conjectural (Maloney--Witten $2009$
sum plus wormhole corrections, not convergent without off-shell
Maxfield--Turiaci $2020$ contributions). The target object
\eqref{eq:hcs-all-loop-resummation} has the automorphic boundary
character $(\Phi_{10}^{\mathrm{un}})^{-1}$ along the large-$\hbar$ BV
filtration limit below; identifying it with the Costello 6d hCS
effective action requires the same recognition gates.
\end{remark}

\begin{proposition}[Large-$\hbar$ target limit and the
twisted-holography boundary character]
\label{prop:hcs-large-hbar-limit}
\index{large-hbar@large-$\hbar$!BV limit}
For the scalar target of Theorem~\ref{thm:hcs-all-loop-resummation},
the pole along the Humbert discriminant
$\{\Phi_{10}^{\mathrm{un}}=0\}\subset\mathcal{A}_{2}$
controls every finite truncation of
$Z_{\hCS}^{\mathrm{target}}(X,\hbar)$. For
\[
C_{K3}^{\leq N}(Z;\hbar)=\sum_{1\leq g\leq N}\hbar^g b_g(Z)
\]
one has
\[
\log Z_{\hCS,N}^{\mathrm{target}}(X,\hbar)
=
C_{K3}^{\leq N}(Z;\hbar)
\log\frac{\Phi_{10}^{\mathrm{un}}(Z)}
{\eta(\tau)^{24}\eta(\tau')^{24}}.
\]
Thus the large-\(\hbar\) behaviour of a finite truncation is governed
by its highest retained loop coefficient; there is no collapse to a
single factor \(\hbar\,c_{K3}(Z)\). The
identification of $(\Phi_{10}^{\mathrm{un}})^{-1}$ as the Gritsenko--Nikulin--Borcherds
singular theta lift of the K3 elliptic-genus generator
$\phi_{0,1}^{K3}$ is the content of the Borcherds multiplicative theta
correspondence~\cite{GritsenkoNikulin1998}; no stationary-phase or
Laplace-method saddle-point approximation is performed, and the
phrase ``large-$\hbar$ saddle'' denotes only the BV-grading limit on
the effective-action complex.  The same asymptotic statement for
$Z_{\hCS}$ itself is conditional on the gates in
Theorem~\ref{thm:hcs-all-loop-resummation}.
\end{proposition}

\begin{proof}[Source package for Theorem~\ref{thm:hcs-all-loop-resummation}]
The target coefficient package \eqref{eq:bv-obstruction-jacobi-matching}
enters the scalar comparison of \eqref{eq:hcs-feynman-expansion}: the
target coefficient at $\hbar^{g}$ is the \(p^g\)-Fourier--Jacobi
coefficient of \(\mathcal L_{\Phi}\) after Bruinier projection. The
Bruinier--Borcherds regularised lift of \(\phi_{0,1}^{K3}\) is supplied by
the source-recognition comparison; the BV graph calculation supplies the
Costello-side coefficient source.  At each loop order, the target
bubble-chain contribution is governed by
\(\log\bigl(\Phi_{10}^{\mathrm{un}}/(\eta(\tau)^{24}\eta(\tau')^{24})\bigr)\).
Near the separating divisor \(z=0\) in the Siegel period, the quotient has
the expansion
\[
 \frac{\Phi_{10}^{\mathrm{un}}(Z)}
 {\eta(\tau)^{24}\eta(\tau')^{24}}
 \;=\;
 u_{\mathrm{GN}}\,z^{2}+O(z^{4})
\]
with \(u_{\mathrm{GN}}\in\mathbb C^\times\) fixed by the
Gritsenko--Nikulin normalisation
\cite{GritsenkoNikulin1998,DijkgraafVerlindeVerlinde1997},
so the separating-boundary quotient carries the residual automorphy
described above after division by both
$\eta(\tau)^{24}$ and $\eta(\tau')^{24}$; the
$24$-fold $I_{1}$-fibre combinatorics of
Remark~\ref{rem:pf-deep-2} supplies each $\eta^{24}$ factor from
$24$ M5-brane one-loop integrals threading the respective elliptic
fibre as target evidence at the Costello--Gaiotto--Paquette
twisted-holography scope.  Summing the loop-degree target coefficients
gives \eqref{eq:hcs-all-loop-resummation}; the final comparison with
$Z_{\hCS}$ requires the finite Hall--Borcherds source-recognition gates.
\end{proof}

\begin{proof}[Proof of Proposition~\ref{prop:hcs-large-hbar-limit}]
From the finite-truncation identity in
Proposition~\ref{prop:hcs-large-hbar-limit},
\[
\log Z_{\hCS,N}^{\mathrm{target}}(X,\hbar)
=
\left(\sum_{1\leq g\leq N}\hbar^g b_g(Z)\right)
\log\frac{\Phi_{10}^{\mathrm{un}}}
{\eta(\tau)^{24}\eta(\tau')^{24}},
\]
the leading term along any ray \(\hbar\to\infty\) is the highest
nonzero retained loop coefficient times the displayed logarithm. The
pole structure along $\{\Phi_{10}^{\mathrm{un}}=0\}\subset
\mathcal{A}_2$ is controlled by the quadratic-divisor theorem of
Gritsenko--Nikulin~\cite{GritsenkoNikulin1998}:
$\Phi_{10}^{\mathrm{un}}$ vanishes to order $2$ along the Humbert divisor
$H_{1}\subset\mathcal A_{2}$.
The identification of $(\Phi_{10}^{\mathrm{un}})^{-1}$ as the Borcherds multiplicative
lift of the K3 Borcherds input \(\phi_{0,1}^{K3}\) is
the content of the singular theta correspondence of
Harvey--Moore~\cite{HarveyMoore1996} and
Borcherds~\cite{Borcherds1998}; no residue operation and no
stationary-phase approximation is performed. The phrase
``large-$\hbar$ saddle'' in the BV context denotes the filtration
limit on the effective-action complex, not a Laplace-method
approximation of a convergent integral. Identification with the
symmetric-product BPS generating function (DVV $1997$) and with the
attractor-point BPS counting (MMS $1999$) is the content of
Remark~\ref{rem:hcs-dvv-identification-scope}, at the stated
twisted-holography scope.
\end{proof}

\begin{remark}[Loop order versus Casimir order in the Borcherds target]
\label{rem:hcs-loop-vs-casimir}
The Fourier--Jacobi coefficient \([p^n]\mathcal L_{\Phi}\) in
\eqref{eq:bv-obstruction-jacobi-matching} indexes
\emph{loop order} $n$, not Casimir order. The target BV coefficient
package is organised by loop degree: loop~$1$ carries the symmetric
quartic Casimir \eqref{eq:one-loop-anomaly} of
Proposition~\ref{prop:hcs-one-loop-anomaly}, loop~$n \geq 2$ carries
higher adjoint-invariant polynomials built from iterated tadpole and
box contractions of the cubic BRST vertex. At the 3d-gravity specialisation
$\fg = \fgl_{1}$, every adjoint-polynomial structure is trivial and
the target product reduces to the K3-normal-bundle combinatorics: the
$24$ $I_{1}$-fibres of Remark~\ref{rem:pf-deep-2} supply
$\eta(\tau)^{24}$ along the $\tau$-modular fibre and $\eta(\tau')^{24}$
along the $\tau'$-modular fibre of the Siegel period
$Z=\begin{pmatrix}\tau&z\\z&\tau'\end{pmatrix}$, and the
$\phi_{0,1}^{K3}$ Fourier data supplies the Borcherds numerator.
The all-loop scalar target
\[
\left(\frac{\Phi_{10}^{\mathrm{un}}}
{\eta(\tau)^{24}\eta(\tau')^{24}}\right)^{C_{K3}(Z;\hbar)}
\]
is therefore Casimir-independent at the gravitational locus and
belongs to the coefficient/source package without recourse to the
quadratic--quartic distinction of
Remark~\ref{rem:hcs-anomaly-locus}.  It is an identity for the BV tower
only after the gates of Theorem~\ref{thm:hcs-all-loop-resummation}.
\end{remark}

\begin{remark}[Costello-side versus Vol~III CY-side scalar comparisons]
\label{rem:hcs-two-sides}
Theorem~\ref{thm:hcs-all-loop-resummation} is the perturbative
Costello-side comparison for the scalar Borcherds all-loop target: the BV
master equation \eqref{eq:cme}, the renormalised QME \eqref{eq:qme},
the Feynman expansion \eqref{eq:hcs-feynman-expansion}, the
configuration-space integrals \eqref{eq:hcs-conf-integral} with motivic
content \ref{thm:hcs-motivic}, the obstruction-Jacobi matching
\ref{prop:hcs-obstruction-jacobi}, and the all-loop target identity
\eqref{eq:hcs-all-loop-resummation} stand in a single comparison
diagram. The complementary Vol~III CY-side recognition package uses
the CY-to-chiral functor $\Phi\colon\mathrm{Perf}(K3)\to\mathrm{ChirAlg}$
to compare the graded dimensions of $\mathrm{Ext}^{g}$ on the derived
category of K3 with the Fourier--Jacobi coefficients
\([p^g]\mathcal L_{\Phi}\), equivalently with the root-label coefficients
\(c_{\phi_{0,1}^{K3}}(a_g,r_g)\) after the Bruinier projection. The two
comparison packages agree only after
the Vol~III finite Hall source, pairing, PBW/no-extra-relations,
inverse-limit, and Borcherds comparison gates, together with the
Bruinier regularised-lift comparison, identify their sources with the
same universal $\Phi$-trace identity
(Remark~\ref{rem:3dqg-phi-trace-identity}).
\end{remark}

\section{Synthesis (continued)}
\label{sec:pf-deep}

\begin{remark}[Perturbative finiteness of twisted 11D SUGRA on $K3\times T^2$]
\label{rem:pf-deep-1}
At the K3 base, perturbative finiteness has the one-loop scalar target
\[
Z_{\mathrm{twisted\ 11D}}^{K3\times T^2}\big|_{\mathrm{all\ loops}} \;=\; Z_{\mathrm{twisted\ 11D}}^{K3\times T^2}\big|_{\mathrm{1\text{-}loop}} \;=\; L(s,\Phi_{10}^{\mathrm{un}})
\]
via Witten 1995 M-theory/$K3\times T^2$ duality.  The displayed
equality is a theorem only under the compact Hall--Borcherds
source-recognition hypotheses of Vol~III and the Vol~II vanishing
hypothesis that the completed HT graph-complex obstruction classes
outside the recognised comparison line vanish.  Under those hypotheses
the first obstruction maps by Bruinier's Heegner characteristic
comparison and Etingof--Kazhdan super-quasi-Hopf functoriality to
\[
H^2_{\mathrm{Lie}}\bigl(\mathfrak q_{K(1)};\mathrm{grt}_1\bigr)
\;\cong\; \mathbb C\cdot\Delta_5,
\]
while higher-loop contributions land in higher Hochschild degree
$(\ge 3)$ and are killed by Theorem~H concentration.  Without those
recognition and vanishing inputs, $L(s,\Phi_{10}^{\mathrm{un}})$ is the
honest scalar target comparison, not a direct construction of
$\mathbf H_{\Delta_5}$ and not a direct BKM Lie-cohomology
classification.  Primary-lit: Witten 1995, Bruinier 2002,
Etingof--Kazhdan 2007 Part~V.
\end{remark}

\begin{remark}[$24$ M$5$-brane loop-integral scalar target]
\label{rem:pf-deep-2}
The perturbative expansion on $K3\times T^2$ has $24$-fold loop
structure: each $M$-theory loop threading a single Kodaira-I$_1$ fibre
contributes an independent one-loop integral, summed over the $24$
Steiner-rigid fibres. The scalar target comparison asks for the
normalised identity
\[
\sum_{a=1}^{24}\, \mathcal I^{(1)}_a
\;\stackrel{\mathrm{target}}{=}\;
C_{\mathrm{M5}}\,\Phi_{10}^{\mathrm{un}}
\]
with a normalisation constant $C_{\mathrm{M5}}$ fixed by the
coefficient/source package. The $S(5,8,24)$ octad structure supplies
evidence for this closure: each $5$-subset of fibres extends to a
unique octad contribution with Wilson-coefficient
$c_\sigma\cdot m_\sigma$ (Golay weight, umbral shift). It becomes a
source-recognised identity only after the finite Hall--Borcherds gates
and the Bruinier regularised-lift comparison are supplied. Primary-lit:
Conway-Sloane 1988, Eguchi-Ooguri-Tachikawa 2011.
\end{remark}
